% !TeX root = ../../Book3.tex
% 纲要标题：## 第12章　弱解与局部椭圆估计
% 纲要位置：几何学计划纲要.md，第 2550 行。
% 正文状态：初稿；尚未正式审核。

% 纲要内容开始
% 
% > 本章目标与证明范围：把可解性、内部估计和正则性分开。完整建立后续 Hodge 与 Cauchy–Riemann 模块实际使用的 L² 和所选 Lᵖ 估计；标量最大值原理不得替代系统论证。
% 
% 纲要内容结束

\chapter{弱解与局部椭圆估计}
\label{g3:ch:12}
\BookChapterNumber{g3:ch:12}

\begin{chapterabstract}
椭圆方程为什么能从较弱的右端恢复较强的导数控制？本章先由 Hilbert 空间的投影、表示与能量方法构造弱解，再以主符号和 Fourier 变换揭示高频增益。冻结系数、截断交换子与差商把常系数估计移到局部；奇异积分及 Hölder 振荡估计进一步给出所需的 \(L^p\) 与 Schauder 正则性。随后把这些结果用于 Cauchy--Riemann 型算子，最后用屏障与 Moser 迭代建立实标量的最大值原理和 Harnack 不等式。每一种结论分别说明算子、系数、函数空间与区域条件，使可解性、内部正则性和次序性质各有明确依据。
\end{chapterabstract}

\begin{learninggoals}
\begin{itemize}
	\item 由投影与 Riesz 表示证明 Lax--Milgram 定理，用直接方法处理弱下半连续的约束极小问题。
	\item 写出弱 Poisson 方程，辨别零边界与闭流形上零均值条件的作用。
	\item 计算主符号与形式伴随，区分椭圆性、强椭圆性和能量型的强制性。
	\item 完整使用 Fourier、冻结系数与差商证明内部 \(L^2\) 估计，并从分布解开始正则性提升。
	\item 通过奇异积分建立 \(1<p<\infty\) 的局部估计，追踪 Hölder 估计中的尺度、插值与系数条件。
	\item 推导 Cauchy--Riemann 型系统的局部 \(W^{1,p}\) 估计，说明闭曲面、带边界域与无限条带的差别。
	\item 证明标量最大值原理、Hopf 引理与所选 Harnack 不等式，并辨明它们不能直接用于一般系统的原因。
\end{itemize}
\end{learninggoals}

给定一个源项 \(f\)，考虑零边界问题 \(-\sum_j\partial_j^2u=f\)。若只知道 \(f\in L^2\)，一开始便要求 \(u\) 有连续二阶导数，会把求解空间限制得过早。分部积分把方程改写为
\[
	\int_\Omega \mathrm Du\cdot\overline{\mathrm Dv}\,\mathrm dx
	=\int_\Omega f\overline v\,\mathrm dx,
	\qquad v\in H_0^1(\Omega).
\]
现在只需一阶弱导数，能量下界和完备性就可能给出解。但这一等式并未直接说明二阶导数属于 \(L^2\)，更未保证边界附近的全部导数可控。存在性与正则性由此成为相连而各有证明任务的两个问题。

\ProseParagraphBreak

在全空间上，Fourier 变换把 \(-\sum_j\partial_j^2\) 变成乘数 \(|\xi|^2\)。当频率很高时，由右端除以这个乘数，恰好获得两阶导数的增益；低频部分仍须用函数本身的范数控制。变系数使频率不再独立，但在足够小的球中，最高阶系数接近其中心值。把偏差作为可吸收的小项，再控制截断带来的低一阶交换子，就能保留这种局部增益。这一思路对矩阵系统同样有效，其条件是每个非零实余切方向上的主符号可逆。

\ProseParagraphBreak

从 \(L^2\) 转到 \(L^p\)，频率乘法还须转化为积分范数控制；从积分范数转到 Hölder 范数，则须估计核的消去与函数的振荡。本章分别给出所需的奇异积分证明和 Schauder 证明，不把端点 \(p=1,\infty\) 混入 \(1<p<\infty\) 的结论。对一般光滑系数椭圆系统，局部 \(L^2\) 与 \(L^p\) 正则性从分布解出发；Schauder 部分处理实标量正定主部及其适用的分量形式。标量最大值原理与 Harnack 不等式的证明还使用实数的次序结构，需要单独建立。

\ProseParagraphBreak

本章承接第十一章的弱导数、嵌入、Poincaré 不等式、紧性及迹，补足能量方法与 Fourier 分析的实际证明。Hilbert 空间的投影、表示与变分方法可对照 Brezis 第5.1—5.3节\cite[第5.1—5.3节]{Brezis2011}。本章的正则性估计以内部区域为对象；闭流形可用有限坐标覆盖拼接，边界与非紧端部则须另列条件。下一章将在这些内部工具上建立全局 Fredholm 理论、椭圆复形的调和模型及非线性映射的共同工具。

% !TeX root = ../../Book3.tex
% 纲要标题：### 12.1 能量极小化与弱解
% 纲要位置：几何学计划纲要.md，第 2554 行。
% 正文状态：初稿；尚未正式审核。

% 纲要内容开始
% 
% * Hilbert 投影和 Riesz 表示
% * Lax–Milgram 的证明
% * 约束最小化与弱下半连续
% * Poisson 方程的存在与唯一性
% 
% 纲要内容结束

\section{能量极小化与弱解}
\label{g3:sec:1201}

求解 \(\Delta_g u=f\) 时，直接寻找二阶可微函数往往过早要求了正则性。能量方法先在 \(H^1\) 中求一个满足积分恒等式的函数，再研究它是否有更多导数。存在性依赖能量的正性及允许变化的空间；这两个条件可以先在抽象 Hilbert 空间中看清。

\subsection{Hilbert 投影和 Riesz 表示}
\label{g3:subsec:1201:01}
在完备内积空间中，最短距离问题给出投影，投影又把一个有界线性泛函表示为内积。以下实、复空间都可使用；复内积 \(\langle u,v\rangle\) 对第一变量线性、第二变量共轭线性，范数为 \(\|u\|^2=\langle u,u\rangle\)。

\begin{theorem}\label{g3:thm:hilbert-projection}
	设 \(H\) 为实或复 Hilbert 空间，\(K\subset H\) 为非空闭凸集。对每个 \(x\in H\)，存在唯一 \(y=P_Kx\in K\) 使 \(\|x-y\|=\inf_{z\in K}\|x-z\|\)。它由条件
	\[
		\operatorname{Re}\langle x-y,z-y\rangle\le0\qquad(z\in K)
	\]
	刻画。投影满足 \(\|P_Kx-P_Kx'\|\le\|x-x'\|\)。若 \(K=V\) 是闭线性子空间，则 \(H=V\oplus V^\perp\)，\(P_V\) 线性且范数不超过一。
\end{theorem}
\begin{proof}
	令 \(d=\inf_K\|x-z\|\)，取极小化列 \(z_j\)。中点仍在 \(K\)，平行四边形恒等式给出
	\[
		\|z_j-z_k\|^2
		=2\|x-z_j\|^2+2\|x-z_k\|^2
		-4\left\|x-\frac{z_j+z_k}{2}\right\|^2
		\le2\|x-z_j\|^2+2\|x-z_k\|^2-4d^2.
	\]
	右侧趋零，故该列在完备空间中收敛；闭性使极限 \(y\) 属于 \(K\)，且达到距离 \(d\)。同一恒等式用于两个极小点，给出唯一性。

	对 \(z\in K\)，把 \(y+t(z-y)\)、\(0\le t\le1\) 代入最短距离条件，展开平方、除以 \(t>0\)，再令 \(t\downarrow0\)，得到所列不等式。反之，由
	\[
		\|x-z\|^2=\|x-y\|^2+\|z-y\|^2
		-2\operatorname{Re}\langle x-y,z-y\rangle
	\]
	即可恢复最短性。分别对 \(x,x'\) 的投影使用该条件，相加得到
	\(\|y-y'\|^2\le\operatorname{Re}\langle x-x',y-y'\rangle\)，再用 Cauchy--Schwarz 不等式便得非扩张性。

	若 \(K=V\) 线性，取 \(z=y+tv\) 的正、负实数 \(t\)，复情形再取 \(tv\) 的纯虚倍数，得到 \(x-y\perp V\)。正交分解的唯一性说明投影保持加法和数乘；范数估计已包含在非扩张性中。
\end{proof}

\begin{theorem}\label{g3:thm:hilbert-riesz}
	设 \(H\) 为实或复 Hilbert 空间。每个有界线性泛函 \(\ell:H\to\mathbb K\) 都唯一写成
	\(\ell(v)=\langle v,f\rangle\)，其中 \(f\in H\)、\(\|f\|=\|\ell\|\)。相应地，每个有界共轭线性泛函 \(F\) 都唯一写成 \(F(v)=\langle f,v\rangle\)。实空间中两种表述相同。
\end{theorem}
\begin{proof}
	若 \(\ell=0\)，取 \(f=0\)。否则 \(V=\ker\ell\) 为真闭子空间。选 \(x\notin V\)，令 \(e=x-P_Vx\ne0\)。对任意 \(v\)，
	\(v-\ell(v)e/\ell(e)\in V\)，因此
	\[
		\langle v,e\rangle=\frac{\ell(v)}{\ell(e)}\|e\|^2.
	\]
	取 \(f=\overline{\ell(e)}e/\|e\|^2\)，便有 \(\ell(v)=\langle v,f\rangle\)。Cauchy--Schwarz 给出 \(\|\ell\|\le\|f\|\)，以 \(v=f/\|f\|\) 测试给出反向不等式。两个表示之差对所有 \(v\) 配对为零，取 \(v\) 为该差即得唯一性。对共轭线性 \(F\)，将前述结论用于 \(\overline F\) 后取共轭即可。
\end{proof}

这个结论称为 Riesz 表示定理。\footnote{Frigyes Riesz（1880--1956），匈牙利数学家，是泛函分析的奠基者之一，研究函数空间、积分表示与算子。}在复空间中区分线性对偶 \(H^*\) 和共轭线性对偶 \(H^\dagger\)，可以使弱方程的右端与测试变量的共轭线性一致；这只是配对的记法，范数与可解性不受影响。
\symidx{\(H^\dagger\)：复 Hilbert 空间的有界共轭线性对偶}

\begin{example}\label{g3:exm:hilbert-mean-projection}
	在体积为 \(V>0\) 的紧黎曼流形上，常数函数构成 \(L^2\) 的一维闭子空间，其正交投影是 \(u\mapsto u_M=V^{-1}\int_Mu\,\mu_g\)。因为 \(\int_M(u-u_M)\,\mu_g=0\)，零均值空间正是常数的正交补。第十一章零均值 Poincaré 不等式中的归一化因而也是一个正交投影。
\end{example}

\subsection{Lax–Milgram 的证明}
\label{g3:subsec:1201:02}
积分型方程一般写成 \(a(u,v)=F(v)\)。即使 \(a\) 不对称，只要它在同一向量的配对上控制范数，也可以求解。对称性只在把解解释为二次能量的极小点时才需要。

\begin{theorem}\label{g3:thm:lax-milgram}
	设 \(H\) 为实或复 Hilbert 空间，\(a:H\times H\to\mathbb K\) 对第一变量线性、第二变量共轭线性，且存在 \(C<\infty\)、\(\alpha>0\) 满足
	\[
		|a(u,v)|\le C\|u\|\|v\|,\qquad
		\operatorname{Re}a(u,u)\ge\alpha\|u\|^2.
	\]
	则对每个 \(F\in H^\dagger\)，存在唯一 \(u\in H\) 使
	\[
		a(u,v)=F(v)\quad(v\in H),\qquad
		\|u\|\le\alpha^{-1}\|F\|.
	\]
	若另有 \(a(u,v)=\overline{a(v,u)}\)，此解是
	\(J(v)=\tfrac12a(v,v)-\operatorname{Re}F(v)\) 的唯一极小点。
\end{theorem}
\begin{proof}
	由\cref{g3:thm:hilbert-riesz}，定义有界线性算子 \(A:H\to H\) 及 \(f\in H\)，使
	\(a(u,v)=\langle Au,v\rangle\)、\(F(v)=\langle f,v\rangle\)。强制性与 Cauchy--Schwarz 给出 \(\|Au\|\ge\alpha\|u\|\)。所以 \(A\) 单射，且像闭：若 \(Au_j\) 收敛，该估计使 \(u_j\) 为 Cauchy 列，连续性使其极限仍映到指定极限。

	若 \(z\perp\operatorname{Ran}A\)，则 \(a(v,z)=0\) 对所有 \(v\) 成立；取 \(v=z\)，强制性给出 \(z=0\)。由闭子空间的正交分解，像既闭又具有零正交补，只能等于 \(H\)。因此 \(Au=f\) 唯一可解，范数界也由同一强制估计得到。

	在 Hermitian 情形，对任意 \(h\in H\)，利用方程展开可得
	\[
		J(u+h)-J(u)=\tfrac12a(h,h)\ge\tfrac\alpha2\|h\|^2.
	\]
	这给出唯一极小性。反之，极小点沿 \(th\) 的实方向导数为零，给出 \(\operatorname{Re}(a(u,h)-F(h))=0\)；复情形再取 \(\mathrm ih\) 即得到整个复数等式。
\end{proof}

第一条估计称为连续性，第二条称为\term[qiangzhixing]{强制性}[coercivity]。上述结果是 Lax--Milgram 定理。\footnote{Peter Lax（1926--2025），匈牙利出生的美国数学家，研究偏微分方程、计算数学与积分系统；Arthur N. Milgram（1912--1961），美国数学家，研究拓扑与分析。两人共同建立了这一弱方程可解性结果。}强制性是一个带统一正下界的条件，不能仅用 \(a(u,u)>0\)（\(u\ne0\)）代替。

\begin{example}\label{g3:exm:positive-not-coercive}
	在 \(\ell^2\) 上令 \(Ae_j=j^{-1}e_j\)。则 \(a(u,v)=\langle Au,v\rangle\) 连续、Hermitian 且对每个非零 \(u\) 有 \(a(u,u)>0\)，但其强制常数只能为零。向量 \(f=(1/j)_{j\ge1}\) 属于 \(\ell^2\)，方程 \(Au=f\) 却要求每个 \(u_j=1\)，所以没有 \(\ell^2\) 解。高频方向的能量越来越弱，使正性不足以保证满射。
\end{example}

若 \(a\) 含反 Hermitian 部分，\(\operatorname{Re}a(v,v)\) 只看见其 Hermitian 部分。直接极小化这个实能量一般会丢掉方程中的反对称项；Lax--Milgram 的存在性证明保留了整个算子。Hilbert 投影与上述可解性可对照 Brezis 的第5.1—5.3节。\cite[定理5.2、5.5，推论5.8]{Brezis2011}

\subsection{约束最小化与弱下半连续}
\label{g3:subsec:1201:03}
约束可能是均值为零、取值在一个凸集内，或满足指定的线性条件。弱收敛能保留哪一种约束，决定了极小化列的极限是否仍可使用。

\begin{lemma}\label{g3:lem:hilbert-weak-minimization-tools}
	设 \(H\) 为实或复 Hilbert 空间。每个有界序列都有弱收敛子列；每个范数闭凸集在弱拓扑下也闭。若 \(J:H\to(-\infty,+\infty]\) 为凸且范数下半连续的泛函，则它弱下半连续。
\end{lemma}
\begin{proof}
	对给定有界序列，取它所张成的闭线性子空间 \(V\)。有理系数有限组合（复情形同时取有理实、虚部）给出可数稠密集。Riesz 表示说明 \(V^*\) 也可分。对其单位球的可数稠密集逐次选取标量配对收敛的子列，再取对角列；统一范数界使配对极限延伸为 \(V^*\) 上的有界线性泛函。Riesz 表示使用两次说明 \(V\) 反射，即该极限由某个 \(v\in V\) 的求值给出。在 \(H\) 中测试时，可先将表示测试泛函的向量投影到 \(V\)，故所得子列也弱收敛于 \(H\)。这与\cref{g3:thm:sobolev-weak-compactness}中可分对偶的对角论证相同。

	若 \(K\) 非空、闭且凸，\(x\notin K\)，令 \(y=P_Kx\)。投影不等式给出
	\[
		\operatorname{Re}\langle x-y,z-y\rangle\le0\quad(z\in K),
		\qquad \operatorname{Re}\langle x-y,x-y\rangle>0.
	\]
	一个连续实线性泛函于是将 \(x\) 与 \(K\) 分开，其严格半空间是弱开邻域。这证明 \(K\) 弱闭；空集无须另证。\(J\) 的每个下水平集 \(\{J\le c\}\) 都范数闭且凸，故弱闭，这正是弱下半连续性。
\end{proof}

\begin{theorem}\label{g3:thm:direct-method-hilbert}
	设 \(H\) 为实或复 Hilbert 空间，\(K\subset H\) 非空且弱序列闭，\(J:K\to(-\infty,+\infty]\) 弱序列下半连续。假定 \(-\infty<\inf_KJ<\infty\)，并且至少有一个极小化列在 \(H\) 中有界。则 \(J\) 在 \(K\) 上达到下确界。若 \(K\) 凸且 \(J\) 严格凸，则极小点唯一。
\end{theorem}
\begin{proof}
	从有界极小化列提取弱收敛子列 \(u_j\rightharpoonup u\)。约束的弱序列闭性给出 \(u\in K\)，而
	\[
		\inf_KJ\le J(u)\le\liminf_jJ(u_j)=\inf_KJ.
	\]
	所以 \(u\) 为极小点。若有两个不同极小点，严格凸性使其中点的能量严格小于下确界，矛盾。
\end{proof}

\begin{corollary}\label{g3:cor:quadratic-convex-constraint}
	在\cref{g3:thm:lax-milgram}的 Hermitian 条件下，若 \(K\subset H\) 非空、范数闭且凸，则 \(J(v)=\tfrac12a(v,v)-\operatorname{Re}F(v)\) 在 \(K\) 上有唯一极小点。它由变分不等式
	\[
		\operatorname{Re}\bigl(a(u,v-u)-F(v-u)\bigr)\ge0\qquad(v\in K)
	\]
	刻画。
\end{corollary}
\begin{proof}
	强制性给出 \(J(v)\ge\alpha\|v\|^2/2-\|F\|\|v\|\)，故下确界有限且极小化列有界。\(J\) 为连续严格凸泛函，上一引理与定理给出存在唯一性。沿线段 \(u+t(v-u)\) 求右导数，得到必要的变分不等式。反之，展开
	\(J(v)-J(u)=\operatorname{Re}(a(u,v-u)-F(v-u))+a(v-u,v-u)/2\)，即知其非负。
\end{proof}

在障碍条件 \(u\ge\psi\) 下，可行集合通常是凸集，极小点满足变分不等式；只有允许向两个方向变化的测试才产生等式。另一类约束如 \(\|u\|_{L^2}=1\) 并非凸集。闭流形上若极小化列在 \(H^1\) 中有界，\cref{g3:thm:compact-bundle-rellich}给出强 \(L^2\) 子列，才可保证这一归一化在极限中保持。无限维 Hilbert 空间的正交单位向量列弱趋零，已经表明仅凭弱收敛不能保住单位长度。

\subsection{Poisson 方程的存在与唯一性}
\label{g3:subsec:1201:04}
先给出两个可解性模型：有界区域上的零边界问题，以及闭流形上去掉常数核后的问题。这里的 Laplace 算子始终采用 \(\Delta=-\operatorname{div}\operatorname{grad}\)。Poisson 方程的右端因而与非负 Dirichlet 能量配对。\footnote{Siméon Denis Poisson（1781--1840），法国数学家，研究势论、力学与概率；带源项的 Laplace 方程以其命名。}

\begin{definition}\label{g3:def:energy-dual}
	设 \(\Omega\subset\mathbb R^n\) 为有界开集。记
	\[
		H^{-1}(\Omega)=\bigl(H_0^1(\Omega)\bigr)^\dagger,
		\qquad \|F\|_{H^{-1}}=\sup_{\|v\|_{H^1}\le1}|F(v)|.
	\]
	实空间取通常的线性对偶。若 \(F\in H^{-1}\)，函数 \(u\in H_0^1(\Omega)\) 满足
	\(\int_\Omega\langle\mathrm Du,\mathrm Dv\rangle\,\mathrm dx=F(v)\) 对每个 \(v\in H_0^1\) 成立，就称它是零边界 Poisson 方程的弱解。\(f\in L^2\) 所确定的右端是 \(F(v)=\int f\overline v\)。
\end{definition}
\symidx{\(H^{-1}(\Omega)\)：零边界 \(H^1\) 空间的能量对偶}

\begin{theorem}\label{g3:thm:weak-poisson-dirichlet}
	设 \(\Omega\subset\mathbb R^n\) 为非空有界开集，\(n\ge1\)，\(F\in H^{-1}(\Omega)\)。则零边界 Poisson 方程有唯一弱解 \(u\in H_0^1(\Omega)\)，且
	\[
		\|u\|_{H^1}\le C_\Omega\|F\|_{H^{-1}}.
	\]
	它也是 \(J(v)=\tfrac12\int_\Omega|\mathrm Dv|^2-\operatorname{Re}F(v)\) 在 \(H_0^1\) 上的唯一极小点。
\end{theorem}
\begin{proof}
	第十一章定义的 \(H_0^1\) 是 \(H^1\) 的闭子空间，因而完备。\Cref{g3:thm:poincare-zero-boundary}给出 \(\|v\|_2\le C_P\|\mathrm Dv\|_2\)，从而
	\[
		\int_\Omega|\mathrm Dv|^2\ge(1+C_P^2)^{-1}\|v\|_{H^1}^2.
	\]
	能量型连续且 Hermitian，故\cref{g3:thm:lax-milgram}直接给出全部结论。这里零边界由闭包定义，对任意有界开集都有意义；当边界为 \(C^1\) 时，\cref{g3:thm:zero-trace-characterization}进一步将它识别为迹为零。
\end{proof}

若指定一个已给定的延拓 \(w\in H^1(\Omega)\)，则令 \(u=w+z\)、\(z\in H_0^1\)，把右端改为 \(F(v)-\int\langle\mathrm Dw,\mathrm Dv\rangle\) 即得到唯一的仿射边界类解。若 \(w\) 换成同一模 \(H_0^1\) 类中的另一个代表，所得仿射空间相同，解也相同。

\begin{theorem}\label{g3:thm:weak-poisson-closed}
	设 \((M,g)\) 为非空、连通、紧且无边界的光滑黎曼流形，\(F\in(H^1(M))^\dagger\)。方程
	\[
		\int_M\langle\mathrm du,\mathrm dv\rangle\,\mu_g=F(v)
		\qquad(v\in H^1(M))
	\]
	有弱解当且仅当 \(F(1)=0\)。满足 \(\int_Mu\,\mu_g=0\) 的解唯一，且
	\(\|u\|_{H^1}\le C\|F\|_{(H^1)^\dagger}\)；所有其他解与它相差常数。
\end{theorem}
\begin{proof}
	测试常数一立即给出必要条件。反过来，在闭子空间
	\(V=\{v\in H^1(M):\int_Mv\,\mu_g=0\}\) 上，\cref{g3:thm:poincare-compact-connected}使梯度能量强制，Lax--Milgram 给出唯一的 \(u\in V\)。每个测试函数分解为零均值部分与常数；\(F(1)=0\) 使这两部分都满足方程。若两解之差为 \(h\)，以 \(h\) 测试得到 \(\mathrm dh=0\)，\cref{g3:prop:weak-gradient-constant}及连通性说明 \(h\) 为常数。
\end{proof}

\begin{example}\label{g3:exm:circle-poisson}
	在标准圆周 \(\mathbb R/(2\pi\mathbb Z)\) 上，\(f(x)=\sin x\) 的平均为零，解为 \(u(x)=\sin x+C\)，因为 \(-u''=\sin x\)。若 \(f\equiv1\)，对常数测试便得到 \(0=2\pi\)，故无解。同一微分表达式是否可解，已经取决于整体的常数核及右端的配对。
\end{example}

更一般地，在 \(H_0^1(\Omega;\mathbb C^r)\) 上若积分型
\(a(u,v)=\int\langle A(x)\mathrm Du,\mathrm Dv\rangle\) 的矩阵 \(A\) 有界，且对所有 \(Z\in\mathbb C^{nr}\) 满足 \(\operatorname{Re}\langle A(x)Z,Z\rangle\ge\lambda|Z|^2\)，同一证明给出系统的弱可解性。这里要求的是整个梯度矩阵上的正性。下一节将说明主符号的椭圆性是一种不同的条件，而之后的局部估计负责把这些弱解提升到更高正则性。
% 纲要条目排为编号小节标题，并进入目录与书签；说明段仍保留于纲要注释中。

% !TeX root = ../../Book3.tex
% 纲要标题：### 12.2 主符号和常系数估计
% 纲要位置：几何学计划纲要.md，第 2561 行。
% 正文状态：初稿；尚未正式审核。

% 纲要内容开始
% 
% * 微分算子、形式伴随与主符号
% * 椭圆性、强椭圆性及 Laplace 型系统
% * Fourier 变换与 Sobolev 空间
% * 常系数 L² 估计和微分阶数
% 
% 纲要内容结束

\section{主符号和常系数估计}
\label{g3:sec:1202}

一个二阶方程为什么能够控制全部二阶导数？对常系数算子，振荡函数 \(e^{\mathrm ix\cdot\xi}v\) 把求导变成乘以频率 \(\xi\)，于是这个问题化为一个矩阵在每个非零频率处是否可逆。主符号记录最高阶的这部分矩阵；Fourier 变换把逐频率的代数估计转回函数范数。

\subsection{微分算子、形式伴随与主符号}
\label{g3:subsec:1202:01}
令 \(E,F\) 为光滑流形 \(M\) 上的有限秩实或复向量丛。阶数不超过整数 \(m\) 的线性微分算子 \(P:C^\infty(E)\to C^\infty(F)\)，在坐标与标架中写为
\[
	P=\sum_{|\alpha|\le m}A_\alpha(x)\mathrm D^\alpha,
	\qquad \mathrm D_j=-\mathrm i\partial_j,
	\qquad \mathrm D^\alpha=\prod_j\mathrm D_j^{\alpha_j}.
\]
这里 \(\alpha=(\alpha_1,\ldots,\alpha_n)\) 是非负整数多重指标，\(|\alpha|=\sum\alpha_j\)。实算子作复化后也采用此写法，因而某些 \(A_\alpha\) 可含 \(\mathrm i\)。本节中带下标或多重指标的 \(\mathrm D_j,\mathrm D^\alpha\) 固定表示这种 Fourier 归一化；\(\mathrm Du\) 作为完整梯度时仍指普通 Euclid 导数，两个版本的逐分量模长相同。

\begin{definition}\label{g3:def:principal-symbol}
	设 \(M\) 为光滑流形，\(E,F\to M\) 为有限秩实或复光滑向量丛，\(m\ge0\) 为整数。设 \(P:C^\infty(M;E)\to C^\infty(M;F)\) 在坐标及标架下写成 \(P=\sum_{|\alpha|\le m}A_\alpha(x)\mathrm D^\alpha\)，其中 \(\mathrm D_j=-\mathrm i\partial_j\)（实算子先复化）。其\term[zhufuhao]{主符号}[principal symbol]为余切向量 \(\xi\in\mathrm T_x^*M\) 上的齐次丛映射
	\[
		\sigma_m(P)(x,\xi)=\sum_{|\alpha|=m}A_\alpha(x)\xi^\alpha:E_x\longrightarrow F_x.
	\]
	它在 \(\xi\) 上为 \(m\) 次齐次；若实际阶数低于 \(m\)，此符号为零。
\end{definition}
\symidx{\(\sigma_m(P)(x,\xi)\)：\(m\) 阶微分算子的主符号}

\begin{proposition}\label{g3:prop:symbol-calculus}
	设 \(M\) 为光滑无边界流形，\(E,F,G\to M\) 为有限秩实或复光滑向量丛，\(P:C^\infty(E)\to C^\infty(F)\)、\(Q:C^\infty(F)\to C^\infty(G)\) 为阶数分别不超过整数 \(m\ge1,l\ge0\) 的光滑微分算子。按\cref{g3:def:principal-symbol}定义的主符号不依赖坐标和标架，且有
	\[
		\sigma_{m+l}(QP)=\sigma_l(Q)\sigma_m(P).
	\]
	固定光滑正密度和两丛的纤维度量后，唯一的\term[xingshi bansui]{形式伴随}[formal adjoint] \(P^*\) 满足
	\[
		\int_M\langle Pu,v\rangle\,\mu
		=\int_M\langle u,P^*v\rangle\,\mu
	\]
	对紧支光滑 \(u,v\) 成立，且 \(\sigma_m(P^*)=\sigma_m(P)^*\)。对标量光滑函数 \(\chi\)，交换子 \([P,\chi]=P\chi-\chi P\) 的阶数至多为 \(m-1\)。
\end{proposition}
\begin{proof}
	在 \(x\) 附近取实函数 \(\phi\) 满足 \(\mathrm d\phi_x=\xi\)，截面 \(s\) 满足 \(s(x)=v\)。逐次求导给出
	\[
		\lim_{t\to+\infty}t^{-m}e^{-\mathrm it\phi(x)}P(e^{\mathrm it\phi}s)(x)
		=\sigma_m(P)(x,\xi)v.
	\]
	只有每次导数都落在指数上的项保留 \(t^m\)；其他项至多含 \(t^{m-1}\)。左侧有坐标无关的意义，右侧只依赖 \(\mathrm d\phi_x,s(x)\)，因此得到内在的丛映射。复合时同样只保留全部最高阶导数，给出符号乘法。

	在一张坐标图中令 \(\mu=\rho|\mathrm dx|\)，以正定 Hermitian 矩阵 \(h_E,h_F\) 写纤维度量。由于 \(\mathrm D_j=-\mathrm i\partial_j\) 对普通 \(L^2\) 配对形式自伴，分部积分得到
	\[
		P^*v=\rho^{-1}h_E^{-1}\sum_{|\alpha|\le m}
		\mathrm D^\alpha\bigl(\rho A_\alpha^*h_Fv\bigr).
	\]
	最高阶项的系数是 \(h_E^{-1}A_\alpha^*h_F\)，正是纤维伴随；导数落在系数或密度上的项都降阶。配对恒等式保证这些局部表达式相容。唯一性由紧支测试截面的任意性得到。有边界时这里只测试内部紧支截面，故尚未引入伴随算子的边界定义域。

	最后，由多重指标乘积公式，\(P(\chi u)\) 中所有导数都落在 \(u\) 上的项恰等于 \(\chi Pu\)，剩余项至少有一次导数落在 \(\chi\)，所以对 \(u\) 的阶数至多为 \(m-1\)。
\end{proof}

\begin{example}\label{g3:exm:symbol-laplace}
	设 \((M^n,g)\) 为光滑黎曼流形，局部体积密度为 \(\rho\,|\mathrm dx|\)，其中 \(\rho=\sqrt{\det(g_{ij})}\)。函数上的非负 Laplace 算子为
	\[
		\Delta_g=-\rho^{-1}\partial_i(\rho g^{ij}\partial_j),
		\qquad \sigma_2(\Delta_g)(x,\xi)=g^{ij}\xi_i\xi_j=|\xi|_g^2.
	\]
	联络只在协变导数中增加零阶矩阵，因此度量联络的 \(\nabla^*\nabla\) 在任意有限秩丛上具有主符号 \(|\xi|_g^2I_E\)。两者的低阶系数不同，却具有相同的最高阶正性。
\end{example}

\subsection{椭圆性、强椭圆性及 Laplace 型系统}
\label{g3:subsec:1202:02}
符号可逆使每个非零频率都能由方程右端恢复；能量正性还要比较符号与纤维内积。这是两个相关而不同的要求。

\begin{definition}\label{g3:def:elliptic-system}
	设 \(M\) 为光滑流形，\(E,F\to M\) 为有限秩实或复光滑向量丛，\(P:C^\infty(M;E)\to C^\infty(M;F)\) 是光滑系数、阶数为 \(m\ge1\) 的微分算子，\(E,F\) 秩相同。若每个 \(x\) 和每个实非零余切向量 \(\xi\) 都使 \(\sigma_m(P)(x,\xi)\) 可逆，就称 \(P\) 为\term[tuoyuan suanzi]{椭圆算子}[elliptic operator]。在固定紧坐标区域上，这等价于存在 \(c>0\)，使
	\[
		|\sigma_m(P)(x,\xi)v|\ge c|\xi|^m|v|.
	\]
	若 \(E=F\)、\(m=2k\) 为偶数，且存在 \(\lambda>0\) 满足
	\[
		\operatorname{Re}\langle\sigma_{2k}(P)(x,\xi)v,v\rangle
		\ge\lambda|\xi|^{2k}|v|^2,
	\]
	就称它在此区域上\term[qiang tuoyuan]{强椭圆}[strongly elliptic]。若二阶算子的符号恰为 \(|\xi|_g^2I_E\)，称其为\term[Laplace xing suanzi]{Laplace 型算子}[Laplace-type operator]。
\end{definition}

紧集上的等价性来自 \((x,\xi,v)\) 在两个单位球面上的紧性。强椭圆性蕴含椭圆性；符号可逆也允许不定号。例如 \(\operatorname{diag}(-\sum\partial_i^2,\sum\partial_i^2)\) 的符号为 \(\operatorname{diag}(|\xi|^2,-|\xi|^2)\)，可逆却不正定。标量算子 \(\partial_1^2-\partial_2^2\) 则在 \(\xi_1=\xi_2\ne0\) 时符号为零，连椭圆性也不成立。

\ProseParagraphBreak

对一阶矩形系统，也常研究符号为单射的情形。此时 \(P^*P\) 的符号为 \(\sigma_1(P)^*\sigma_1(P)\)，在紧集上满足二阶强椭圆估计。以下将“椭圆系统”的主证明用于等秩可逆符号；矩形算子的结论须通过所写的单射条件或相应的椭圆复形另行判断。

\begin{proposition}\label{g3:prop:constant-strong-energy}
	设 \(P_0=\sum_{|\alpha|=2k}A_\alpha\mathrm D^\alpha\) 为 \(\mathbb R^n\) 上的齐次常系数系统，其中 \(n,r,k\ge1\)，\(A_\alpha\) 为 \(r\times r\) 复矩阵，\(\mathrm D_j=-\mathrm i\partial_j\)。假定存在 \(\lambda>0\)，使 \(\operatorname{Re}\langle\sum_{|\alpha|=2k}A_\alpha\xi^\alpha z,z\rangle\ge\lambda|\xi|^{2k}|z|^2\) 对所有 \(\xi\in\mathbb R^n,z\in\mathbb C^r\) 成立；Fourier 变换取 \((2\pi)^{-n/2}\) 归一化。则对 \(u\in C_c^\infty(\mathbb R^n;\mathbb C^r)\)，
	\[
		\operatorname{Re}\int\langle P_0u,u\rangle\,\mathrm dx
		\ge\lambda\int|\xi|^{2k}|\widehat u(\xi)|^2\,\mathrm d\xi.
	\]
	右侧与全部 \(k\) 阶导数的 \(L^2\) 范数平方之和等价。
\end{proposition}
证明将在下一小节建立 Fourier 等距性后直接给出。二阶情形中，这个条件只检验梯度矩阵 \(Z=\xi\otimes v\) 的特殊形状；上一节对整个 \(Z\) 的点态正性更强。可用能量方法构造弱解的系统，仍须按其实际积分型检查强制性。

\subsection{Fourier 变换与 Sobolev 空间}
\label{g3:subsec:1202:03}
\term[Fourier bianhuan]{Fourier 变换}[Fourier transform]把平移不变的求导运算变为频率乘法。\footnote{Jean Baptiste Joseph Fourier（1768--1830），法国数学家，研究热传导与三角级数，发展了以频率分解研究函数的方法。}我们采用保持 \(L^2\) 范数的归一化
\begin{equation}\label{g3:eq:fourier-normalization}
	\widehat u(\xi)=\mathcal Fu(\xi)
	=(2\pi)^{-n/2}\int_{\mathbb R^n}e^{-\mathrm ix\cdot\xi}u(x)\,\mathrm dx.
\end{equation}

\begin{definition}\label{g3:def:schwartz-tempered}
	所有对任意多重指标 \(\alpha,\beta\) 均满足 \(\sup_x|x^\alpha\partial^\beta u(x)|<\infty\) 的光滑函数组成 Schwartz 空间，记为 \(\mathcal S(\mathbb R^n)\)。\footnote{Laurent Schwartz（1915--2002），法国数学家，建立分布理论，并将其用于偏微分方程与 Fourier 分析。}在这些半范数下连续的线性泛函称为\term[huanzeng fenbu]{缓增分布}[tempered distribution]，其空间记为 \(\mathcal S'\)。连续性等价于存在有限个上述半范数，其和乘常数控制泛函的绝对值。
\end{definition}
\symidx{\(\mathcal S(\mathbb R^n),\mathcal S^\prime(\mathbb R^n)\)：Schwartz 函数与缓增分布空间}

紧支光滑函数属于 \(\mathcal S\)；高斯函数及其各阶导数也属于其中。若 \(u\in\mathcal S\)，积分求导和分部积分给出
\[
	\widehat{\mathrm D_j u}=\xi_j\widehat u,
	\qquad \widehat{x_ju}=\mathrm i\partial_{\xi_j}\widehat u,
	\qquad \widehat{u*v}=(2\pi)^{n/2}\widehat u\,\widehat v.
\]
各式的绝对可积性由快速衰减保证。将这两个求导恒等式迭代，所得 Fourier 半范数均由原函数有限个半范数控制，所以 \(\mathcal F:\mathcal S\to\mathcal S\) 连续。

\begin{theorem}\label{g3:thm:fourier-plancherel}
	在 \(\mathbb R^n\)、\(n\ge1\) 上，取 \(\widehat u(\xi)=(2\pi)^{-n/2}\int e^{-\mathrm ix\cdot\xi}u(x)\,\mathrm dx\)，其中 \(\mathcal S\) 为 Schwartz 空间，\(\mathcal S^\prime\) 为其连续对偶。若 \(u\in\mathcal S\)，则
	\[
		u(x)=(2\pi)^{-n/2}\int e^{\mathrm ix\cdot\xi}\widehat u(\xi)\,\mathrm d\xi,
		\qquad \|\widehat u\|_2=\|u\|_2.
	\]
	Fourier 变换唯一延拓为 \(L^2(\mathbb R^n;\mathbb C^r)\) 上的满射等距映射，保持内积；它及其逆也在 \(\mathcal S'\) 上连续定义，并与分布微分相容。
\end{theorem}
\begin{proof}
	先证明反演，而不预先使用 \(L^2\) 等距性。对 \(\varepsilon>0\) 插入因子 \(e^{-\varepsilon|\xi|^2}\)。一维高斯积分给出
	\[
		(2\pi)^{-n}\int e^{\mathrm iz\cdot\xi-\varepsilon|\xi|^2}\,\mathrm d\xi
		=(4\pi\varepsilon)^{-n/2}e^{-|z|^2/(4\varepsilon)}=:G_\varepsilon(z).
	\]
	这一公式可逐坐标验证：积分对 \(z\) 求导并对 \(\xi\) 分部积分，得到 \(I'(z)=-zI(z)/(2\varepsilon)\)；初值 \(I(0)=\sqrt{\pi/\varepsilon}\) 由将高斯积分平方后作平面极坐标积分得到。于是交换绝对收敛的积分可知
	\[
		(2\pi)^{-n/2}\int e^{\mathrm ix\cdot\xi-\varepsilon|\xi|^2}\widehat u(\xi)\,\mathrm d\xi
		=G_\varepsilon*u(x).
	\]
	右侧趋于 \(u(x)\)：\(G_\varepsilon\ge0\)、积分为一，任何固定小球外的质量趋零；在球内使用 \(u\) 的连续性，球外使用有界性。左侧由 \(\widehat u\in L^1\) 及支配收敛趋于反演积分。

	对 \(u,v\in\mathcal S\)，反演与积分换序给出
	\(\int\widehat u\,\overline{\widehat v}=\int u\overline v\)。例如先在 \(\int u\overline v\) 中以 \(\widehat v\) 的反演表示 \(v\)，被积绝对值由 \(|u(x)||\widehat v(\xi)|\) 控制，所得等式正是所求。取 \(u=v\) 即得等距性。\Cref{g3:lem:lp-smooth-density}保证 \(\mathcal S\) 在 \(L^2\) 中稠密，故有唯一等距延拓；其像闭，且由 \(\mathcal S\) 上的反演包含稠密的 \(\mathcal S\)，所以满射。有限维向量值情形逐分量相加。

	对 \(T\in\mathcal S'\) 定义 \(\widehat T(\phi)=T(\widehat\phi)\)。这采用双线性的分布测试配对，与函数的变换一致。\(\mathcal F\) 及反向变换在 \(\mathcal S\) 上连续互逆，转置后也互逆；对测试函数的分部积分给出 \(\widehat{\mathrm D_jT}=\xi_j\widehat T\)。
\end{proof}

等距性通常称为 Plancherel 定理。\footnote{Michel Plancherel（1885--1967），瑞士数学家，研究分析、数学物理与代数；Fourier 变换的平方可积等距定理以其命名。}

\begin{proof}[\cref{g3:prop:constant-strong-energy}的证明]
	由等距性和微分公式，左侧等于
	\(\int\operatorname{Re}\langle\sigma_{2k}(P_0)(\xi)\widehat u,\widehat u\rangle\,\mathrm d\xi\)。逐频率使用强椭圆条件，便得到所列下界。多项式 \(\sum_{|\alpha|=k}\xi^{2\alpha}\) 在单位球面上严格为正，且为 \(2k\) 次齐次，故与 \(|\xi|^{2k}\) 双向可比；再次用等距性即得导数范数比较。
\end{proof}

\termidx[shishujie Sobolev kongjian]{实数阶 Sobolev 空间}
\foreignidx{Sobolev space of real order}
\begin{definition}\label{g3:def:fourier-sobolev}
	设 \(n,r\ge1\)，对实数 \(s\)，记 \(\langle\xi\rangle=(1+|\xi|^2)^{1/2}\)。所有满足 \(\langle\xi\rangle^s\widehat u\in L^2\) 的缓增分布构成 \(H^s(\mathbb R^n;\mathbb C^r)\)，范数为
	\[
		\|u\|_{H^s}^2=\int\langle\xi\rangle^{2s}|\widehat u(\xi)|^2\,\mathrm d\xi.
	\]
	这里要求 \(\widehat u\) 由相应加权平方可积函数表示。对开集 \(U\)，\(H^s_{\mathrm{loc}}(U)\) 表示每个 \(\chi\in C_c^\infty(U)\) 都有 \(\chi u\in H^s(\mathbb R^n)\)。
\end{definition}
\symidx{\(H^s(\mathbb R^n)\)：以 Fourier 权定义的实数阶 Sobolev 空间}

\begin{lemma}\label{g3:lem:fourier-sobolev-tools}
	设 \(n,r\ge1\)、\(s\in\mathbb R\)，\(H^s=H^s(\mathbb R^n;\mathbb C^r)\) 为满足 \(\langle\xi\rangle^s\widehat u\in L^2\) 的缓增分布空间，\(\langle\xi\rangle=(1+|\xi|^2)^{1/2}\)。此空间完备，\(\mathcal S\) 在其中稠密。整数 \(s=k\ge0\) 时，它与第十一章的 \(H^k=W^{k,2}\) 相同，范数等价。另有以下性质。
	\begin{enumerate}
		\item 对 \(a\in C_c^\infty(\mathbb R^n)\)、实数 \(s\)，乘法 \(u\mapsto au\) 在 \(H^s\) 上有界。若 \(a\) 支于固定紧集，存在整数 \(N=N(n,s)\) 使乘子范数至多为 \(C\|a\|_{C^N}\)。矩阵值系数亦然。
		\item \(\mathrm D^\alpha:H^s\to H^{s-|\alpha|}\) 有界。若 \(m\ge1\)、\(\varepsilon>0\)，则
		\[
			\|u\|_{H^{s+m-1}}\le\varepsilon\|u\|_{H^{s+m}}+C_\varepsilon\|u\|_{H^s}.
		\]
		\item 每个紧支分布属于某个 \(H^{-N}\)。每个开集上的分布在相对紧子区域中因而具有某个有限的负 Sobolev 阶数。
	\end{enumerate}
\end{lemma}
\begin{proof}
	乘以 \(\langle\xi\rangle^s\) 后作 Fourier 变换，将 \(H^s\) 等距对应到 \(L^2\)；逆向得到的多项式增长加权 \(L^2\) 函数定义缓增分布。完备性随之成立。先用光滑紧支函数在频率空间逼近加权 \(L^2\) 函数，乘 \(\langle\xi\rangle^{-s}\) 后仍光滑紧支，逆变换属于 \(\mathcal S\)，这给出稠密性。整数阶比较来自
	\(\sum_{|\alpha|\le k}|\xi^\alpha|^2\asymp\langle\xi\rangle^{2k}\) 与分布微分公式。

	对乘法先取 \(u\in\mathcal S\)。由乘积的 Fourier 卷积公式和初等权重不等式
	\[
		\langle\xi\rangle^s\le C_s\langle\xi-\eta\rangle^{|s|}\langle\eta\rangle^s,
	\]
	得到
	\[
		\|au\|_{H^s}\le C_s\left\|\langle\zeta\rangle^{|s|}\widehat a(\zeta)\right\|_{L^1_\zeta}\|u\|_{H^s}.
	\]
	这里卷积的 \(L^1*L^2\to L^2\) 界由积分三角不等式和 \(L^2\) 的平移不变性给出。为估计 \(a\)，取整数 \(q\) 满足 \(2q>|s|+n\)，将 \((1-\sum\partial_j^2)^q\) 移到 \(a\) 上，便有
	\(|\widehat a(\zeta)|\le C\langle\zeta\rangle^{-2q}\|a\|_{C^{2q}}\)，常数只另依赖其固定支集体积。于是上述 \(L^1\) 范数有限，并有明确的有限阶系数界。由稠密性延拓到 \(H^s\)，同时在分布意义下仍为真实乘积。

	微分估计直接来自 \(|\xi^\alpha|\le\langle\xi\rangle^{|\alpha|}\)。插值估计将频率分成 \(\langle\xi\rangle\ge R\) 与其余部分：高频项用 \(R^{-1}\|u\|_{H^{s+m}}\) 控制，低频项用 \(R^{m-1}\|u\|_{H^s}\) 控制；选 \(R\ge\varepsilon^{-1}\) 即可。

	最后，设 \(T\) 支于紧集 \(K\)，取 \(\eta\in C_c^\infty\) 在其邻域为一。分布的有限阶估计使
	\[
		\widehat T(\xi)=(2\pi)^{-n/2}T\bigl(\eta(x)e^{-\mathrm ix\cdot\xi}\bigr),
		\qquad |\widehat T(\xi)|\le C\langle\xi\rangle^q
	\]
	对某个有限整数 \(q\) 成立。求 \(\xi\) 导数可移入测试函数，故它是光滑函数；把 \(T\) 与 Schwartz 测试函数的 Fourier 积分交换，可核对这确为分布变换。取 \(N>q+n/2\)，便有 \(\langle\xi\rangle^{-N}\widehat T\in L^2\)。局部情形先乘紧支截断，便归约于此。
\end{proof}

负阶空间在这里用来容纳尚未具有函数代表的分布。全空间的 \(H^{-1}\) 可由 \(H^1\) 的有界共轭线性对偶等距识别：Fourier 配对和加权 Cauchy--Schwarz 给出双向范数表示。这与上一节能量对偶的记法相容；在有界域上仍按 \((H_0^1)^\dagger\) 定义，不能直接把它等同于全空间中任意补零的分布。

\subsection{\texorpdfstring{常系数 \(L^2\) 估计和微分阶数}{常系数 L² 估计和微分阶数}}
\label{g3:subsec:1202:04}
现在逐频率解最高阶矩阵方程。非零频率处可逆，零频率则需要另一个低阶范数控制；这个区别也解释了局部估计中为什么常常出现 \(\|u\|\)。

\begin{theorem}\label{g3:thm:constant-elliptic-l2}
	设 \(P=\sum_{|\alpha|\le m}A_\alpha\mathrm D^\alpha\) 为 \(\mathbb R^n\) 上的 \(r\times r\) 常系数椭圆系统，\(m\ge1\)、\(n\ge1\)。则对每个实数 \(s\)，存在 \(C\) 使
	\begin{equation}\label{g3:eq:constant-elliptic-l2}
		\|u\|_{H^{s+m}}\le C\bigl(\|Pu\|_{H^s}+\|u\|_{H^s}\bigr).
	\end{equation}
	若分布 \(u\in H^s\) 且 \(Pu\in H^s\)，则 \(u\in H^{s+m}\)，并满足同一估计。齐次主部 \(P_m\) 还满足
	\[
		\sum_{|\alpha|=m}\|\partial^\alpha u\|_2^2\le C\|P_mu\|_2^2
		\qquad(u\in\mathcal S).
	\]
\end{theorem}
\begin{proof}
	写完整符号 \(p(\xi)=\sum A_\alpha\xi^\alpha=p_m(\xi)+p_{<m}(\xi)\)。由单位球面上的最小奇异值正下界，\(|p_m(\xi)v|\ge c|\xi|^m|v|\)。而低阶部分满足 \(|p_{<m}(\xi)v|\le C_0\langle\xi\rangle^{m-1}|v|\)，所以当 \(|\xi|\) 大于某个 \(R\) 时，\(|p(\xi)v|\ge(c/2)|\xi|^m|v|\)。小频率部分直接用 \(v\) 控制，两者合并得到
	\[
		\langle\xi\rangle^m|v|\le C\bigl(|p(\xi)v|+|v|\bigr).
	\]
	代入 \(v=\widehat u(\xi)\)，乘 \(\langle\xi\rangle^s\)，取 \(L^2\) 范数，即得第一式。对分布假设，\(\widehat u\) 及 \(p\widehat u\) 已具有加权 \(L^2\) 代表；同一逐点估计直接证明更高加权可积性，无须预设 \(u\) 的额外导数。齐次结论由 \(|p_m(\xi)v|\ge c|\xi|^m|v|\) 在全部非零频率处成立而得到。
\end{proof}

\begin{lemma}\label{g3:lem:constant-high-frequency-inverse}
	设 \(n,r,m\ge1\)，\(P_0=\sum_{|\alpha|=m}A_\alpha\mathrm D^\alpha\) 为 \(\mathbb R^n\) 上作用于 \(\mathbb C^r\) 值函数的常系数系统，\(\mathrm D_j=-\mathrm i\partial_j\)，\(p_m(\xi)=\sum A_\alpha\xi^\alpha\) 对 \(\xi\ne0\) 可逆。记 \(\mathcal F\) 为 \((2\pi)^{-n/2}\) 归一化的 Fourier 变换。固定 \(\chi\in C_c^\infty(\mathbb R^n)\)，在零的邻域为一。令
	\[
		q(\xi)=(1-\chi(\xi))p_m(\xi)^{-1},\qquad
		Q_0=\mathcal F^{-1}q\mathcal F,\qquad S_0=\mathcal F^{-1}\chi\mathcal F.
	\]
	则 \(q\) 在零附近按零延拓后光滑，且
	\[
		|\partial_\xi^\alpha q(\xi)|\le C_\alpha\langle\xi\rangle^{-m-|\alpha|},
		\qquad Q_0P_0=P_0Q_0=I-S_0.
	\]
	对每个实数 \(s\)，\(Q_0:H^s\to H^{s+m}\) 有界；对任意实数 \(s,t\)，\(S_0:H^s\to H^t\) 有界。\(S_0\) 还把任意缓增分布映成光滑函数。
\end{lemma}
\begin{proof}
	在单位球面上，逆矩阵光滑，求导公式为 \(\partial(p_m^{-1})=-p_m^{-1}(\partial p_m)p_m^{-1}\)。其齐次次数为 \(-m\)，每次求导再降低一次齐次次数，故远离零点有相应导数界；截断附近是固定紧环带，零附近恒为零，得到全空间估计。两个乘积恒等式直接在频率空间相乘即可验证。

	范数映射由 \(\langle\xi\rangle^m|q(\xi)|\le C\) 和 \(\chi\) 的紧支性给出。对一般缓增分布，\(\chi\widehat T\) 为紧支分布。其逆变换由它作用于 \(e^{\mathrm ix\cdot\xi}\) 乘一个固定紧支截断来表示，对 \(x\) 的任意次求导可移入这个测试函数，所以结果光滑。
\end{proof}

\begin{example}\label{g3:exm:laplace-frequency-gain}
	对 \(P=1-\sum_j\partial_j^2\)，完整符号就是 \(1+|\xi|^2\)。因此对每个实数 \(s\) 和 \(f\in H^s\)，
	\[
		\widehat u(\xi)=\frac{\widehat f(\xi)}{1+|\xi|^2}
	\]
	给出唯一的 \(u\in H^{s+2}\)，且 \(\|u\|_{H^{s+2}}=\|f\|_{H^s}\)。删掉常数一后，高频仍增益两阶，零频率处却没有同样的可逆界。
\end{example}

阶数 \(m\) 决定恢复多少导数，而不由方程分量数决定。下一节将固定一点的主系数，用这个常系数逆控制小球上的变系数误差；交换子降一阶，使截断引入的误差能够留在较低的 Sobolev 范数中。
% 纲要条目排为编号小节标题，并进入目录与书签；说明段仍保留于纲要注释中。

% !TeX root = ../../Book3.tex
% 纲要标题：### 12.3 变系数估计的局部化
% 纲要位置：几何学计划纲要.md，第 2568 行。
% 正文状态：初稿；尚未正式审核。

% 纲要内容开始
% 
% * 冻结系数与小球选择
% * 截断交换子和低阶项吸收
% * 差商方法与弱解的二阶正则性
% * 椭圆系统的同阶估计
% * 自举和分布解的光滑性
% 
% 纲要内容结束

\section{变系数估计的局部化}
\label{g3:sec:1203}

常系数估计控制每一个频率，变系数算子却会把不同频率混合。局部化的做法是在小球上固定最高阶系数，把它们的变化留作可吸收的扰动；截断产生的误差少一个导数，因而可以逐阶处理。还须区分先验估计与正则性：前者先对已有足够导数的函数计算，后者要证明一个弱解或分布确实拥有这些导数。本节先取得局部估计，再分别用差商和可逆的小扰动完成这一步。

\subsection{冻结系数与小球选择}
\label{g3:subsec:1203:01}
沿用上一节的记号，在开集 \(\Omega\subset\mathbb R^n\) 上写
\[
	P=\sum_{|\alpha|\le m}A_\alpha(x)\mathrm D^\alpha,
	\qquad \mathrm D_j=-\mathrm i\partial_j,
	\qquad p_m(x,\xi)=\sum_{|\alpha|=m}A_\alpha(x)\xi^\alpha.
\]
这里 \(m\ge1\)，各系数是光滑 \(N\times N\) 矩阵；标量情形取 \(N=1\)。若 \(P\) 椭圆，则在每个固定紧子集 \(K\Subset\Omega\) 上存在 \(c_K>0\)，使
\[
	|p_m(x,\xi)z|\ge c_K|\xi|^m|z|,
	\qquad x\in K,\quad \xi\in\mathbb R^n,\quad z\in\mathbb C^N.
\]
这是单位余向量与单位纤维向量上的紧性给出的统一下界。在一点 \(x_0\) 冻结系数，得到齐次常系数算子
\[
	P_0=\sum_{|\alpha|=m}A_\alpha(x_0)\mathrm D^\alpha.
\]
最高阶部分的差虽然系数小，却仍作用于 \(m\) 阶导数，所以必须先保证这个差能被左侧的同阶范数吸收。

\begin{lemma}\label{g3:lem:scaled-leading-perturbation}
	设 \(\Omega\subset\mathbb R^n\) 开，\(n,N,m\ge1\)，且 \(P=\sum_{|\alpha|\le m}A_\alpha(x)\mathrm D^\alpha\) 的系数为光滑 \(N\times N\) 矩阵。固定 \(x_0\in\Omega\)，取 \(\rho>0\)，使 \(\overline B_{3\rho}(x_0)\subset\Omega\)，并令 \(y=(x-x_0)/\rho\)。在 \(y\) 坐标中乘以 \(\rho^m\) 后，算子成为
	\[
		P_\rho=\sum_{|\alpha|\le m}\rho^{m-|\alpha|}
		A_\alpha(x_0+\rho y)\mathrm D_y^\alpha.
	\]
	固定 \(\zeta\in C_c^\infty(B_3)\)，使 \(\zeta=1\) 于 \(B_2\)。则在 \(B_2\) 上可写
	\[
		P_\rho=P_0+V_\rho+L_\rho,\qquad
		V_\rho=\sum_{|\alpha|=m}a_{\alpha,\rho}(y)\mathrm D_y^\alpha,
	\]
	其中 \(a_{\alpha,\rho}=\zeta(y)[A_\alpha(x_0+\rho y)-A_\alpha(x_0)]\) 紧支于 \(B_3\)，而 \(L_\rho\) 是系数也紧支于 \(B_3\) 的至多 \(m-1\) 阶算子。对每个固定实数 \(s\)，有
	\begin{equation}\label{g3:eq:scaled-leading-smallness}
		\|V_\rho v\|_{H^s}\le C_s\rho\,\|v\|_{H^{s+m}},
		\qquad
		\|L_\rho v\|_{H^s}\le C_s\|v\|_{H^{s+m-1}}
	\end{equation}
	对全空间 \(v\in H^{s+m}(\mathbb R^n;\mathbb C^N)\) 成立。这里 \(\rho\) 限于一个固定的小区间，常数由所用系数的有限阶导数界控制。
\end{lemma}
\begin{proof}
	在一个预先固定的小邻域中，对任意固定整数 \(q\ge0\)，乘积求导与微积分基本定理给出
	\[
		\max_{|\beta|\le q}\|\partial_y^\beta a_{\alpha,\rho}\|_\infty
		\le C_q\rho.
	\]
	零阶差由 \(A_\alpha(x_0+\rho y)-A_\alpha(x_0)\) 的线段积分估计；至少一次落在 \(A_\alpha\) 上的导数带因子 \(\rho^{|\beta|}\)，落在 \(\zeta\) 上的其余项仍乘以同一个零阶差。取 \(\rho\le1\)，便得到所写上界。低阶系数取
	\(\zeta(y)\rho^{m-|\alpha|}A_\alpha(x_0+\rho y)\)，它们的每个固定阶导数一致有界。

	由\cref{g3:lem:fourier-sobolev-tools}的光滑乘子估计，固定紧支集上的乘法算子范数受有限个这样的导数上确界控制；同时 \(\mathrm D^\alpha:H^{s+|\alpha|}\to H^s\) 连续。逐项相加就得到式\eqref{g3:eq:scaled-leading-smallness}。缩放把原坐标中未必小的系数导数一并变小，这是只写 \(\|A(x)-A(x_0)\|_\infty\ll1\) 所没有包含的信息。
\end{proof}

\begin{proposition}\label{g3:prop:frozen-coefficient-apriori}
	设 \(\Omega\subset\mathbb R^n\) 开，\(P\) 是光滑系数的椭圆 \(N\times N\) 系统，阶数为 \(m\ge1\)，并固定 \(s\in\mathbb R\)、\(x_0\in\Omega\)。存在 \(\rho>0\) 与 \(C>0\)，使对所有 \(v\in C_c^\infty(B_\rho(x_0);\mathbb C^N)\)，
	\begin{equation}\label{g3:eq:small-ball-elliptic-estimate}
		\|v\|_{H^{s+m}(\mathbb R^n)}
		\le C\bigl(\|Pv\|_{H^s(\mathbb R^n)}+\|v\|_{H^s(\mathbb R^n)}\bigr).
	\end{equation}
	系数只需在该球附近定义，式中的紧支函数及其算子像均在外部补零。
\end{proposition}
\begin{proof}
	先在上一引理的单位尺度上计算。对 \(P_0\) 应用\cref{g3:thm:constant-elliptic-l2}，得
	\[
	\begin{aligned}
		\|v\|_{H^{s+m}}
		&\le C_0\bigl(\|P_\rho v\|_{H^s}+\|V_\rho v\|_{H^s}
		+\|L_\rho v\|_{H^s}+\|v\|_{H^s}\bigr)\\
		&\le C_0\|P_\rho v\|_{H^s}+C_1\rho\|v\|_{H^{s+m}}
		+C_2\|v\|_{H^{s+m-1}}+C_0\|v\|_{H^s}.
	\end{aligned}
	\]
	先选 \(\rho\) 使 \(C_1\rho\le1/4\)，再用\cref{g3:lem:fourier-sobolev-tools}中的插值估计选择 \(\varepsilon\)，使
	\(C_2\|v\|_{H^{s+m-1}}\le\tfrac14\|v\|_{H^{s+m}}+C_\varepsilon\|v\|_{H^s}\)。将两项移到左侧，得到单位尺度上的结论。对已经固定的 \(\rho\)，变量替换在每个所用的 \(H^t\) 范数之间给出有界同构，因而变回 \(x\) 坐标即得式\eqref{g3:eq:small-ball-elliptic-estimate}。该常数允许依赖所选尺度，并未把非齐次范数当成齐次缩放量。
\end{proof}

\subsection{截断交换子和低阶项吸收}
\label{g3:subsec:1203:02}
将任意函数限制到小球，需要乘一个截断函数。由乘积法则，\(m\) 阶算子与截断的交换子只有至多 \(m-1\) 阶：
\begin{equation}\label{g3:eq:localized-commutator}
	[P,\eta]u=P(\eta u)-\eta Pu
	=\sum_{|\alpha|\le m}
	\sum_{0<\beta\le\alpha}\binom\alpha\beta
	A_\alpha(\mathrm D^\beta\eta)\mathrm D^{\alpha-\beta}u.
\end{equation}
这里 \(\eta\) 是标量光滑函数，因此它与矩阵系数可交换。每一项至少有一次求导落在 \(\eta\) 上，正是少掉一个 \(u\) 的导数的原因。

\begin{lemma}\label{g3:lem:cutoff-elliptic-apriori}
	设 \(P\) 是开集 \(\Omega\subset\mathbb R^n\) 上的光滑椭圆 \(N\times N\) 系统，阶数 \(m\ge1\)，\(s\in\mathbb R\)。给定 \(\chi,\psi\in C_c^\infty(\Omega)\)，且 \(\psi=1\) 于 \(\operatorname{supp}\chi\) 的某个邻域，则存在常数 \(C\)，使对每个光滑 \(u:\Omega\to\mathbb C^N\)，
	\begin{equation}\label{g3:eq:cutoff-elliptic-apriori}
		\|\chi u\|_{H^{s+m}}
		\le C\bigl(\|\psi Pu\|_{H^s}+\|\psi u\|_{H^{s+m-1}}\bigr).
	\end{equation}
	全部范数在局部截断后按全空间定义。
\end{lemma}
\begin{proof}
	用有限个满足\cref{g3:prop:frozen-coefficient-apriori}的小球覆盖 \(\operatorname{supp}\chi\)，在这些球中选光滑函数 \(\theta_a\)，使 \(\sum_a\theta_a=1\) 于该支集附近，且各支集都位于 \(\psi=1\) 的区域。对每个 \(v_a=\theta_a\chi u\) 使用式\eqref{g3:eq:small-ball-elliptic-estimate}。算子像为
	\[
		Pv_a=\theta_a\chi\,Pu+[P,\theta_a\chi]u.
	\]
	第一项由 \(\|\psi Pu\|_{H^s}\) 控制；第二项由式\eqref{g3:eq:localized-commutator}及光滑乘子估计，受 \(C\|\psi u\|_{H^{s+m-1}}\) 控制。零阶余项也受同一范数控制。最后用 \(\chi u=\sum_av_a\) 和有限项三角不等式即得结论。
\end{proof}

对二阶算子 \(P=-\sum_{i,j}\partial_i(a^{ij}\partial_j)\)，交换子也可以直接写成
\[
	[P,\eta]u
	=-\sum_{i,j}\partial_i(a^{ij}\partial_j\eta)u
	-\sum_{i,j}a^{ij}(\partial_j\eta)\partial_i u
	-\sum_{i,j}a^{ij}(\partial_i\eta)\partial_j u.
\]
当 \(a^{ij}=a^{ji}\) 时，最后两项合并。截断的导数支于两个同心球之间的区域；估计内球上的解时，须保留稍大球中的信息，如\cref{g3:fig:elliptic-localization}所示。
\begin{figure}[htbp]
	\centering
	% !TeX root = ../Book3.tex
\begin{tikzpicture}[scale=1.05,>=stealth]
	\fill[black!9] (0,0) circle (1.45);
	\fill[white] (0,0) circle (.75);
	\draw[thick] (0,0) circle (.75);
	\draw[thick] (0,0) circle (1.45);
	\draw[dashed] (0,0) circle (2.05);
	\fill (0,0) circle (1.2pt);
	\node[below] at (0,0) {\(x_0\)};
	\node at (0,.4) {\(\chi=1\)};
	\node[below left] at (-.52,-.55) {\(B_r\)};
	\node[below left] at (-1.04,-1.04) {\(B_R\)};
	\draw[->] (3,.9)--(1.15,.55);
	\node[right,align=left] at (3,.9) {\(\partial\chi\) 可能非零\\交换子至多 \(m-1\) 阶};
	\draw[->] (3,-.9)--(1.65,-.7);
	\node[right,align=left] at (3,-.9) {外层截断 \(\psi\)\\在 \(\operatorname{supp}\chi\) 附近为一};
\end{tikzpicture}

	\caption[椭圆估计中的嵌套截断]{内球上保留原函数，外侧过渡区产生低一阶交换子。再选外层截断，使其在整个内层支集附近等于一，便可继续估计这些误差。}
	\label{g3:fig:elliptic-localization}
\end{figure}

吸收低阶项时要区分两种情形。对同一个紧支函数，可以直接用全空间插值，把 \(H^{s+m-1}\) 项分成任意小的 \(H^{s+m}\) 项与 \(H^s\) 项；若这两个范数分别属于 \(\chi u\) 与 \(\psi u\)，则它们位于不同区域，不能直接移项。后面的分布正则性证明将通过有限串嵌套截断逐级提高阶数，并得到只含外层低阶范数的内部估计。

\subsection{差商方法与弱解的二阶正则性}
\label{g3:subsec:1203:03}
变分法给出的解通常只有一阶弱导数。如果直接把 \(\partial_k u\) 代入弱方程，测试函数未必具有所需的一阶导数。差商保留了有限平移，在每个非零步长下仍可构造合法的测试函数，再从统一估计恢复二阶导数。

\begin{lemma}\label{g3:lem:difference-quotient-criterion}
	设 \(B_r\Subset B_R\subset\mathbb R^n\) 为同心球，\(e_k\) 为第 \(k\) 个坐标向量。对足够小的实数 \(h\ne0\)，定义
	\[
		\delta_h^ku(x)=\frac{u(x+he_k)-u(x)}h.
	\]
	若 \(u\in H^1(B_R;\mathbb C^N)\)，则
	\(\|\delta_h^ku\|_{L^2(B_r)}\le\|\partial_ku\|_{L^2(B_R)}\) 对 \(|h|<R-r\) 成立。反之，若 \(u\in L^2(B_R;\mathbb C^N)\)，且这些差商在 \(L^2(B_r)\) 中一致有界，则 \(\partial_ku\in L^2(B_r)\)，其范数不超过该统一上界。
\end{lemma}
\begin{proof}
	第一项沿坐标线使用绝对连续代表，写成
	\[
		\delta_h^ku(x)=\int_0^1\partial_ku(x+the_k)\,\mathrm dt,
	\]
	再用积分平均的凸性与换序；这也正是\cref{g3:lem:sobolev-translation-bound}的局部证明。对第二项，由\cref{g3:thm:sobolev-weak-compactness}的 \(k=0,p=2\) 情形，从 \(h_j\to0\) 中选子列使差商弱收敛于 \(v\in L^2(B_r)\)。对 \(\phi\in C_c^\infty(B_r;\mathbb C^N)\)，平移换元给出
	\[
		\int\langle\delta_{h_j}^ku,\phi\rangle\,\mathrm dx
		=-\int\langle u,\delta_{-h_j}^k\phi\rangle\,\mathrm dx
		\longrightarrow-\int\langle u,\partial_k\phi\rangle\,\mathrm dx.
	\]
	所以 \(v\) 就是弱导数 \(\partial_ku\)，弱下半连续性给出范数界。
\end{proof}

\begin{theorem}\label{g3:thm:weak-interior-h2}
	设 \(\Omega\subset\mathbb R^n\) 开，\(u\in H^1_{\mathrm{loc}}(\Omega;\mathbb C^N)\)。设矩阵系数 \(A^{ij}\in W^{1,\infty}_{\mathrm{loc}}(\Omega)\)，\(b^j,c\in L^\infty_{\mathrm{loc}}(\Omega)\)，且在每个紧子区域上存在 \(\lambda>0\)，使
	\begin{equation}\label{g3:eq:full-gradient-coercivity}
		\operatorname{Re}\sum_{i,j=1}^n\langle A^{ij}(x)Z_j,Z_i\rangle
		\ge\lambda\sum_{j=1}^n|Z_j|^2
	\end{equation}
	对几乎处处 \(x\) 及任意 \(Z_1,\ldots,Z_n\in\mathbb C^N\) 成立。若 \(f\in L^2_{\mathrm{loc}}(\Omega;\mathbb C^N)\)，且 \(u\) 满足弱方程
	\begin{equation}\label{g3:eq:weak-system-divergence}
		-\sum_{i,j}\partial_i(A^{ij}\partial_ju)
		+\sum_jb^j\partial_ju+cu=f,
	\end{equation}
	则 \(u\in H^2_{\mathrm{loc}}(\Omega;\mathbb C^N)\)。对 \(B_r\Subset B_R\Subset\Omega\)，有
	\begin{equation}\label{g3:eq:interior-h2-estimate}
		\|u\|_{H^2(B_r)}
		\le C\bigl(\|f\|_{L^2(B_R)}+\|u\|_{L^2(B_R)}\bigr),
	\end{equation}
	其中 \(C\) 依赖于 \(n,N,r,R,\lambda\)、\(A^{ij}\) 的一阶有界导数范数及低阶系数的有界范数。
\end{theorem}
\begin{proof}
	先保留右侧的 \(H^1\) 范数，证明二阶导数存在；再用一次局部能量估计把它换成 \(L^2\) 范数。令
	\(F=f-\sum_jb^j\partial_ju-cu\)，则 \(F\in L^2\) 于任意相对紧子区域。取 \(0\le\eta\le1\) 的紧支光滑截断，在所求内球上等于一。对足够小的 \(h\)，令 \(z=\delta_h^ku\)，并在弱方程中取
	\[
		v=-\delta_{-h}^k(\eta^2z).
	\]
	固定 \(h\ne0\) 时，\(z\) 是两个 \(H^1\) 函数的差，故 \(v\in H^1_0\)，且支集仍位于原区域内部。第十一章的内部光滑逼近使弱方程可用于这个测试函数。

	使用上一引理证明中的离散分部积分，以及
	\[
		\delta_h^k(A^{ij}\partial_ju)
		=A^{ij}(x+he_k)\partial_jz+(\delta_h^kA^{ij})\partial_ju,
	\]
	得到
	\[
		\sum_{i,j}\int\left\langle
		A^{ij}(x+he_k)\partial_jz+(\delta_h^kA^{ij})\partial_ju,
		\partial_i(\eta^2z)\right\rangle\,\mathrm dx
		=\int\langle F,-\delta_{-h}^k(\eta^2z)\rangle\,\mathrm dx.
	\]
	令 \(X=\|\eta\partial z\|_2\)，其中 \(\partial z=(\partial_1z,\ldots,\partial_nz)\)。展开 \(\partial_i(\eta^2z)=\eta^2\partial_iz+2\eta(\partial_i\eta)z\)，主项的实部由式\eqref{g3:eq:full-gradient-coercivity}至少为 \(\lambda X^2\)。其余三类项分别含有 \(\delta_h^kA^{ij}\)、\(\partial_i\eta\)，或两者。由于
	\[
		\|\delta_h^kA^{ij}\|_\infty\le\|\partial_kA^{ij}\|_\infty,
		\qquad \|z\|_2\le\|\partial_ku\|_{L^2(\text{稍大区域})},
	\]
	这些项的绝对值之和不超过
	\(C\|u\|_{H^1}X+C\|u\|_{H^1}^2\)。这里所有原函数范数都取在一个固定的稍大区域，因而常数与 \(h\) 无关。

	右侧也无须对 \(F\) 求导。把 \(\eta^2z\) 补零为全空间 \(H^1\) 函数，由平移估计有
	\[
	\begin{aligned}
		\left|\int\langle F,-\delta_{-h}^k(\eta^2z)\rangle\right|
		&\le\|F\|_2\,\|\partial_k(\eta^2z)\|_2\\
		&\le C\|F\|_2\bigl(X+\|u\|_{H^1}\bigr).
	\end{aligned}
	\]
	由 \(ab\le\varepsilon a^2+(4\varepsilon)^{-1}b^2\) 逐项吸收 \(X\)，得到
	\[
		\|\eta\partial(\delta_h^ku)\|_2
		\le C\bigl(\|f\|_2+\|u\|_{H^1}\bigr).
	\]
	弱导数与平移交换，所以左侧控制每个 \(\delta_h^k\partial_ju\)。对有限个 \(j,k\) 应用差商判据，得到内球中的全部二阶弱导数以及相应估计。

	最后选内、外球之间的中间球，并取另一个截断 \(\theta\)，使其在中间球上为一、支于 \(B_R\)。将 \(v=\theta^2u\) 代入原弱方程，取实部。主项为正，其余项由系数界满足
	\[
	\begin{aligned}
		\lambda\|\theta\partial u\|_2^2
		\le{}&\|f\|_2\|u\|_2
		+C\|\theta\partial u\|_2\|u\|_2
		+C\|u\|_2^2.
	\end{aligned}
	\]
	截断导数的界已吸收到 \(C\) 中。再次吸收，得到中间球上
	\(\|u\|_{H^1}\le C(\|f\|_{L^2(B_R)}+\|u\|_{L^2(B_R)})\)。将此前的二阶估计应用于内球和中间球，便得式\eqref{g3:eq:interior-h2-estimate}。整个论证都不要求 \(u\) 在 \(\partial\Omega\) 上满足边界条件。
\end{proof}

当 \(A^{ij}=a^{ij}I_N\)，\(a=(a^{ij})\) 为实对称、一致正定矩阵时，式\eqref{g3:eq:full-gradient-coercivity}直接逐分量成立。这覆盖光滑 Laplace 型系统在坐标中的弱方程。Lipschitz 系数保证差商 \(\delta_hA^{ij}\) 一致有界；若系数跨越界面发生跳跃，解的一阶导数可以相应跳跃，二阶 \(L^2\) 正则性便可能失效。

\subsection{椭圆系统的同阶估计}
\label{g3:subsec:1203:04}
系统中的未知量是一个向量，最高阶矩阵可以把不同分量及不同微分方向耦合。冻结系数时应反转整个主符号矩阵，不能只逐个观察它的对角元素。前两小节的证明一直采用矩阵算子范数，因此已适用于同阶的等秩椭圆系统：每个未知分量用同一阶 \(m\)，每个方程也取这个共同的最高阶。

\begin{example}\label{g3:exm:coupled-elliptic-symbol}
	在 \(\mathbb R^2\) 上考虑两个复值未知函数的系统
	\[
		P\binom uv=
		\binom{-\partial_1^2u-\partial_2^2u+\partial_1\partial_2v}
		{-\partial_1^2v-\partial_2^2v+\partial_1\partial_2u}.
	\]
	其主符号为
	\[
		p_2(\xi)=
		\begin{pmatrix}|\xi|^2&-\xi_1\xi_2\\-\xi_1\xi_2&|\xi|^2\end{pmatrix}.
	\]
	在向量 \((1,1)\)、\((1,-1)\) 上的特征值分别为 \( |\xi|^2-\xi_1\xi_2\)、\( |\xi|^2+\xi_1\xi_2\)，均至少为 \( |\xi|^2/2\)。所以冻结估计可以同时控制 \(u,v\) 的所有二阶导数。把未知量换成 \(u+v,u-v\) 可在这个例子中对角化；一般变系数系统未必存在同一个固定的对角化，主符号矩阵的统一逆界仍足以给出局部估计。
\end{example}

\begin{proposition}\label{g3:prop:uniform-system-localization}
	设一族光滑 \(N\times N\) 的 \(m\) 阶系统 \(P_\tau\) 定义在同一个开集 \(\Omega\subset\mathbb R^n\) 上，固定实数 \(s\) 及截断 \(\chi,\psi\in C_c^\infty(\Omega)\)，其中 \(\psi=1\) 于 \(\operatorname{supp}\chi\) 附近。若在 \(\operatorname{supp}\psi\) 的固定邻域内，主符号满足统一下界
	\[
		|p_{m,\tau}(x,\xi)z|\ge c|\xi|^m|z|,
	\]
	且证明所需的有限阶系数导数有统一上界，则式\eqref{g3:eq:cutoff-elliptic-apriori}中的 \(C\) 可与 \(\tau\) 无关。所需导数阶数可只依赖于 \(n,m,s\)。
\end{proposition}
\begin{proof}
	常系数 Fourier 估计中的常数由 \(c^{-1}\) 及系数上界统一控制。\Cref{g3:lem:scaled-leading-perturbation}中 \(C_s\rho\) 的系数由同一组有限阶导数界控制，故所有成员可用同一个充分小的球半径。固定紧支集的有限球覆盖、单位分解和各阶截断导数也可一次选定。将这些统一界代入\cref{g3:lem:cutoff-elliptic-apriori}的证明，即得共同常数。
\end{proof}

前述差商定理采用了对任意梯度数组的正性条件\eqref{g3:eq:full-gradient-coercivity}。一般强椭圆性只要求主符号在单个实余向量上的正性，两者不能混同。例如在 \(\mathbb R^2\) 的向量场上，取
\[
	A^{ij}=\delta_{ij}I_2-\frac34 e_i\otimes e_j,
	\qquad
	p_2(\xi)=|\xi|^2I_2-\frac34\xi\otimes\xi.
\]
主符号特征值为 \( |\xi|^2\) 与 \( |\xi|^2/4\)，所以算子强椭圆；若梯度数组 \(Z_j=e_j\)，则
\[
	\sum_{i,j}\langle A^{ij}Z_j,Z_i\rangle=2-\frac34\,4=-1.
\]
因此不能把差商证明中的逐点梯度正性自动套到所有强椭圆系统。光滑同阶系统的内部正则性由主符号的可逆性得到，下面直接按这一条件证明。

\ProseParagraphBreak

在向量丛的平凡化中，联络和标架变化只产生光滑的低阶矩阵项。第十一章的范数等价\cref{g3:prop:sobolev-norm-equivalence}使局部整数阶估计与丛值 Sobolev 范数一致。若各方程和未知量具有不同最高阶，则需要另行指定权重；本节的同阶符号与估计只用于已明确的共同阶数。

\subsection{自举和分布解的光滑性}
\label{g3:subsec:1203:05}
光滑函数的先验估计尚不能直接用于一个没有 \(m\) 阶导数的分布。关键是把最高阶扰动移到左侧，并在两个不同的 Sobolev 阶数上反转同一个算子：较低阶空间保证解的识别，较高阶空间提供新增导数。

\begin{lemma}\label{g3:lem:two-scale-inverse}
	设 \(t_0<t_1\) 为实数，线性算子 \(T\) 在 \(H^{t_0}(\mathbb R^n;\mathbb C^N)\) 与 \(H^{t_1}(\mathbb R^n;\mathbb C^N)\) 上均有界，且两个作用在交集上相同。如果它在这两个空间上的算子范数都不超过 \(1/2\)，那么 \(I+T\) 在两空间上均可逆，逆算子范数均不超过二。若 \(v\in H^{t_0}\) 且 \((I+T)v=g\in H^{t_1}\)，则 \(v\in H^{t_1}\)，并有 \(\|v\|_{H^{t_1}}\le2\|g\|_{H^{t_1}}\)。
\end{lemma}
\begin{proof}
	在每个空间中，算子级数 \(\sum_{j=0}^\infty(-T)^j\) 绝对收敛。有限几何级数恒等式给出
	\[
		(I+T)\sum_{j=0}^{J}(-T)^j=I-(-T)^{J+1},
	\]
	右端按算子范数趋于 \(I\)，反向乘法也相同，故级数确为双侧逆。对 \(g\in H^{t_1}\)，先在较高阶空间得到解 \(w\)。连续包含 \(H^{t_1}\subset H^{t_0}\) 保证 \(w\) 在较低阶空间也是解；该空间的唯一性使 \(v=w\)。
\end{proof}

这里使用的几何级数称为\term[Neumann jishu]{Neumann 级数}[Neumann series]。\footnote{Carl Gottfried Neumann（1832--1925），德国数学家，研究位势论、积分方程与数学物理；这里的级数及 Neumann 边值问题以其命名。}两个空间中的逆来自同一列部分和，这一点保证了所构造的较光滑解就是原来的分布。

\begin{theorem}\label{g3:thm:elliptic-distribution-l2}
	设 \(\Omega\subset\mathbb R^n\) 开，\(P\) 是光滑系数的椭圆 \(N\times N\) 微分系统，阶数 \(m\ge1\)，\(s\in\mathbb R\)。若 \(u\in\mathcal D'(\Omega;\mathbb C^N)\) 且 \(Pu\in H^s_{\mathrm{loc}}(\Omega;\mathbb C^N)\)，则
	\[
		u\in H^{s+m}_{\mathrm{loc}}(\Omega;\mathbb C^N).
	\]
	此外，给定 \(t\in\mathbb R\)、\(\chi,\psi\in C_c^\infty(\Omega)\)，且 \(\psi=1\) 于 \(\operatorname{supp}\chi\) 的邻域，则对满足 \(\psi u\in H^t\)、\(\psi Pu\in H^s\) 的分布，有
	\begin{equation}\label{g3:eq:elliptic-distribution-estimate}
		\|\chi u\|_{H^{s+m}}
		\le C\bigl(\|\psi Pu\|_{H^s}+\|\psi u\|_{H^t}\bigr).
	\end{equation}
	常数依赖于两个截断、\(P,s,t\) 及所用的有限阶系数界。
\end{theorem}
\begin{proof}
	先说明一个小球中怎样增加至多一阶正则性。使用\cref{g3:lem:scaled-leading-perturbation}的单位尺度写法，固定一个当前已知的阶数 \(a\)，并设 \(a<s+m\)。令
	\(b=\min\{a+1,s+m\}>a\)。取\cref{g3:lem:constant-high-frequency-inverse}中的高频逆 \(Q_0\)，满足
	\[
		Q_0P_0=I-S_0,\qquad
		Q_0:H^q\longrightarrow H^{q+m}\quad(q\in\mathbb R),
	\]
	其中 \(S_0\) 为低频截断算子，从任意 \(H^q\) 映到任意 \(H^{q'}\)。由式\eqref{g3:eq:scaled-leading-smallness}，\(T_\rho=Q_0V_\rho\) 在 \(H^a,H^b\) 上的范数均为 \(O(\rho)\)。因此可选同一个 \(\rho\)，使两者都不超过 \(1/2\)。这些算子均由同一个 Fourier 乘子和微分、乘法运算定义，其在不同空间中的作用相容。

	记缩放后的分布为 \(u_\rho(y)=u(x_0+\rho y)\)，右端为 \(f_\rho(y)=\rho^m(Pu)(x_0+\rho y)\)。对分布，这个写法以变量替换定义，即
	\(\langle u_\rho,\phi\rangle=\rho^{-n}\langle u,\phi((x-x_0)/\rho)\rangle\)；于是 \(P_\rho u_\rho=f_\rho\)。取 \(\eta,\vartheta\in C_c^\infty(B_2)\)，使 \(\vartheta=1\) 于 \(\operatorname{supp}\eta\) 的邻域，令 \(v=\eta u_\rho\)。假设已知 \(\vartheta u_\rho\in H^a\)，则 \(v\in H^a\)。由分布乘积公式及 \(P_\rho=P_0+V_\rho+L_\rho\)，得到真正的分布恒等式
	\begin{equation}\label{g3:eq:elliptic-neumann-identity}
		(I+T_\rho)v
		=Q_0\bigl(\eta f_\rho+[P_\rho,\eta]u_\rho\bigr)
		-Q_0L_\rho v+S_0v.
	\end{equation}
	第一项的右端部分 \(Q_0(\eta f_\rho)\) 属于 \(H^{s+m}\)。交换子至多 \(m-1\) 阶，且其系数支于 \(\vartheta=1\) 的区域，所以
	\[
		[P_\rho,\eta]u_\rho=[P_\rho,\eta](\vartheta u_\rho)
		\in H^{a-m+1}.
	\]
	故它经过 \(Q_0\) 后属于 \(H^{a+1}\)。低阶项 \(Q_0L_\rho v\) 也属于 \(H^{a+1}\)，而 \(S_0v\) 属于所有 Sobolev 阶数。右侧因而属于 \(H^b\)，且其范数由
	\(C(\|\vartheta f_\rho\|_{H^s}+\|\vartheta u_\rho\|_{H^a})\) 控制。\Cref{g3:lem:two-scale-inverse}识别两空间的解，得到 \(v\in H^b\) 以及同样形式的估计。

	现在迭代这一增益。对给定 \(s,t\)，令 \(a_0=t\)、\(a_{j+1}=\min\{a_j+1,s+m\}\)，只要 \(a_j<s+m\) 就继续。至多需要 \(\lceil s+m-t\rceil\) 步。预先把球半径选得足够小，使 \(T_\rho\) 在这有限多个阶数上同时小，再选同样多层截断：每个外层截断在下一内层的支集附近等于一。最外层的 \(u\) 用已知 \(H^t\) 范数控制，每应用一次式\eqref{g3:eq:elliptic-neumann-identity}，就在下一内层增加一阶，最后达到 \(H^{s+m}\)。若 \(t\ge s+m\)，所需估计直接来自光滑乘法与包含 \(H^t\subset H^{s+m}\)。

	为取得一般截断 \(\chi,\psi\) 的估计，在 \(\psi=1\) 的开区域内用有限个这样的球覆盖 \(\operatorname{supp}\chi\)，每个球及全部嵌套截断都留在这一区域。将有限次增益估计代入，再以从属单位分解相加，得到式\eqref{g3:eq:elliptic-distribution-estimate}。所有缩放因子在选定球以后都是固定常数。

	最后处理未给出初始 Sobolev 阶数的任意分布。对每个固定紧支截断，\cref{g3:lem:fourier-sobolev-tools}保证截断后的分布属于某个 \(H^{-q}\)：有限阶分布的 Fourier 变换至多多项式增长，乘以足够高的负次权重就平方可积。取这个负阶作为起点，应用刚证明的有限步论证，即得每个紧子区域上的 \(H^{s+m}\) 正则性。
\end{proof}

\begin{corollary}\label{g3:cor:elliptic-smoothness}
	设 \(E,F\to M\) 是光滑无边界流形上的同秩有限秩实或复向量丛，\(P:C^\infty(M;E)\to C^\infty(M;F)\) 为光滑系数、正整数阶的椭圆算子。若分布截面 \(u\) 满足 \(Pu\in C^\infty(M;F)\)，则 \(u\in C^\infty(M;E)\)。在任意开集 \(U\subset M\) 上，\(u\) 光滑当且仅当 \(Pu\) 光滑于 \(U\)。
\end{corollary}
\begin{proof}
	在每张相对紧平凡化中，上一结果使 \(u\) 属于所有整数阶局部 \(H^k\)。对任意整数 \(r\ge0\)，选 \(k>r+n/2\)，由\cref{g3:cor:higher-sobolev-embedding}的局部版本，得到 \(C^r\) 代表。各阶代表在重叠处作为分布相等，连续性使它们逐点相等，所以这些代表相容为光滑截面。反向由光滑微分算子保持光滑性直接得到。实向量丛可先复化并使用同一局部估计。
\end{proof}

例如光滑度量的 Laplace 型算子作用于 \(L^2_{\mathrm{loc}}\) 截面时，若右端局部属于 \(H^k\)，则截面局部属于 \(H^{k+2}\)。特别地，分布调和截面也是光滑的。形式伴随的主符号是原主符号的纤维伴随，仍可逆，因此伴随方程具有同样的内部正则性。可解性仍由所选全局空间、边界条件及算子的核决定；本节的结论负责从一个已满足方程的分布中恢复局部导数。

% 纲要条目排为编号小节标题，并进入目录与书签；说明段仍保留于纲要注释中。

% !TeX root = ../../Book3.tex
% 纲要标题：### 12.4 Lᵖ 椭圆估计
% 纲要位置：几何学计划纲要.md，第 2576 行。
% 正文状态：初稿；尚未正式审核。

% 纲要内容开始
% 
% * 基本解和奇异积分核
% * Calderón–Zygmund 分解与奇异积分估计
% * 常系数估计与 Fourier 乘子
% * 变系数冻结、截断和吸收
% * 分布解到 W^{m,p} 的正则性
% * 椭圆系统与系数正则性
% 
% 纲要内容结束

\section{\texorpdfstring{\(L^p\)}{Lᵖ} 椭圆估计}
\label{g3:sec:1204}

在 \(L^2\) 中，频率空间的逐点界可以直接积分；换成 \(L^p\) 后，Fourier 变换不再保持范数。关键是把逆主符号的频率估计转成核在空间中的消去与衰减。下面先建立所需的奇异积分工具，再把它用于任意有限秩、等秩的椭圆系统。最后通过构造一个具有指定正则性的特解，处理原本只给定为分布的解。

\subsection{基本解和奇异积分核}
\label{g3:subsec:1204:01}
以非负算子 \(-\sum_i\partial_i^2\) 为例，其逆的核在原点奇异。记 \(\sigma_{n-1}\) 为单位球面面积。当 \(n\ge2\) 时，令
\[
	\Gamma(x)=
	\begin{cases}
	\displaystyle\frac{|x|^{2-n}}{(n-2)\sigma_{n-1}},&n\ge3,\\[4pt]
	\displaystyle-\frac1{2\pi}\log|x|,&n=2.
	\end{cases}
\]
这两个函数局部可积，在原点以外调和，且 \(\partial_r\Gamma=-1/(\sigma_{n-1}r^{n-1})\)。在穿孔球中分部积分，对测试函数令内球半径趋零，内边界的通量给出 \(\varphi(0)\)，含 \(\Gamma\partial_\nu\varphi\) 的项趋零。因此 \(-\sum_i\partial_i^2\Gamma=\delta_0\)。这样的分布称为\term[jiben jie]{基本解}[fundamental solution]。在一维，对应公式为 \(\Gamma(x)=-|x|/2\)，可直接用斜率的跳跃验证。

\ProseParagraphBreak

对紧支光滑 \(f\)，卷积 \(u=\Gamma*f\) 满足 \(-\sum_i\partial_i^2u=f\)。一阶导数核尚局部可积，二阶导数则具有临界奇性。对 \(n\ge2\)，在 \(x\ne0\) 处有
\[
	K_{ij}(x)=\partial_i\partial_j\Gamma(x)
	=\frac{n x_ix_j-\delta_{ij}|x|^2}{\sigma_{n-1}|x|^{n+2}},
	\qquad \int_{S^{n-1}}K_{ij}\,\mathrm dS=0.
\]
其中球面积分为零来自 \(\int_{S^{n-1}}x_ix_j\,\mathrm dS=\delta_{ij}\sigma_{n-1}/n\)。径向积分给出 \(\int_{a<|x|<b}K_{ij}=0\)。于是
\[
	\operatorname{pv}\int K_{ij}(x-y)f(y)\,\mathrm dy
	:=\lim_{\varepsilon\downarrow0}\int_{|x-y|>\varepsilon}K_{ij}(x-y)f(y)\,\mathrm dy
\]
对紧支光滑 \(f\) 存在：在原点附近减去 \(f(x)\)，剩余差为 \(O(|x-y|)\)，抵消了对数发散。这个极限称为\term[zhuzhi]{主值}[principal value]。

\begin{proposition}\label{g3:prop:newtonian-second-kernel}
	设 \(n\ge2\)，\(f\in C_c^\infty(\mathbb R^n)\)，令 \(\Gamma(x)=|x|^{2-n}/((n-2)\sigma_{n-1})\)（\(n\ge3\)）或 \(\Gamma(x)=-(2\pi)^{-1}\log|x|\)（\(n=2\)），其中 \(\sigma_{n-1}\) 为单位球面面积。令 \(K_{ij}(x)=(nx_ix_j-\delta_{ij}|x|^2)/(\sigma_{n-1}|x|^{n+2})\)（\(x\ne0\)），\(u=\Gamma*f\)。则在分布意义下，亦在逐点主值意义下，
	\[
		\partial_i\partial_j u
		=\operatorname{pv}(K_{ij}*f)-\frac{\delta_{ij}}n f.
	\]
\end{proposition}
\begin{proof}
	在 \(|x|>\varepsilon\) 上两次分部积分。由于 \(\partial_j\Gamma=-x_j/(\sigma_{n-1}|x|^n)\)，内边界中含测试函数本身的一项趋于
	\(-\sigma_{n-1}^{-1}\int_{S^{n-1}}\theta_i\theta_j\varphi(0)\,\mathrm dS=-\delta_{ij}\varphi(0)/n\)。另一边界项至多为 \(C\varepsilon\)（二维为 \(C\varepsilon|\log\varepsilon|\)），趋零。体积项用上述消去性质减去测试函数在零点的值后收敛，故
	\(\partial_i\partial_j\Gamma=\operatorname{pv}K_{ij}-(\delta_{ij}/n)\delta_0\)。再与 \(f\) 卷积即得结论。求迹时，所有 \(K_{ii}\) 之和为零，而点质量项之和为 \(-\delta_0\)，与基本解的符号相符。
\end{proof}

核的绝对值大小为 \(|x|^{-n}\)，在原点附近不能积；可用的性质是零均值所带来的消去。下一小节把这一现象写成一个能够复用的算子定理。

\subsection{Calderón–Zygmund 分解与奇异积分估计}
\label{g3:subsec:1204:02}
设 \(f\in L^1(\mathbb R^n;\mathbb C^r)\)。选择平均值较大的区域，把 \(f\) 分成有界部分与局部零均值部分，这就是 Calderón--Zygmund 分解的用途。\footnote{Alberto Pedro Calderón（1920--1998），阿根廷数学家；Antoni Szczepan Zygmund（1900--1992），波兰出生、后在美国工作的数学家。两人在奇异积分、Fourier 分析及偏微分方程方面作出基础性贡献。}

\begin{lemma}\label{g3:lem:cz-decomposition}
	设 \(n,r\ge1\)、\(f\in L^1(\mathbb R^n;\mathbb C^r)\)，\(\lambda>0\)。存在两两不交的半开二进立方体 \(Q_j\)，使
	\[
		\lambda<\frac1{|Q_j|}\int_{Q_j}|f|\le2^n\lambda,
		\qquad\sum_j|Q_j|\le\lambda^{-1}\|f\|_1.
	\]
	写 \(f=g+\sum_jb_j\)，其中 \(g=f\) 于这些立方体之外，\(g=f_{Q_j}\) 于 \(Q_j\)，\(b_j=(f-f_{Q_j})\mathbf1_{Q_j}\)。则
	\[
		\|g\|_\infty\le2^n\lambda,\quad\|g\|_1\le\|f\|_1,
		\int b_j=0,\quad\sum_j\|b_j\|_1\le2\|f\|_1.
	\]
\end{lemma}
\begin{proof}
	二进立方体指边长 \(2^k\)、顶点坐标为 \(2^k\) 整数倍的半开立方体。任意两个这样的立方体或者不交，或者一个包含另一个。由于大立方体的平均值至多为 \(2^{-kn}\|f\|_1\)，平均值超过 \(\lambda\) 的立方体都含于某个极大的此类立方体。取全部极大者作为 \(Q_j\)；其父立方体不满足严格不等式，故平均值至多为 \(2^n\lambda\)。对互不相交的 \(Q_j\) 积分就得体积和的界。

	还须证明其外部几乎处处有 \(|f|\le\lambda\)。说明所用的二进微分事实：对 \(h\in L^1\)，二进最大平均 \(M_dh=\sup_{Q\ni x}\frac1{|Q|}\int_Q|h|\) 满足
	\(|\{M_dh>t\}|\le\|h\|_1/t\)，因为取上述极大立方体即可。取连续紧支 \(h_0\) 逼近 \(h\) 于 \(L^1\)。沿缩向 \(x\) 的二进立方体，\(h_0\) 的平均趋于 \(h_0(x)\)，而 \(h\) 的平均与 \(h(x)\) 的极限上差至多为 \(M_d(h-h_0)(x)+|h-h_0|(x)\)。最大平均估计与积分的 Chebyshev 估计使这个上差超过任意 \(t>0\) 的集合测度至多为 \(2\|h-h_0\|_1/t\)。令逼近误差趋零，即得平均几乎处处趋于函数值。将它用于 \(|f|\)，在全部 \(Q_j\) 外的每个平均都不超过 \(\lambda\)，所以 \(|f|\le\lambda\) 几乎处处成立。

	其余断言由 \(|f_{Q_j}|\le\frac1{|Q_j|}\int_{Q_j}|f|\) 及三角不等式逐块得到。若还有 \(f\in L^2\)，则 \(g,b_j\) 也属于 \(L^2\)，且坏部分之和在 \(L^2\) 中收敛。
\end{proof}

\begin{theorem}\label{g3:thm:cz-lp-module}
	设 \(T\) 是 \(L^2(\mathbb R^n;\mathbb C^r)\) 上有界的线性算子。在原点以外有矩阵值 \(C^1\) 核 \(K\)，满足
	\[
		|K(z)|\le A|z|^{-n},\qquad|\mathrm DK(z)|\le A|z|^{-n-1}.
	\]
	假设对紧支 \(L^2\) 函数，在其支集之外有 \(Tf(x)=\int K(x-y)f(y)\,\mathrm dy\)。则
	\[
		|\{|Tf|>\lambda\}|\le C\lambda^{-1}\|f\|_1
		\quad(f\in L^1\cap L^2),
	\]
	且 \(T\) 唯一延拓为每个 \(L^p\)、\(1<p<\infty\) 上的有界算子。常数依赖 \(n,r,p,A\) 及 \(\|T\|_{2\to2}\)。
\end{theorem}
\begin{proof}
	对 \(f\) 作水平 \(\lambda\) 的分解。好部分满足 \(\|g\|_2^2\le2^n\lambda\|f\|_1\)，所以 \(|\{|Tg|>\lambda/2\}|\le C\|f\|_1/\lambda\)。将每个 \(Q_j\) 同心放大 \(4\sqrt n\) 倍，所有放大立方体之并的测度也至多为 \(C\|f\|_1/\lambda\)。记 \(c_j\) 为其中心、\(\ell_j\) 为边长。在放大立方体外，零均值使
	\[
		Tb_j(x)=\int_{Q_j}[K(x-y)-K(x-c_j)]b_j(y)\,\mathrm dy,
		\quad |Tb_j(x)|\le\frac{C\ell_j}{|x-c_j|^{n+1}}\|b_j\|_1.
	\]
	在外部积分得 \(\int|Tb_j|\le C\|b_j\|_1\)。对有限和求和后，用 \(L^2\) 收敛所给的几乎处处子列与积分下半连续性取极限，得到坏部分在放大立方体之外的 \(L^1\) 界 \(C\|f\|_1\)。再用 Chebyshev 估计，便证明弱型不等式。

	为得到强 \(L^p\) 界，先取 \(1<p<2\) 以及有界紧支 \(f\)。在每个 \(t>0\) 处拆成 \(f_{>t}=f\mathbf1_{\{|f|>t\}}\) 和 \(f_{\le t}\)。弱型界用于前者，\(L^2\) 界用于后者，给出
	\[
		|\{|Tf|>2t\}|\le\frac C t\int_{|f|>t}|f|
		+\frac C{t^2}\int_{|f|\le t}|f|^2.
	\]
	乘 \(p t^{p-1}\) 后积分。由 \(\|h\|_p^p=p\int_0^\infty t^{p-1}|\{|h|>t\}|\,\mathrm dt\) 与非负积分换序，右侧分别等于常数乘
	\[
		\int |f(x)|\int_0^{|f(x)|}t^{p-2}\,\mathrm dt\,\mathrm dx
		=\frac{\|f\|_p^p}{p-1},\qquad
		\int |f(x)|^2\int_{|f(x)|}^\infty t^{p-3}\,\mathrm dt\,\mathrm dx
		=\frac{\|f\|_p^p}{2-p}.
	\]
	故 \(\|Tf\|_p\le C_p\|f\|_p\)。稠密性给出延拓。

	当 \(p>2\) 时，\(L^2\) 伴随 \(T^*\) 的离对角核为 \(K(-z)^*\)：先对支集不交的紧支函数用 Fubini 公式核对，再用 \(L^2\) 逼近。这一核满足相同估计，所以已证结论适用于 \(p'=p/(p-1)<2\)。由 \(L^p\) 积分对偶，
	\( |\langle Tf,h\rangle|=|\langle f,T^*h\rangle|\le C\|f\|_p\|h\|_{p'}\)，得到所需界。有限维向量的 Euclid 范数和矩阵算子范数满足上述全部估计，故证明同时适用于有限秩系统。
\end{proof}

\subsection{常系数估计与 Fourier 乘子}
\label{g3:subsec:1204:03}
对矩阵函数 \(a(\xi)\)，记 \(a(\mathrm D)f=\mathcal F^{-1}(a\widehat f)\)。此处 \(\mathrm D_j=-\mathrm i\partial_j\)，与前节主符号约定相同。光滑齐次零阶符号在每个环带上的变化受尺度控制，因而可以应用刚建立的奇异积分定理。

\begin{theorem}\label{g3:thm:smooth-multiplier-lp}
	设 \(a:\mathbb R^n\setminus\{0\}\to\operatorname{End}(\mathbb C^r)\) 为 \(C^{n+2}\) 函数，且
	\[
		|\partial_\xi^\beta a(\xi)|\le A_\beta|\xi|^{-|\beta|},\qquad |\beta|\le n+2.
	\]
	则 \(a(\mathrm D)\) 在每个 \(L^p(\mathbb R^n;\mathbb C^r)\)、\(1<p<\infty\) 上有界，其范数受这些常数、\(n,r,p\) 控制。特别地，光滑零阶齐次矩阵函数满足此结论。
\end{theorem}
\begin{proof}
	取径向光滑 \(\rho\)，在单位球为一、在半径二球外为零，令 \(\psi(\xi)=\rho(\xi)-\rho(2\xi)\)。则 \(\sum_{j\in\mathbb Z}\psi(2^{-j}\xi)=1\) 当 \(\xi\ne0\)。写 \(a_j=a\psi(2^{-j}\cdot)\)，并令 \(k_j=(2\pi)^{-n/2}\mathcal F^{-1}a_j\)，所以 \(a_j(\mathrm D)f=k_j*f\)。缩放到固定环带后，\(a(2^j\xi)\psi(\xi)\) 的至 \(n+2\) 阶导数一致有界。

	在逆变换积分中，零次估计给出 \(|k_j(x)|\le C2^{jn}\)。当 \(x\ne0\) 时，沿 \(x/|x|\) 的方向对频率变量分部积分 \(n+2\) 次；支集紧含于环带，没有边界项。因此
	\[
		|k_j(x)|\le C2^{jn}(1+2^j|x|)^{-n-2},\qquad
		|\mathrm Dk_j(x)|\le C2^{j(n+1)}(1+2^j|x|)^{-n-2}.
	\]
	第二式只在频率积分中多一个 \(\xi\) 因子；在固定环带上它及所需导数都有界。按 \(2^j|x|\le1\) 与 \(>1\) 将几何级数分开，得到
	\[
		\sum_j|k_j(x)|\le C|x|^{-n},\qquad
		\sum_j|\mathrm Dk_j(x)|\le C|x|^{-n-1}.
	\]
	所以 \(K=\sum_jk_j\) 在原点之外局部 \(C^1\) 收敛，并满足核估计。

	Plancherel 恒等式给出 \(\|a(\mathrm D)\|_{2\to2}\le\|a\|_\infty\)，且环带部分和在 \(L^2\) 中趋于这一算子。对支集分离的紧支光滑 \(f,h\)，上面的核界在其支集乘积上可积，可以交换求和与双重积分，故极限算子的离对角核确为 \(K(x-y)\)。再以 \(L^2\) 逼近得到定理所需的紧支 \(L^2\) 版本。应用\cref{g3:thm:cz-lp-module}便完成证明。
\end{proof}

\begin{corollary}\label{g3:cor:constant-elliptic-lp}
	设 \(P_0=\sum_{|\alpha|=m}A_\alpha\mathrm D^\alpha\) 为 \(\mathbb R^n\) 上的常系数齐次算子，\(m\ge1\)，各 \(A_\alpha\) 为 \(r\times r\) 复矩阵，且 \(p_m(\xi)=\sum A_\alpha\xi^\alpha\) 对所有 \(\xi\ne0\) 可逆。对 \(1<p<\infty\)、\(u\in C_c^\infty(\mathbb R^n;\mathbb C^r)\)，有
	\[
		\sum_{|\alpha|=m}\|\partial^\alpha u\|_p\le C\|P_0u\|_p.
	\]
	若 \(\rho\in C_c^\infty\) 在零点附近为一，则符号 \(q(\xi)=(1-\rho(\xi))p_m(\xi)^{-1}\) 所定的 \(Q_0=q(\mathrm D)\) 有界地将 \(L^p\) 映到 \(W^{m,p}\)，并满足
	\(P_0Q_0=Q_0P_0=I-S_0\)，其中 \(S_0=\rho(\mathrm D)\)。
\end{corollary}
\begin{proof}
	对 \(|\alpha|=m\)，矩阵函数 \(\xi^\alpha p_m(\xi)^{-1}\) 光滑且零阶齐次。将乘子定理用于 \(P_0u\)，得到各最高阶导数估计；\(\partial^\alpha=\mathrm i^{|\alpha|}\mathrm D^\alpha\) 的相位不改变范数。对高频逆，\(\xi^\alpha q(\xi)\)、\(|\alpha|\le m\) 在低频为零，在高频的阶数不超过零，满足同一乘子条件，故控制全部 \(W^{m,p}\) 导数。算子恒等式在 Fourier 侧逐点成立，再延伸到分布。
\end{proof}

在基本解的例子中，二阶导数对应 \(-\xi_i\xi_j/|\xi|^2\)。其 \(L^p\) 有界性于是给出二阶 Poisson 估计。矩阵逆的使用也说明：这里的系统估计取决于整个主符号可逆，而非各个分量方程分别满足某个标量不等式。

\subsection{变系数冻结、截断和吸收}
\label{g3:subsec:1204:04}
最高阶系数在足够小的球上接近常数，因而可由常系数估计吸收；截断所产生的导数都降低了一阶。为了把这些低阶项最后压到函数本身的范数，先记录一个插值估计。

\begin{lemma}\label{g3:lem:lp-derivative-interpolation}
	设 \(m\ge1\)、\(1<p<\infty\)，则对 \(0\le j<m\)、\(0<\delta\le1\) 及 \(v\in W^{m,p}(\mathbb R^n)\)，
	\[
		\|\mathrm D^jv\|_p\le C\delta^{m-j}\|\mathrm D^mv\|_p+C\delta^{-j}\|v\|_p.
	\]
	这里同阶导数的范数表示其有限个分量范数之和。
\end{lemma}
\begin{proof}
	先取 \(\delta=1\)。低频 \(\rho(\mathrm D)\mathrm D^jv\) 是可积光滑核与 \(v\) 的卷积，受 \(\|v\|_p\) 控制。高频可写成最高阶导数的有限个零阶乘子之和。具体地，令 \(b(\xi)=\sum_{|\alpha|=m}\xi^{2\alpha}>0\) 当 \(\xi\ne0\)，则对 \(|\beta|=j\)，
	\[
		(1-\rho)\xi^\beta
		=\sum_{|\alpha|=m}\frac{(1-\rho)\xi^{\beta+\alpha}}{b(\xi)}\xi^\alpha.
	\]
	系数在高频的阶数 \(j-m<0\)，满足乘子定理。故 \(\|\mathrm D^jv\|_p\le C(\|\mathrm D^mv\|_p+\|v\|_p)\)。对 \(v(\delta x)\) 应用它，消去共同体积因子，就得带 \(\delta\) 的公式。光滑稠密性保证它适用于所述空间。
\end{proof}

\begin{theorem}\label{g3:thm:local-elliptic-lp-apriori}
	设 \(\Omega\subset\mathbb R^n\) 开，\(P=\sum_{|\alpha|\le m}A_\alpha(x)\mathrm D^\alpha\) 为光滑系数、有限秩等秩椭圆系统，\(m\ge1\)。若 \(1<p<\infty\)、\(U\Subset V\Subset\Omega\) 为开集，则
	\[
		\|u\|_{W^{m,p}(U)}\le C\bigl(\|Pu\|_{L^p(V)}+\|u\|_{L^p(V)}\bigr)
	\]
	对所有 \(u\in W^{m,p}(V;\mathbb C^r)\) 成立。常数依赖这些固定区域、\(p,m\)、系数界与主符号逆的球面界。此先验估计实际只需最高阶系数连续、其余系数局部有界。
\end{theorem}
\begin{proof}
	在 \(x_0\) 冻结最高阶系数。由椭圆性和球面的紧性，常系数估计有有限常数 \(C_0\)。把球缩小到使最高阶系数的振幅小于 \((2C_0)^{-1}\) 除以最高阶项数。对支于此球的 \(v\)，有
	\[
		\|\mathrm D^mv\|_p\le C_0\|Pv\|_p+\tfrac12\|\mathrm D^mv\|_p+C\|v\|_{W^{m-1,p}}.
	\]
	移项后，用上一引理选择足够小的 \(\delta\) 吸收低阶导数，得到
	\(\|v\|_{W^{m,p}}\le C(\|Pv\|_p+\|v\|_p)\)。

	说明局部化后如何去掉较大球上的导数项。固定一个上述小球 \(B_R\)，取 \(R/2\le r<t<R\) 和截断 \(\eta\)，使 \(\eta=1\) 于 \(B_r\)、支于 \(B_t\)，并使 \(|\mathrm D^k\eta|\le C_k(t-r)^{-k}\)。乘积公式使 \([P,\eta]\) 的阶数至多为 \(m-1\)。记 \(F(r)=\|u\|_{W^{m,p}(B_r)}\)，则
	\[
		F(r)\le C\|Pu\|_{L^p(B_R)}+C\sum_{j<m}(t-r)^{j-m}\|\mathrm D^ju\|_{L^p(B_t)}.
	\]
	球上的延拓算子来自\cref{g3:prop:higher-order-boundary-extension}；其有限反射与截断构造还满足零阶界 \(\|Eu\|_p\le C\|u\|_p\)。经缩放，延拓的各阶范数在 \(R/2\le t\le R\) 一致受控；将上一引理用于延拓，便对任意固定 \(\varepsilon>0\) 得
	\[
		F(r)\le\varepsilon F(t)+C_\varepsilon(t-r)^{-m}
		\bigl(\|Pu\|_{L^p(B_R)}+\|u\|_{L^p(B_R)}\bigr),
	\]
	其中可将固定的 \(R\) 因子并入常数。取 \(r_j=R(1-2^{-j-1})\)，选 \(\varepsilon<2^{-m-1}\) 并逐次代入。几何级数 \(\sum_j\varepsilon^j2^{mj}\) 收敛，余项 \(\varepsilon^NF(r_N)\le\varepsilon^N\|u\|_{W^{m,p}(B_R)}\) 趋零，故得 \(B_{R/2}\) 上所需估计。最后以有限个这种小球覆盖 \(\overline U\)，并使其大球均含于 \(V\)，相加即可。
\end{proof}

每次冻结均在选定紧集内进行。若一族算子具有相同的主符号逆界、最高阶系数连续模和低阶系数界，就能统一选择小球尺度，并得到统一的局部估计常数。

\subsection{\texorpdfstring{分布解到 \(W^{m,p}\) 的正则性}{分布解到 W\textasciicircum{}\{m,p\} 的正则性}}
\label{g3:subsec:1204:05}
上一不等式以 \(u\in W^{m,p}\) 为前提。对于尚无任何函数代表的分布解，直接代入没有意义。下面先为同一右端构造一个 \(W^{m,p}\) 特解，再利用上一节的分布椭圆光滑性识别剩余部分。

\begin{theorem}\label{g3:thm:elliptic-distribution-lp}
	设 \(\Omega\subset\mathbb R^n\) 开，\(P\) 是阶 \(m\ge1\)、光滑系数的有限秩等秩椭圆系统，\(1<p<\infty\)。若 \(u\in\mathcal D'(\Omega;\mathbb C^r)\) 且 \(Pu\in L^p_{\mathrm{loc}}\)，则 \(u\in W^{m,p}_{\mathrm{loc}}\)。更一般地，若整数 \(k\ge0\) 且 \(Pu\in W^{k,p}_{\mathrm{loc}}\)，则 \(u\in W^{k+m,p}_{\mathrm{loc}}\)。若还已知 \(u\in L^p(V)\)，则前一定理的局部估计直接适用。
\end{theorem}
\begin{proof}
	固定 \(x_0\)，令 \(x=x_0+\rho y\)，并把算子乘以 \(\rho^m\)。其最高阶系数为 \(A_\alpha(x_0+\rho y)\)，低阶系数为 \(\rho^{m-|\alpha|}A_\alpha(x_0+\rho y)\)。在固定单位球中，前者减去 \(A_\alpha(x_0)\) 后为 \(O(\rho)\)，所有低阶系数也为 \(O(\rho)\)。用固定截断将这些差与低阶系数延拓到全空间，得到
	\[
		\widetilde P=P_0+B,\qquad
		\|Bv\|_p\le C\rho\|v\|_{W^{m,p}(\mathbb R^n)}.
	\]
	这里 \(P_0\) 为冻结的齐次主部，\(B\) 的系数光滑紧支；\(\widetilde P\) 在较小单位球中等于缩放后的原算子。取上一小节的高频逆 \(Q_0:L^p\to W^{m,p}\)。选 \(\rho\) 使 \(\|Q_0B\|_{W^{m,p}\to W^{m,p}}<1/2\)。

	在这一小球中将缩放后的右端截断为紧支 \(f_0\in L^p(\mathbb R^n)\)，使它在更小球上与原右端相同。级数
	\[
		v=(I+Q_0B)^{-1}Q_0f_0
		=\sum_{j=0}^\infty(-Q_0B)^jQ_0f_0
	\]
	在完备空间 \(W^{m,p}\) 中绝对收敛。它满足 \(v+Q_0Bv=Q_0f_0\)。施以 \(P_0\) 并用 \(P_0Q_0=I-S_0\)，得到精确恒等式
	\[
		\widetilde Pv=f_0+S_0(Bv-f_0).
	\]
	\(S_0\) 的符号光滑紧支，因此其卷积核是快速下降光滑函数。由于 \(Bv-f_0\) 紧支且属于 \(L^p\)，卷积的每一阶导数都可移到这个核上，故附加项光滑。

	在更小球中，原分布 \(u\)（换到 \(y\) 坐标）与 \(v\) 的差满足同一个光滑椭圆算子作用后为光滑函数。由\cref{g3:thm:elliptic-distribution-l2}的分布光滑性，该差本身光滑。因此 \(u\) 在这个球中属于 \(W^{m,p}\)。\(x_0\) 任意，得到局部结论。这一步没有对未知分布预先假定任何 \(L^p\) 范数。

	最后证明高阶版本。若已得 \(u\in W^{m+j-1,p}_{\mathrm{loc}}\) 且 \(f=Pu\in W^{j,p}_{\mathrm{loc}}\)，对 \(|\beta|=j\) 有
	\[
		P(\partial^\beta u)=\partial^\beta f
		-\sum_{\alpha}\sum_{0<\gamma\le\beta}
		\binom\beta\gamma(\partial^\gamma A_\alpha)\mathrm D^\alpha\partial^{\beta-\gamma}u.
	\]
	右边涉及的 \(u\) 的阶数至多为 \(m+j-1\)，故属于局部 \(L^p\)。已经证明的分布结论适用于 \(\partial^\beta u\)，得到它属于 \(W^{m,p}_{\mathrm{loc}}\)。对 \(j\) 归纳，即得 \(W^{m+k,p}_{\mathrm{loc}}\)。
\end{proof}

\subsection{椭圆系统与系数正则性}
\label{g3:subsec:1204:06}
本节的主符号可以是非对角矩阵。矩阵求逆、有限维核范数和向量值分解都保留了分量间耦合，因此所得估计适用于光滑等秩椭圆系统，包括 Laplace 型系统和下面的一阶 Cauchy--Riemann 型系统。对实系数算子先复化，所得估计再限制回实截面即可。丛上的结论则在平凡化中使用：过渡函数和局部截断的有限阶导数在紧集上有界，故不同表示给出等价的局部范数。

\ProseParagraphBreak

各结论所用的系数条件可以分开。\cref{g3:thm:local-elliptic-lp-apriori}对已经属于 \(W^{m,p}\) 的函数，只用最高阶系数连续以及低阶系数有界；最高阶系数的连续模决定冻结半径。\cref{g3:thm:elliptic-distribution-lp}则对任意分布给出正则性，其特解构造之后调用了光滑系数的分布椭圆正则性，因此这里保留全部系数光滑。高阶自举中，求导系数的阶数也随所求正则性增加。有限正则性系数的更精细版本须逐阶核对，不能仅删除陈述中的“光滑”。

\ProseParagraphBreak

指数 \(1<p<\infty\) 同样进入证明。弱型与 \(L^2\) 插值在 \(p=1\) 处出现发散积分，对偶论证也不能给出 \(L^\infty\) 有界性。核的消去保证中间指数的估计，却一般不保证端点估计。下一节采用 Hölder 模连续性，得到另一种适合经典解的控制；它与端点 \(L^\infty\) 估计有不同的假设。
% 纲要条目排为编号小节标题，并进入目录与书签；说明段仍保留于纲要注释中。

% !TeX root = ../../Book3.tex
% 纲要标题：### 12.5 Schauder 方法
% 纲要位置：几何学计划纲要.md，第 2585 行。
% 正文状态：初稿；尚未正式审核。

% 纲要内容开始
% 
% * Hölder 半范数及尺度
% * 常系数的振荡控制
% * 局部插值不等式
% * 变系数估计与高阶自举
% 
% 纲要内容结束

\section{Schauder 方法}
\label{g3:sec:1205}

积分估计控制导数的平均大小，Hölder 估计则追踪两个邻近点的导数怎样变化。若方程右端在尺度 \(r\) 上的振荡为 \(r^\alpha\)，椭圆方程通常把同一指数传给解的最高阶导数；奇异核的平均抵消使这个传递成为可能。\term[Schauder fangfa]{Schauder 方法}[Schauder method]\footnote{Juliusz Paweł Schauder（1899--1943），波兰数学家，研究泛函分析、偏微分方程与非线性分析。}在常系数估计之后，通过缩小坐标球使最高阶系数的变化可以吸收。本节的主部取实标量正定矩阵，允许它逐分量作用于有限维向量值函数。

\subsection{Hölder 半范数及尺度}
\label{g3:subsec:1205:01}
固定 \(0<\alpha<1\)。对 \(B_r(x_0)\subset\mathbb R^n\)，采用带尺度的范数
\[
	\|u\|_{k,\alpha;r}
	=\sum_{j=0}^k r^j\|\mathrm D^ju\|_{L^\infty(B_r)}
	+r^{k+\alpha}[\mathrm D^ku]_{\alpha;B_r},
	\qquad
	[v]_{\alpha;A}=\sup_{x\ne y\in A}\frac{|v(x)-v(y)|}{|x-y|^\alpha}.
\]
\symidx{\(\lVert u\rVert_{k,\alpha;r}\)：带尺度的 Hölder 范数}
这里各阶导数取有限个分量的 Euclid 范数；复值函数按实分量处理。换元 \(x=x_0+ry\) 后，这恰与单位球上的范数一致。对二阶方程，右端 \(f\) 应乘 \(r^2\)，一阶与零阶系数则分别乘 \(r\) 与 \(r^2\)。

\begin{lemma}\label{g3:lem:holder-product-scale}
	设 \(0<\alpha<1\)，\(a,v\) 为同一集合上的有界 \(\alpha\)-Hölder 函数，取值允许为相容大小的矩阵与向量。则
	\[
		[av]_\alpha\le\|a\|_\infty[v]_\alpha+[a]_\alpha\|v\|_\infty.
	\]
	因此带尺度的零阶 Hölder 范数满足乘法估计。若 \(a\in C^{0,\alpha}(B_r(x_0))\)，则
	\[
		\|a-a(x_0)\|_\infty+r^\alpha[a]_\alpha
		\le2r^\alpha[a]_\alpha.
	\]
\end{lemma}
\begin{proof}
	把 \(a(x)v(x)-a(y)v(y)\) 写成
	\(a(x)(v(x)-v(y))+(a(x)-a(y))v(y)\)，除以 \(|x-y|^\alpha\) 后取上确界，即得第一式。第二式由 \(|x-x_0|\le r\) 直接得到。
\end{proof}

冻结系数时，要小的是 \(a-a(x_0)\) 的带尺度 Hölder 范数。仅有 \(\|a-a(x_0)\|_\infty\) 小，并不能直接删去乘积估计中的第二项。后面会把该项放到二阶导数的上确界范数中，再用插值处理。

\subsection{常系数的振荡控制}
\label{g3:subsec:1205:02}
为保持最大值原理中常用的符号，暂记
\[
	\mathcal L_0=\sum_{i=1}^n\partial_i^2=-\Delta_{\mathrm{Eucl}}.
\]
\(n\ge2\) 时，令 \(\sigma_{n-1}\) 为单位球面面积，取上一节非负 Laplace 算子基本解的相反数，记为
\[
	\Phi(x)=
	\begin{cases}
		\displaystyle\frac{|x|^{2-n}}{(2-n)\sigma_{n-1}},&n\ge3,\\[4pt]
		\displaystyle\frac1{2\pi}\log|x|,&n=2.
	\end{cases}
\]
这里 \(\mathcal L_0\Phi=\delta_0\)。确实，零点以外 \(\mathcal L_0\Phi=0\)，而 \(\partial_r\Phi=\sigma_{n-1}^{-1}r^{1-n}\) 的球面通量恒为一；在穿孔球上对测试函数分部积分，再令内半径趋零，就得到这个分布恒等式。

\begin{lemma}\label{g3:lem:newton-holder-estimate}
	设 \(n\ge2\)、\(0<\alpha<1\)，\(f\in C_c^{0,\alpha}(\mathbb R^n)\) 支于 \(B_2\)。则 \(w=\Phi*f\) 属于 \(C^{2,\alpha}_{\mathrm{loc}}(\mathbb R^n)\)，满足 \(\mathcal L_0w=f\)，且
	\[
		\|w\|_{C^{2,\alpha}(B_3)}
		\le C(n,\alpha)\|f\|_{C^{0,\alpha}(\mathbb R^n)}.
	\]
\end{lemma}
\begin{proof}
	真正需要估计的是二阶导数。对 \(z\ne0\)，写
	\[
		K_{ij}(z)=\partial_i\partial_j\Phi(z)
		=\frac1{\sigma_{n-1}}\left(\frac{\delta_{ij}}{|z|^n}
		-\frac{n z_i z_j}{|z|^{n+2}}\right).
	\]
	它满足 \(|K(z)|\le C|z|^{-n}\)、\(|\mathrm DK(z)|\le C|z|^{-n-1}\)，并且在每个以零为中心的球面、因而每个同心球壳上的积分为零。对穿孔区域再次分部积分，内球面上的 \(z_i z_j\) 平均为 \(\delta_{ij}/n\)，于是
	\begin{equation}\label{g3:eq:newton-second-derivative}
		\partial_i\partial_jw(x)
		=\frac{\delta_{ij}}n f(x)
		+\operatorname{pv}\int K_{ij}(x-y)f(y)\,\mathrm dy.
	\end{equation}
	在 \(|x-y|<1\) 中把 \(f(y)\) 换成 \(f(y)-f(x)\)，不改变主值；此时绝对值积分不超过
	\(C[f]_\alpha\int_0^1 t^{\alpha-1}\,\mathrm dt\)。外部积分在 \(x\in B_3\) 上由固定支集和 \(\|f\|_\infty\) 控制。这已给出二阶导数的上确界估计。

	说明 Hölder 估计中怎样避免一个多余的对数。令 \(h=|x-x'|>0\)，把积分分成距两点至多 \(3h\) 的部分和其余部分。在近区分别减去 \(f(x)\)、\(f(x')\)，其积分大小不超过
	\[
		C[f]_\alpha\int_0^{4h}t^{\alpha-1}\,\mathrm dt
		\le C\alpha^{-1}[f]_\alpha h^\alpha.
	\]
	在公共远区，差核满足
	\(|K(x-y)-K(x'-y)|\le Ch|x-y|^{-n-1}\)。将它乘以 \(f(y)-f(x)\) 后，积分不超过
	\[
		Ch[f]_\alpha\int_{2h}^{\infty}t^{\alpha-2}\,\mathrm dt
		\le C(1-\alpha)^{-1}[f]_\alpha h^\alpha.
	\]
	还会出现 \((f(x)-f(x'))\) 乘以远区核积分。这个核积分有与 \(h\) 无关的界：以 \(x'\) 为中心的完整球壳积分为零，而被另一近球切掉的部分位于半径在 \(2h\) 与 \(4h\) 之间的球壳，其绝对值积分至多为 \(C\log2\)。外截断半径趋于无穷时的两球壳差趋零。因而这一项同样至多为 \(C[f]_\alpha h^\alpha\)。式\eqref{g3:eq:newton-second-derivative}中的局部项也有同样估计。

	上述公式先对光滑 \(f\) 成立。对一般 \(f\)，卷积得到 \(f_\varepsilon\)，其 \(C^{0,\alpha}\) 范数一致有界，且对每个 \(0<\beta<\alpha\) 在 \(C^{0,\beta}\) 中趋于 \(f\)：近点用统一的 \(\alpha\)-半范数，远点用一致收敛即可证明。用刚证的 \(\beta\) 估计知 \(\mathrm D^2(\Phi*f_\varepsilon)\) 局部一致收敛；\(\Phi\) 及其一阶导数局部可积，所以函数与一阶导数也局部一致收敛。沿线段积分导数关系，极限确为 \(w\) 的经典二阶导数。最后对任意两点取极限保留 \(\alpha\)-半范数界。函数及一阶导数的上确界由可积核直接估计，证明完成。
\end{proof}

\begin{theorem}\label{g3:thm:constant-schauder}
	设 \(n\ge1\)、\(0<\alpha<1\)，实对称常矩阵 \(A=(a^{ij})\) 满足 \(\lambda I\le A\le\Lambda I\)，其中 \(0<\lambda\le\Lambda\)。若 \(u\in C^2(B_{2r}(x_0);\mathbb R^N)\) 且
	\(a^{ij}\partial_i\partial_j u=f\in C^{0,\alpha}(B_{2r})\)，则 \(u\in C^{2,\alpha}(B_r)\)，并有
	\begin{equation}\label{g3:eq:constant-schauder}
		\|u\|_{2,\alpha;r}
		\le C\bigl(\|u\|_{L^\infty(B_{2r})}
		+r^2\|f\|_{L^\infty(B_{2r})}
		+r^{2+\alpha}[f]_{\alpha;B_{2r}}\bigr),
	\end{equation}
	其中 \(C\) 只依赖 \(n,N,\alpha,\lambda,\Lambda\)。
\end{theorem}
\begin{proof}
	先令 \(A=I\)、\(r=1\)。取支于 \(B_2\)、在 \(B_{3/2}\) 上为一的光滑截断 \(\eta\)，并令 \(w=\Phi*(\eta f)\)。上一引理给出 \(w\) 的所需估计，\(h=u-w\) 在 \(B_{3/2}\) 上调和。

	调和函数的高阶导数可直接由平均值公式控制。球面平均对半径的导数等于球内 \(\mathcal L_0h\) 的积分除以球面面积，故平均值等于中心值。将球面平均对半径作光滑径向平均，得到 \(h=\rho_\delta*h\)，其中 \(\rho_\delta\) 为支于 \(B_\delta\) 的固定光滑径向核，积分为一。对 \(x\in B_{5/4}\) 取固定小 \(\delta\)，把任意阶导数放到核上，得到
	\[
		\|\mathrm D^jh\|_{L^\infty(B_{5/4})}
		\le C_j\|h\|_{L^\infty(B_{3/2})}.
	\]
	\(j=2,3\) 的估计分别处理远点与近点，给出 \([\mathrm D^2h]_{\alpha;B_1}\le C\|h\|_\infty\)。与 \(w\) 相加，得到结论。

	一般常矩阵用 \(x=A^{1/2}y\) 化为上述情形。由特征值界，变换后的椭球包含与外切的球半径之比受 \(\lambda,\Lambda\) 控制；用有限个固定比例的小球覆盖较小椭球，近点用同一球内的估计，远点用二阶导数上确界，便得到相同形式的 Hölder 界。缩放恢复 \(r\)，向量值情况逐分量处理。\(n=1\) 时，\(u''=f/a^{11}\) 直接给出二阶半范数，沿区间积分并用函数上确界控制一次积分常数，得到同一估计。
\end{proof}

核的抵消把不可积的 \(|x-y|^{-n}\) 与增量 \(f(y)-f(x)\) 配对。条件 \(\alpha>0\) 保证近区积分收敛，\(\alpha<1\) 保证差核的远区积分具有 \(h^\alpha\) 大小；两个端点需要不同的连续性尺度。

\subsection{局部插值不等式}
\label{g3:subsec:1205:03}
截断会引入 \(\mathrm Du\) 与 \(u\)，冻结系数的乘积还会产生 \(\|\mathrm D^2u\|_\infty\)。这些量虽低于最高的 Hölder 半范数，却必须带上可任意缩小的系数。下面用有限差分给出这种插值，而不预先假定更高阶导数。

\begin{lemma}\label{g3:lem:weighted-holder-interpolation}
	设 \(\Omega\subset\mathbb R^n\) 为有界开集，\(0<\alpha<1\)，\(u\in C^{2,\alpha}(\overline\Omega)\)。记 \(d(x)=\operatorname{dist}(x,\partial\Omega)\)、\(d_{xy}=\min\{d(x),d(y)\}\)，并令
	\[
		M_j=\sup_{x\in\Omega}d(x)^j|\mathrm D^ju(x)|\quad(0\le j\le2),
		\qquad
		H_2=\sup_{x\ne y}d_{xy}^{2+\alpha}
		\frac{|\mathrm D^2u(x)-\mathrm D^2u(y)|}{|x-y|^\alpha}.
	\]
	则对每个 \(\varepsilon>0\)，存在 \(C_\varepsilon=C(n,\alpha,\varepsilon)\)，使
	\begin{equation}\label{g3:eq:weighted-holder-interpolation}
		M_1+M_2\le\varepsilon H_2+C_\varepsilon M_0.
	\end{equation}
	此外，按同样权重定义的 \(H_1,H_0\) 满足
	\(H_1\le C(M_1+M_2)\)、\(H_0\le C(M_0+M_1)\)。
\end{lemma}
\begin{proof}
	取 \(0<\tau<1/4\)，\(h=\tau d(x)\)。沿单位向量 \(e\) 作二阶中心差分，Taylor 公式的积分余项给出
	\[
		\left|\partial_e^2u(x)-
		\frac{u(x+he)-2u(x)+u(x-he)}{h^2}\right|
		\le C h^\alpha[\mathrm D^2u]_{\alpha;B_h(x)}.
	\]
	球内的边界距离至少为 \((1-\tau)d(x)\)，故乘以 \(d(x)^2\) 后取上确界，得到
	\(M_2\le C\tau^\alpha H_2+C\tau^{-2}M_0\)。这里先取坐标方向，再取 \((e_i+e_j)/\sqrt2\)，用极化恒等式恢复混合导数。

	一阶差分同理给出 \(M_1\le C\tau M_2+C\tau^{-1}M_0\)。先选 \(\tau\) 使 \(H_2\) 的系数足够小，再吸收固定有限维常数，即得式\eqref{g3:eq:weighted-holder-interpolation}。

	为估计 \(H_1\)，若 \(|x-y|\ge d_{xy}/2\)，用两个一阶导数的上确界；若 \(|x-y|<d_{xy}/2\)，连接线段位于 \(\Omega\)，沿线段积分二阶导数，且其边界距离与 \(d_{xy}\) 可比。两种估计分别给出 \(CM_1\) 与 \(CM_2\)。对 \(H_0\) 用同样的零阶、一阶论证即可。
\end{proof}

在两个严格嵌套的球上，这给出通常的局部插值形式。例如固定 \(B_r\Subset B_R\) 后，\(\|\mathrm Du\|_\infty\)、\(\|\mathrm D^2u\|_\infty\) 及所需低阶 Hölder 半范数，都由较大球上最高阶半范数的任意小倍数与 \(\|u\|_\infty\) 控制；常数同时记录两球之间的距离。边界距离权重把全部内球的估计装进一个有限量中，便于随后一次完成吸收。

\subsection{变系数估计与高阶自举}
\label{g3:subsec:1205:04}
现在考虑
\begin{equation}\label{g3:eq:schauder-operator}
	Lu=a^{ij}(x)\partial_i\partial_ju+b^i(x)\partial_i u+c(x)u.
\end{equation}
主系数 \(a^{ij}\) 为实标量，对称且一致正定；\(b^i,c\) 可为 \(N\times N\) 实矩阵。因而这一估计适用于主符号为正定标量乘恒等的二阶系统，低阶耦合保留在右侧进行估计。

\begin{theorem}[内部 Schauder 估计]\label{g3:thm:interior-schauder}
	设 \(n,N\ge1\)、\(0<\alpha<1\)，式\eqref{g3:eq:schauder-operator}的系数在 \(B_2\) 上属于 \(C^{0,\alpha}\)，其范数至多为 \(K\)，且
	\(\lambda I\le(a^{ij})\le\Lambda I\)，\(\lambda>0\)。若 \(u\in C^{2,\alpha}(\overline B_2;\mathbb R^N)\)，则
	\begin{equation}\label{g3:eq:interior-schauder}
		\|u\|_{C^{2,\alpha}(B_1)}
		\le C\bigl(\|Lu\|_{C^{0,\alpha}(B_2)}+\|u\|_{L^\infty(B_2)}\bigr),
	\end{equation}
	其中 \(C=C(n,N,\alpha,\lambda,\Lambda,K)\)。同一结论也适用于 \(u\in C^2(B_2)\)、\(Lu\in C^{0,\alpha}(B_2)\)：这时结论还包括 \(u\in C^{2,\alpha}_{\mathrm{loc}}(B_2)\)。
\end{theorem}
\begin{proof}
	先作已知 \(C^{2,\alpha}\) 的先验估计。在 \(B_2\) 上用上一引理的 \(M_j,H_2\)。固定 \(x\)，令 \(r=\theta d(x)\)，其中 \(0<\theta<1/8\) 待定。把主系数冻结在 \(x\)，有
	\[
		a^{ij}(x)\partial_i\partial_j u
		=f+(a^{ij}(x)-a^{ij})\partial_i\partial_j u
		-b^i\partial_i u-cu,\qquad f=Lu.
	\]
	对 \(B_r(x)\) 使用常系数估计，以 \(B_{r/2}(x)\) 为内球。冻结误差对二阶半范数产生两项：
	\[
		Cr^{2+\alpha}[a]_\alpha\|\mathrm D^2u\|_{0;B_r(x)}
		+Cr^{2+2\alpha}[a]_\alpha[\mathrm D^2u]_{\alpha;B_r(x)}.
	\]
	第一项来自系数自身的 Hölder 增量，第二项来自 \(\|a-a(x)\|_\infty\le r^\alpha[a]_\alpha\)。这一区别决定了两者的吸收方式。

	取任意两点，令较小边界距离对应点为 \(x\)。若 \(|x-y|<\theta d(x)/4\)，它们位于上述内球。将半范数估计乘以 \((d(x)/r)^{2+\alpha}\)，冻结误差分别不超过 \(CKM_2\) 与 \(CK\theta^\alpha H_2\)，因为 \(d\le2\) 且球内边界距离一致可比。若两点距离较大，直接由 \(M_2\) 得到 \(C\theta^{-\alpha}M_2\)。低阶项利用乘法公式与 \(H_1\le C(M_1+M_2)\)、\(H_0\le C(M_0+M_1)\) 控制。取所有两点的上确界，遂有
	\[
		H_2\le C_\theta\bigl(\|f\|_{C^{0,\alpha}(B_2)}+M_0\bigr)
		+C_{\theta,K}(M_1+M_2)+CK\theta^\alpha H_2.
	\]
	先选 \(\theta\) 使最后一项系数不超过 \(1/4\)，再在插值式\eqref{g3:eq:weighted-holder-interpolation}中选 \(\varepsilon\)，使中间项贡献的 \(H_2\) 系数也不超过 \(1/4\)。移项后，\(M_1+M_2+H_2\) 均受 \(C(\|f\|_{C^{0,\alpha}}+M_0)\) 控制。\(B_1\) 内边界距离至少为一，便得到所述内估计。

	还须把经典 \(C^2\) 解提升到 \(C^{2,\alpha}\)，才能取消预设的半范数。固定 \(B_R\Subset B_2\)，在稍大的球内卷积 \(u_\varepsilon=\rho_\varepsilon*u\)。其方程右端等于 \(f_\varepsilon\) 加有限个形如
	\[
		R_\varepsilon(x)=\int\rho_\varepsilon(x-z)
		\bigl(a(x)-a(z)\bigr)v(z)\,\mathrm dz
	\]
	的项，其中 \(v\) 是 \(u\) 的至多二阶导数，在该较大球上有界。它们满足
	\begin{equation}\label{g3:eq:holder-mollifier-commutator}
		\|R_\varepsilon\|_{C^{0,\alpha}(B_R)}
		\le C\|a\|_{C^{0,\alpha}}\|v\|_\infty,
	\end{equation}
	且常数不依赖小 \(\varepsilon\)。证明这一点时，上确界由 \(C[a]_\alpha\varepsilon^\alpha\|v\|_\infty\) 控制；对于两点差 \(h\le\varepsilon\)，先分离 \(a(x)-a(x')\)，剩余差放到平滑核上，其 \(L^1\) 范数至多为 \(Ch/\varepsilon\)，而系数差至多为 \(C[a]_\alpha\varepsilon^\alpha\)。于是增量不超过 \(C[a]_\alpha\|v\|_\infty(h^\alpha+h\varepsilon^{\alpha-1})\le Ch^\alpha[a]_\alpha\|v\|_\infty\)。\(h>\varepsilon\) 时用上确界即可。

	对光滑 \(u_\varepsilon\) 使用已证估计，得到更小球上一致的 \(C^{2,\alpha}\) 界。又因 \(u\in C^2\)，这些卷积在 \(C^2\) 中局部收敛，故极限保留 \(\alpha\)-半范数界。先这样得到局部 Hölder 正则性，再在逐渐增大的内球上应用先验估计；其半径可保持在固定正下界以上，令外球趋于 \(B_2\)，得到式\eqref{g3:eq:interior-schauder}。
\end{proof}

\begin{corollary}\label{g3:cor:schauder-higher-bootstrap}
	设 \(\Omega\subset\mathbb R^n\) 开，\(L\) 为式\eqref{g3:eq:schauder-operator}的算子，主部在紧子集上一致正定。若 \(a,b,c,f\in C^{k,\alpha}_{\mathrm{loc}}(\Omega)\)，\(k\ge0\) 为整数、\(0<\alpha<1\)，且 \(u\in C^2(\Omega;\mathbb R^N)\) 满足 \(Lu=f\)，则 \(u\in C^{k+2,\alpha}_{\mathrm{loc}}(\Omega)\)。任意 \(\Omega'\Subset\Omega''\Subset\Omega\) 上有
	\[
		\|u\|_{C^{k+2,\alpha}(\Omega')}
		\le C\bigl(\|f\|_{C^{k,\alpha}(\Omega'')}+\|u\|_{L^\infty(\Omega'')}\bigr),
	\]
	常数依赖这些集合、椭圆性及所列有限阶系数界。
\end{corollary}
\begin{proof}
	\(k=0\) 已证。假定系数与 \(f\) 多一阶 Hölder 导数，并已知 \(u\in C^{2,\alpha}\)。令 \(\delta_hu(x)=(u(x+he_l)-u(x))/h\)。它满足
	\[
		L(\delta_hu)=\delta_hf
		-(\delta_ha^{ij})\partial_i\partial_j u(x+he_l)
		-(\delta_hb^i)\partial_i u(x+he_l)
		-(\delta_hc)u(x+he_l).
	\]
	\(\delta_ha=\int_0^1\partial_la(x+the_l)\,\mathrm dt\)，故右端在固定内球的 \(C^{0,\alpha}\) 中一致有界；\(\|\delta_hu\|_\infty\le\|\partial_lu\|_\infty\)。对差商应用先验估计，得到一致的 \(C^{2,\alpha}\) 界。有限网与等度连续性给出各阶导数一致收敛的子列，函数极限是 \(\partial_lu\)，沿线段积分识别其导数，且极限保留 \(\alpha\)-半范数界。因此 \(u\in C^{3,\alpha}\)。继续对方程求导，最高阶项仍为 \(L\) 作用于该导数，其余项由乘积公式和已经获得的阶数控制；有限次归纳给出结论与范数估计。
\end{proof}

\begin{corollary}\label{g3:cor:schauder-distribution-smooth-coefficients}
	设 \(L\) 的主部如上，全部系数在开集 \(\Omega\) 上光滑，\(u\) 为 \(\mathbb R^N\) 值分布。若 \(Lu=f\in C^{k,\alpha}_{\mathrm{loc}}(\Omega)\)，其中整数 \(k\ge0\)、\(0<\alpha<1\)，则 \(u\in C^{k+2,\alpha}_{\mathrm{loc}}(\Omega)\)。
\end{corollary}
\begin{proof}
	由\cref{g3:thm:elliptic-distribution-lp}，\(u\in W^{2,p}_{\mathrm{loc}}\) 对每个有限 \(p>1\) 成立。取 \(p\) 足够大，使 \(1-n/p>\alpha\)；\cref{g3:cor:higher-sobolev-embedding}的局部版本给出 \(u\in C^{1,\alpha}_{\mathrm{loc}}\)。这时卷积交换子中的 \(\mathrm D^2u\) 尚未有界，但可以把一个导数移到核和系数上。具体地，若 \(v=\partial_ju\)，则
	\[
	\begin{aligned}
		[a,\rho_\varepsilon*]\partial_i v(x)
		={}&\int\partial_i\rho_\varepsilon(y)\,y_l
		\int_0^1\partial_la(x-ty)v(x-y)\,\mathrm dt\,\mathrm dy\\
		&+\int\rho_\varepsilon(y)\partial_i a(x-y)v(x-y)\,\mathrm dy.
	\end{aligned}
	\]
	此式由分部积分及 \(a(x)-a(x-y)=\int_0^1\partial_la(x-ty)y_l\,\mathrm dt\) 得到。核 \(y_l\partial_i\rho_\varepsilon\) 的 \(L^1\) 范数一致有界；乘法估计因此把上式的 \(C^{0,\alpha}\) 范数控制在 \(C\|a\|_{C^{1,\alpha}}\|u\|_{C^{1,\alpha}}\) 内。低阶交换子更直接。于是 \(Lu_\varepsilon\) 在 \(C^{0,\alpha}\) 中局部一致有界，先验估计给出 \(u_\varepsilon\) 的局部 \(C^{2,\alpha}\) 界。取子列、识别分布极限得到 \(u\in C^{2,\alpha}\)，再用上一推论完成高阶提升。
\end{proof}

对光滑度量及光滑丛联络，Laplace 型算子在局部标架中具有这样的标量主部，故上述估计可逐图使用。一般椭圆系统的 \(L^p\) 理论已经在前节建立；本节的 Hölder 证明则使用了正定实标量基本解，主部若为任意矩阵系统，应另核其对应核的 Hölder 估计。
% 纲要条目排为编号小节标题，并进入目录与书签；说明段仍保留于纲要注释中。

% !TeX root = ../../Book3.tex
% 纲要标题：### 12.6 一阶 Cauchy–Riemann 型算子
% 纲要位置：几何学计划纲要.md，第 2592 行。
% 正文状态：初稿；尚未正式审核。

% 纲要内容开始
% 
% * 椭圆主符号与局部 W¹ᵖ 估计
% * p>2 时的连续性及非线性复合
% * 闭曲面与条带上的全局条件
% 
% 纲要内容结束

\section{一阶 Cauchy–Riemann 型算子}
\label{g3:sec:1206}

在二维区域中，一阶方程已经可以是椭圆的。复值函数的方程 \(\partial_su+\mathrm i\partial_tu=f\) 同时控制两个方向的变化；其齐次解是全纯函数。几何中把乘以 \(\mathrm i\) 换成纤维上的复结构矩阵，就得到 Cauchy--Riemann 型算子。\footnote{Augustin Louis Cauchy（1789--1857），法国数学家，在实与复分析、微分方程及数学物理中作出基础性贡献。}本节用主符号直接核对椭圆性，再把上一节的估计转为后续曲面方程可以使用的局部结论。

\subsection{\texorpdfstring{椭圆主符号与局部 \(W^{1,p}\) 估计}{椭圆主符号与局部 W¹ᵖ 估计}}
\label{g3:subsec:1206:01}
设 \(J_0\) 是 \(\mathbb R^{2r}\) 上的线性复结构，即 \(J_0^2=-I\)。对
\[
	D_{J_0}u=\partial_su+J_0\partial_tu
\]
及非零实余向量 \(\xi=(\sigma,\tau)\)，有
\[
	(\sigma I-\tau J_0)(\sigma I+\tau J_0)=(\sigma^2+\tau^2)I.
\]
按 \(\mathrm D=-\mathrm i\partial\) 的约定，主符号是 \(\mathrm i(\sigma I+\tau J_0)\)，其逆为 \(-\mathrm i(\sigma I-\tau J_0)/|\xi|^2\)。取相容的实标架可把 \(J_0\) 写成 \(r\) 个标准二维复结构块，但上述计算已经给出不依赖这一换基的可逆性验证。

\begin{theorem}\label{g3:thm:cr-local-w1p}
	设 \(\Omega\subset\mathbb R^2\) 开，\(J:\Omega\to\operatorname{End}(\mathbb R^{2r})\) 光滑且 \(J^2=-I\)，\(A:\Omega\to\operatorname{End}(\mathbb R^{2r})\) 光滑。定义
	\[
		D_Ju=\partial_su+J(s,t)\partial_tu+A(s,t)u.
	\]
	若 \(1<p<\infty\)、\(U\Subset V\Subset\Omega\)，则
	\[
		\|u\|_{W^{1,p}(U)}\le C\bigl(\|D_Ju\|_{L^p(V)}+\|u\|_{L^p(V)}\bigr)
	\]
	对 \(u\in W^{1,p}(V;\mathbb R^{2r})\) 成立。若 \(u\) 仅为分布且 \(D_Ju\in W^{k,p}_{\mathrm{loc}}\)，则 \(u\in W^{k+1,p}_{\mathrm{loc}}\)，其中 \(k\ge0\) 为整数。
\end{theorem}
\begin{proof}
	把上面的逆矩阵公式逐点用于 \(J(s,t)\)。在固定紧集上，\(\|J\|\) 有界，故单位余球上的主符号逆有统一界。零阶项 \(A\) 不改变主符号。将实系统复化，直接应用\cref{g3:thm:local-elliptic-lp-apriori}及\cref{g3:thm:elliptic-distribution-lp}，再限制回实截面，便得全部结论。对于已有 \(W^{1,p}\) 正则性的先验不等式，证明实际上只需 \(J\) 连续、\(A\) 局部有界；从任意分布恢复正则性的陈述保留光滑系数条件。
\end{proof}

\begin{example}\label{g3:exm:cr-factorization}
	当 \(J_0\) 常值且 \(A=0\) 时，
	\[
		(\partial_s-J_0\partial_t)D_{J_0}
		=\partial_s^2+\partial_t^2.
	\]
	若 \(J_0\) 还对 Euclid 内积正交，则 \(J_0^*=-J_0\)，形式伴随为 \(D_{J_0}^*=-\partial_s+J_0\partial_t\)，从而
	\(D_{J_0}^*D_{J_0}=-\partial_s^2-\partial_t^2\)，符号与本卷非负 Laplace 约定一致。若 \(J\) 随点变化，展开乘积会产生 \(\partial_sJ\)、\(\partial_tJ\) 乘一阶导数的项，主部仍是同一个二阶标量 Laplace 主部。
\end{example}

在光滑曲面上，Cauchy--Riemann 型算子的局部表达可以再乘一个光滑可逆的目标丛矩阵；将它移到右端不会改变局部估计。在重叠平凡化中，换标架的导数只增加零阶项。因而上述结论可以用于曲面上有限秩丛的截面，无需把丛整体平凡化。

\subsection{\texorpdfstring{\(p>2\) 时的连续性及非线性复合}{p>2 时的连续性及非线性复合}}
\label{g3:subsec:1206:02}
曲面维数为二，故 \(p>2\) 时 \(\alpha=1-2/p>0\)。由\cref{g3:thm:local-morrey}，在内部小球上每个 \(W^{1,p}\) 函数都有唯一连续代表，且
\[
	\|u\|_{C^{0,\alpha}(B_r)}\le C\|u\|_{W^{1,p}(B_{2r})}.
\]
这里的唯一性来自两个连续代表若几乎处处相等，就处处相等。这使得条件“\(u\) 的像包含于目标坐标图”成为稳定的条件，而不只是对某个可测代表的选择。

\begin{proposition}\label{g3:prop:surface-w1p-composition}
	设 \(\Omega\subset\mathbb R^2\) 为有界光滑域，\(p>2\)，\(u\in W^{1,p}(\Omega;\mathbb R^a)\)，其连续代表的像位于某个紧集 \(K\) 中。若 \(F\) 在 \(K\) 的开邻域为 \(C^2\)，则 \(F(u)\in W^{1,p}\)，且
	\[
		\partial_j(F(u))=\mathrm DF(u)\partial_j u.
	\]
	在满足像位于同一 \(K\)、\(W^{1,p}\) 范数至多为 \(R\) 的函数中，复合映射局部 Lipschitz：对充分接近的 \(u,v\)，
	\[
		\|F(u)-F(v)\|_{W^{1,p}}\le C_{K,F,R,\Omega,p}\|u-v\|_{W^{1,p}}.
	\]
\end{proposition}
\begin{proof}
	在 \(K\) 的邻域取一个紧支截断，将 \(F\) 延拓为全空间 \(C^2\) 函数，使它在 \(K\) 附近不变，且前两阶导数有界。\Cref{g3:prop:sobolev-product-chain}给出弱链式法则。对差值，积分中值公式和链式法则分别给出
	\[
		|F(u)-F(v)|\le C|u-v|,
	\]
	\[
		\partial_j(F(u)-F(v))
		=\mathrm DF(u)\partial_j(u-v)
		+[\mathrm DF(u)-\mathrm DF(v)]\partial_jv.
	\]
	第二项的 \(L^p\) 范数至多为 \(C\|u-v\|_\infty\|\partial_jv\|_p\)。在有界光滑域上，延拓与 Morrey 估计给出 \(\|u-v\|_\infty\le C\|u-v\|_{W^{1,p}}\)，从而得到所需界。若仅在给定开邻域中工作，一致范数足够小就保证连接线段仍落在这一邻域，得到同样的局部结论。
\end{proof}

作为具体应用，若目标坐标中的光滑复结构写为 \(J(z)\)，则
\[
	\mathcal N(u)=\partial_su+J(u)\partial_tu
\]
在像位于固定坐标紧集的 \(W^{1,p}\) 映射上取值于 \(L^p\)，并在有界小邻域内局部 Lipschitz。事实上，差值中的新项为 \([J(u)-J(v)]\partial_tv\)，由上一证明的同一估计控制。线性化沿一个光滑映射 \(u\) 的表达为
\[
	\xi\longmapsto\partial_s\xi+J(u)\partial_t\xi
	+\bigl(\mathrm DJ(u)[\xi]\bigr)\partial_tu,
\]
主符号仍是 Cauchy--Riemann 型，最后一项为零阶项。这一计算说明 \(p>2\) 如何使复合与乘法进入方程；映射空间中的连续微分及高阶光滑性将在下一章按明确的范数和目标空间证明。

\subsection{闭曲面与条带上的全局条件}
\label{g3:subsec:1206:03}
在闭曲面上，有限覆盖可以把局部估计拼成全局估计。边界及无穷远端则引入新的条件，不能由内部小球的估计消除。

\begin{proposition}\label{g3:prop:cr-closed-estimate}
	设 \(\Sigma\) 为紧、光滑、无边界曲面，\(E,F\) 为等秩光滑实或复度量向量丛，\(D:C^\infty(E)\to C^\infty(F)\) 为光滑 Cauchy--Riemann 型一阶算子。对 \(1<p<\infty\)，
	\[
		\|u\|_{W^{1,p}(\Sigma;E)}\le C\bigl(\|Du\|_{L^p(\Sigma;F)}+\|u\|_{L^p(\Sigma;E)}\bigr).
	\]
\end{proposition}
\begin{proof}
	取有限个相对紧平凡化坐标片，并选从属的光滑单位分解 \(\chi_i\)。对 \(\chi_i u\) 使用局部估计。由于 \(D\) 是一阶算子，\([D,\chi_i]u\) 仅含 \(\mathrm d\chi_i\) 乘 \(u\)，因而受全局零阶范数控制。有限求和与局部 Sobolev 范数等价性给出不等式。先对光滑截面证明，再由\cref{g3:thm:sobolev-smooth-density}延伸到 \(W^{1,p}\)。
\end{proof}

零阶范数在这里有实际作用。例如平坦复环面上常数函数满足 \(\bar\partial u=0\)，所以不能无条件删去 \(\|u\|_p\)。核、余核和可解条件须结合全局论证确定。

\ProseParagraphBreak

带边界区域的内部估计不能控制全部边界附近的变化。在单位圆盘取 \(u_k(z)=z^k\)，则 \(\bar\partial u_k=0\)，而极坐标积分给出
\[
	\|u_k\|_p^p=\frac{2\pi}{kp+2},\qquad
	\|\mathrm Du_k\|_p^p=(\sqrt2k)^p\frac{2\pi}{(k-1)p+2}.
\]
所以梯度与函数范数之比随 \(k\) 线性增大。这排除了没有边界条件的整圆盘 \(W^{1,p}\) 估计；在任何严格内部的小圆盘上，上一局部定理仍成立。

\ProseParagraphBreak

在无限条带 \(\mathbb R\times(0,1)\) 上，除两侧边界条件外，还要规定无穷远端的函数空间。以 \(D=\partial_s+\mathrm i\partial_t\) 为例，取非零 \(\eta\in C_c^\infty(\mathbb R;\mathbb R)\)，令
\[
	u_R(s,t)=R^{-1/p}\eta(s/R).
\]
则 \(\|u_R\|_p=\|\eta\|_p\)，而 \(\|Du_R\|_p=R^{-1}\|\eta'\|_p\to0\)。这些函数在两边界都取实值，说明即使给出这种实边界条件，也仍可能在长尺度上出现近零模式。加权空间、两端渐近算子与相应谱条件会改变全局可解性；本节的内部估计负责局部正则性，并不替代这些全局条件。
% 纲要条目排为编号小节标题，并进入目录与书签；说明段仍保留于纲要注释中。

% !TeX root = ../../Book3.tex
% 纲要标题：### 12.7 标量最大值原理与 Harnack 不等式
% 纲要位置：几何学计划纲要.md，第 2598 行。
% 正文状态：初稿；尚未正式审核。

% 纲要内容开始
% 
% * 弱最大值原理及零阶项符号
% * 强最大值原理和屏障
% * Hopf 边界引理的条件
% * Moser 迭代与 Harnack 不等式
% * 椭圆系统中最大值原理的失效
% 
% ---
% 
% 纲要内容结束

\section{标量最大值原理与 Harnack 不等式}
\label{g3:sec:1207}

对实值函数，可以比较两个值的大小，也可以把方程与正部、负幂或对数配对。这些操作产生了前面系统估计之外的工具：边界值控制内部极值，非负解的大小在内球之间相互控制。本节先处理经典解的屏障，再用积分与迭代证明只需可测系数的 Harnack 不等式。\footnote{Karl Gustav Axel von Harnack（1851--1888），德国数学家，研究势论、实分析与代数曲线；正调和函数的比较不等式以其命名。}最后用一个可直接求解的耦合系统说明，标量的次序结构是这些证明的实质条件。

\subsection{弱最大值原理及零阶项符号}
\label{g3:subsec:1207:01}
本小节的经典算子写作
\begin{equation}\label{g3:eq:scalar-maximum-operator}
	Lu=a^{ij}\partial_i\partial_ju+b^i\partial_i u+cu,
\end{equation}
其中 \(a^{ij}=a^{ji}\) 为实数，\(a^{ij}\xi_i\xi_j\ge\lambda|\xi|^2\)，\(\lambda>0\)。这个符号下，\(\sum_i\partial_i^2=-\Delta_{\mathrm{Eucl}}\)；因此 \(Lu\ge0\) 对应通常的下调和方向。

\begin{theorem}[弱最大值原理]\label{g3:thm:weak-maximum-classical}
	设 \(\Omega\subset\mathbb R^n\) 为有界开集，\(Lu=a^{ij}\partial_i\partial_ju+b^i\partial_i u+cu\)，其中 \(a^{ij}=a^{ji},b^i,c\) 为实系数，在 \(\Omega\) 连续且有界，并满足 \(a^{ij}(x)\xi_i\xi_j\ge\lambda|\xi|^2\) 对所有 \(x\in\Omega,\xi\in\mathbb R^n\) 成立，其中 \(\lambda>0\)，\(c\le0\)。若实值 \(u\in C^2(\Omega)\cap C^0(\overline\Omega)\) 且 \(Lu\ge0\)，则
	\[
		\sup_\Omega u\le\max\{0,\sup_{\partial\Omega}u\}.
	\]
	若 \(c=0\)，右侧可直接取 \(\sup_{\partial\Omega}u\)。
\end{theorem}
\begin{proof}
	先设 \(Lu>0\)。若一个正的最大值出现在内部，则该处 \(\mathrm Du=0\)、\(\mathrm D^2u\) 半负定。把正定矩阵 \(a\) 对角化，得 \(a^{ij}u_{ij}\le0\)，又有 \(cu\le0\)，与 \(Lu>0\) 矛盾。

	一般情形取 \(\psi(x)=e^{\kappa x_1}\)。当 \(\kappa\) 足够大时，
	\[
		L\psi=e^{\kappa x_1}(\kappa^2a^{11}+\kappa b^1+c)>0,
	\]
	因为正的二次项压过有界的一次项和零阶项。对 \(u+\varepsilon\psi\) 使用严格情形，再令 \(\varepsilon\downarrow0\)，得到结论。\(c=0\) 时可先减去边界上确界，算子不变。
\end{proof}

应用于两个函数之差便得到比较原理。零阶项的符号不能删去：在 \((0,\pi)\) 上，\(u=\sin x\) 满足 \(u''+u=0\)，边界值为零，内部却为正。即使 \(c\le0\)，若不假定内部最大值非负，也不能在强最大值结论中随意消去零值基准；例如 \(-\cosh x\) 满足 \(u''-u=0\)，在零点达到负的严格最大值。

\begin{proposition}\label{g3:prop:weak-maximum-energy}
	设 \(\Omega\subset\mathbb R^n\) 为有界开集，实矩阵场 \(A\in L^\infty(\Omega)\) 满足
	\(\xi\cdot A(x)\xi\ge\lambda|\xi|^2\) 几乎处处，\(\lambda>0\)，并设 \(c\in L^\infty(\Omega)\)、\(c\ge0\)。若 \(u\in H^1(\Omega;\mathbb R)\)、\(u_+\in H_0^1(\Omega)\)，且
	\[
		\int_\Omega A\mathrm Du\cdot\mathrm D\phi+cu\phi\,\mathrm dx\le0
	\]
	对每个非负 \(\phi\in H_0^1(\Omega)\) 成立，则 \(u\le0\) 几乎处处。
\end{proposition}
\begin{proof}
	取 \(\phi=u_+\)。第十一章的截断链式法则给出 \(\mathrm Du_+=\mathbf1_{\{u>0\}}\mathrm Du\)，所以
	\[
		0\ge\int_\Omega A\mathrm Du_+\cdot\mathrm Du_++c u_+^2\,\mathrm dx
		\ge\lambda\|\mathrm Du_+\|_2^2.
	\]
	零边界 Poincaré 不等式\cref{g3:thm:poincare-zero-boundary}于是给出 \(u_+=0\)。\(A\) 无须对称，因为这里只使用二次型的下界。
\end{proof}

这里采用非负主部 \(-\operatorname{div}(A\mathrm D)+c\)，所以零阶条件写作 \(c\ge0\)，与前一定理的符号完全相容。条件 \(u_+\in H_0^1\) 把边界上“不大于零”写成可用于测试的函数空间条件；边界有第十一章所需正则性时，它与非正的迹相符。

\subsection{强最大值原理和屏障}
\label{g3:subsec:1207:02}
\term[ruo zuidazhi yuanli]{弱最大值原理}[weak maximum principle]只给出大小界。要排除内部一处接触最大值、其他位置仍较小的情形，需要一个从小球边缘向内严格增加的函数。

\begin{lemma}\label{g3:lem:exponential-annular-barrier}
	设 \(R>0\)、\(y\in\mathbb R^n\)，\(Lu=a^{ij}\partial_i\partial_ju+b^i\partial_i u+cu\)，其中 \(a^{ij}=a^{ji},b^i,c\) 为实系数，其系数在球壳 \(R/2<|x-y|<R\) 上有界，主部统一正定，\(c\le0\)。当 \(\gamma\) 足够大时，
	\[
		w(x)=e^{-\gamma|x-y|^2}-e^{-\gamma R^2}
	\]
	在球壳内为正，在外球面为零，并满足 \(Lw>0\)。外球面上的径向外导数为 \(-2\gamma R e^{-\gamma R^2}<0\)。
\end{lemma}
\begin{proof}
	直接求导得
	\[
		Lw=e^{-\gamma|x-y|^2}
		\bigl(4\gamma^2 a^{ij}(x_i-y_i)(x_j-y_j)
		-2\gamma\operatorname{tr}a-2\gamma b\cdot(x-y)\bigr)+cw.
	\]
	在此球壳上，括号内首项至少为 \(\gamma^2\lambda R^2\)，其余两项至多按 \(C\gamma\) 增长；又有 \(cw\ge-\|c\|_\infty e^{-\gamma|x-y|^2}\)。故足够大的 \(\gamma\) 保证严格正性，边界导数的公式直接成立。
\end{proof}

\begin{theorem}[强最大值原理]\label{g3:thm:strong-maximum-classical}
	设 \(\Omega\subset\mathbb R^n\) 为连通开集，\(Lu=a^{ij}\partial_i\partial_ju+b^i\partial_i u+cu\)，其中 \(a^{ij}=a^{ji},b^i,c\) 为实系数，在 \(\Omega\) 连续、在紧子集上有界，主部在紧子集上一致正定，\(c\le0\)。若 \(u\in C^2(\Omega;\mathbb R)\)、\(Lu\ge0\)，且 \(u\) 在某内点达到整个 \(\Omega\) 上的最大值 \(M\ge0\)，则 \(u\equiv M\)。若 \(c=0\)，不要求 \(M\ge0\)。
\end{theorem}
\begin{proof}
	令 \(v=M-u\ge0\)。由 \(M\ge0\)、\(c\le0\)，有 \(Lv=cM-Lu\le0\)。假设 \(v\) 不恒为零。沿连接零点与正值点的紧路径，在正值集合中选一个足够靠近其内部边界的点 \(y\)。从 \(y\) 到零集的距离 \(R\) 为正且足够小，使 \(\overline B_R(y)\subset\Omega\)；最近零点 \(x_0\) 位于球面，而球内 \(v>0\)。最近点存在，因为可把距离问题限制在包含该路径的紧邻域中。

	在 \(R/2<|x-y|<R\) 上取上一引理的 \(w\)。内球面上 \(v\) 有正下界，故可选 \(\varepsilon>0\)，使 \(v-\varepsilon w\ge0\) 在两球面上成立；且 \(L(v-\varepsilon w)<0\)。弱最大值原理应用于 \(-(v-\varepsilon w)\)，得到整个球壳内 \(v\ge\varepsilon w\)。令 \(\nu=(x_0-y)/R\)，沿 \(x_0-t\nu\) 趋向 \(x_0\)，得
	\[
		\liminf_{t\downarrow0}\frac{v(x_0-t\nu)-v(x_0)}t
		\ge2\varepsilon\gamma R e^{-\gamma R^2}>0.
	\]
	但 \(x_0\) 是 \(\Omega\) 内部的 \(v\) 的最小点，\(\mathrm Dv(x_0)=0\)，矛盾。因此正值集合为空。\(c=0\) 时，\(Lv=-Lu\le0\) 对任意 \(M\) 都成立。
\end{proof}

\term[qiang zuidazhi yuanli]{强最大值原理}[strong maximum principle]的证明只在一个内接球壳上使用统一常数，因此不要求区域有界。连通性把一个内点的极值传播到整个区域；不连通时只能得到所在连通分支上的结论。

\subsection{Hopf 边界引理的条件}
\label{g3:subsec:1207:03}
若上述接触点位于区域边界，梯度不必为零，球壳论证留下一个严格的外法向导数。这正是\term[Hopf bianjie yinli]{Hopf 边界引理}[Hopf boundary point lemma]的内容。\footnote{Eberhard Frederich Ferdinand Hopf（1902--1983），奥地利出生的美籍数学家，研究遍历理论、动力系统与偏微分方程；这里的边界引理归于他，与前面 Hopf--Rinow 定理中的 Heinz Hopf 为不同人物。}

\begin{theorem}[Hopf 边界引理]\label{g3:thm:hopf-boundary-lemma}
	设 \(\Omega\subset\mathbb R^n\) 为开集，\(x_0\in\partial\Omega\)，且存在内接球 \(B_R(y)\subset\Omega\)，满足 \(x_0\in\partial B_R(y)\)。设 \(Lu=a^{ij}\partial_i\partial_ju+b^i\partial_i u+cu\)，其中 \(a^{ij}=a^{ji},b^i,c\) 为实系数，在该球附近连续且有界、主部一致正定，\(c\le0\)。若 \(u\in C^2(\Omega)\) 且在 \(\overline B_R(y)\) 连续，\(Lu\ge0\)，并且
	\[
		u(x)<u(x_0)=M\quad(x\in B_R(y)),\qquad M\ge0,
	\]
	则对 \(\nu=(x_0-y)/R\)，有
	\[
		\liminf_{t\downarrow0}\frac{u(x_0)-u(x_0-t\nu)}t>0.
	\]
	特别地，若 \(u\) 在 \(x_0\) 可微，则 \(\partial_\nu u(x_0)>0\)。\(c=0\) 时可删去 \(M\ge0\)。
\end{theorem}
\begin{proof}
	在内接球中令 \(v=M-u>0\)。内球面 \(|x-y|=R/2\) 紧且位于 \(\Omega\)，故 \(v\) 有正下界。与强最大值原理的证明相同，选 \(\varepsilon>0\)，在球壳上比较 \(v\) 与 \(\varepsilon w\)，得到
	\(v(x_0-t\nu)\ge\varepsilon w(x_0-t\nu)\)。除以 \(t\)，令 \(t\downarrow0\)，即得正下界 \(2\varepsilon\gamma R e^{-\gamma R^2}\)。
\end{proof}

若边界为 \(C^2\)，则在每个边界点有足够小的内接球，且球的 \(\nu\) 就是区域外法向；证明实际需要的是内接球条件，而非整个边界的一套光滑图册。图\cref{g3:fig:hopf-barrier}标出这一局部构造。
\begin{figure}[htbp]
	\centering
	% !TeX root = ../Book3.tex
\begin{tikzpicture}[x=1cm,y=1cm,>=stealth]
	\fill[blue!7] (-3,-1.55) rectangle (2,1.55);
	\draw[thick] (2,-1.55)--(2,1.55);
	\node[right] at (2,1.25) {\(\partial\Omega\)};
	\draw[thick] (.7,0) circle (1.3);
	\draw[dashed] (.7,0) circle (.65);
	\fill (.7,0) circle (1.5pt) node[below] {\(y\)};
	\fill (2,0) circle (1.8pt) node[below right] {\(x_0\)};
	\draw[->,thick] (2,0)--(2.85,0) node[right] {\(\nu\)};
	\draw[->] (1.85,.32)--(1.42,.32) node[midway,above] {\(M-u\) 增大};
	\draw[<->] (.7,-.08)--(1.35,-.08);
	\node at (1.02,-.35) {\(R/2\)};
	\node at (-1.85,.95) {\(\Omega\)};
\end{tikzpicture}

	\caption[内接球与 Hopf 屏障]{内接球在 \(x_0\) 接触边界。球壳上的屏障从外球面向内严格增加，因而 \(M-u\) 的内向增长有正下界。}
	\label{g3:fig:hopf-barrier}
\end{figure}

例如在平面第一象限内，\(u(x,y)=-xy\) 调和，在两条坐标边上为零，但顶点的梯度也为零。顶点没有与两条边都相容的内接球，故这一例子并不满足引理的几何条件。再如半区间上的 \(u(x)=-x^2\) 满足 \(xu''-u'=0\)，却有 \(u'(0)=0\)；这里主系数在接触点退化，一致椭圆性已经失效。

\subsection{Moser 迭代与 Harnack 不等式}
\label{g3:subsec:1207:04}
屏障利用二阶点态导数；当主系数只有可测性时，解起初只是 \(H^1\) 函数，必须改用积分。\term[Moser diedai]{Moser 迭代}[Moser iteration]\footnote{Jürgen Kurt Moser（1928--1999），德国出生的美籍数学家，研究 Hamilton 动力系统与偏微分方程，以迭代估计建立了重要的正则性结果。}把一个有限次幂的积分界反复提高到更高次幂，最终得到本性上确界。对正解同时处理正幂与负幂，再以对数的平均振荡连接两者，就得到\term[Harnack budengshi]{Harnack 不等式}[Harnack inequality]。

\ProseParagraphBreak

以下记 \(Q_r(x_0)=x_0+(-r,r)^n\)。在同心立方体上证明估计，是为了让后面的二进分割保持简单；球上的版本可用有限个重叠立方体连接得到。

\begin{lemma}\label{g3:lem:moser-power-estimates}
	设 \(n\ge2\)，记 \(Q_a=(-a,a)^n\)，\(A\) 是 \(Q_4\) 上的实可测矩阵场，满足
	\[
		\xi\cdot A(x)\xi\ge\lambda|\xi|^2,
		\qquad |A(x)\xi|\le\Lambda|\xi|,
		\qquad 0<\lambda\le\Lambda.
	\]
	若 \(u\in H^1(Q_4)\) 非负且弱满足 \(\operatorname{div}(A\mathrm Du)=0\)，令 \(v=u+\varepsilon\)，\(\varepsilon>0\)。则对每个 \(p>0\)，有
	\begin{align}
		\mathop{\mathrm{ess\,sup}}_{Q_1}v
		&\le C_p\left(\frac1{|Q_2|}\int_{Q_2}v^p\right)^{1/p},\label{g3:eq:moser-positive-power}\\
		\mathop{\mathrm{ess\,inf}}_{Q_1}v
		&\ge C_p^{-1}\left(\frac1{|Q_2|}\int_{Q_2}v^{-p}\right)^{-1/p}.\label{g3:eq:moser-negative-power}
	\end{align}
	常数只依赖 \(n,\lambda,\Lambda,p\)，不依赖 \(\varepsilon\)。
\end{lemma}
\begin{proof}
	\(v\) 仍解同一方程，且 \(v\ge\varepsilon\)。先证明正幂的能量估计，并使测试函数的合法性明确。取 \(q\ge2\)、紧支 Lipschitz 截断 \(\eta\)，令
	\[
		F_M(t)=t\min\{t,M\}^{q-2},\qquad
		G_M(t)=t\min\{t,M\}^{q/2-1}.
	\]
	\(F_M\) 在 \([\varepsilon,\infty)\) 上 Lipschitz，故 \(\eta^2F_M(v)\) 是合法的 \(H_0^1\) 测试函数。展开弱方程、用椭圆性及平方配方吸收交叉项，得到
	\[
		\int\eta^2F_M'(v)|\mathrm Dv|^2
		\le C\int\frac{F_M(v)^2}{F_M'(v)}|\mathrm D\eta|^2.
	\]
	在 \(t<M\) 与 \(t>M\) 分别计算，得
	\(|G_M'(t)|^2\le q^2 F_M'(t)\)、\(F_M(t)^2/F_M'(t)\le G_M(t)^2\)。于是
	\[
		\int|\mathrm D(\eta G_M(v))|^2
		\le Cq^2\int G_M(v)^2|\mathrm D\eta|^2
		\le Cq^2\int v^q|\mathrm D\eta|^2.
	\]
	若右端有限，先用 \(M\to\infty\) 和弱紧性识别 \(\eta v^{q/2}\) 的弱导数，再用下半连续性，得到
	\begin{equation}\label{g3:eq:moser-power-energy}
		\int|\mathrm D(\eta v^{q/2})|^2
		\le Cq^2\int v^q|\mathrm D\eta|^2,\qquad q\ge2.
	\end{equation}
	截断过程只在已知 \(v^q\) 可积的那一步使用，随后由 Sobolev 不等式提高下一次可积指数。

	取 \(\chi=n/(n-2)>1\) 当 \(n\ge3\)，在 \(n=2\) 时取 \(\chi=2\)。第十一章的 Euclid Sobolev 不等式给出紧支于 \(Q_2\) 的 \(w\) 满足 \(\|w\|_{2\chi}^2\le C\|\mathrm Dw\|_2^2\)；二维时先用 \(W^{1,4/3}\to L^4\)，再在固定支集上用 Hölder 不等式把梯度换成 \(L^2\)。选 \(\eta=1\) 于 \(Q_\rho\)、支于 \(Q_R\)，\(|\mathrm D\eta|\le C/(R-\rho)\)，则
	\[
		\|v\|_{L^{q\chi}(Q_\rho)}
		\le\left(\frac{Cq^2}{(R-\rho)^2}\right)^{1/q}
		\|v\|_{L^q(Q_R)},\qquad1\le\rho<R\le2.
	\]
	令 \(q_j=2\chi^j\)，在递减半径之间应用此式。因 \(\sum_j1/q_j\) 与 \(\sum_jj/q_j\) 均收敛，常数的乘积有限，得到
	\begin{equation}\label{g3:eq:moser-l2-sup}
		M(\rho):=\mathop{\mathrm{ess\,sup}}_{Q_\rho}v
		\le C(R-\rho)^{-\nu/2}\|v\|_{L^2(Q_R)},
		\qquad \nu=\frac{2\chi}{\chi-1}.
	\end{equation}
	可积指数趋于无穷时，若某个正测度集上的值超过所得统一界，就会与其 \(L^{q_j}\) 界矛盾，故这一步确实得到本性上确界。

	当 \(0<p<2\) 时，记 \(A_p=\|v\|_{L^p(Q_2)}\)。由
	\(\|v\|_{L^2(Q_R)}\le M(R)^{1-p/2}A_p^{p/2}\)，式\eqref{g3:eq:moser-l2-sup}及 Young 不等式给出
	\[
		M(\rho)\le\theta M(R)+C_{p,\theta}(R-\rho)^{-\nu/p}A_p
	\]
	对任意 \(0<\theta<1\) 成立。取 \(\rho_j=3/2-2^{-j-1}\)，选 \(\theta<2^{-\nu/p}\) 后反复代入，几何级数收敛；余项 \(\theta^jM(\rho_j)\) 趋零，因为式\eqref{g3:eq:moser-l2-sup}已使 \(M(3/2)\) 有限。这就得到从任意正幂到上确界的估计。\(p\ge2\) 则由有限体积上的 Hölder 不等式直接降至 \(p=2\)。固定体积因子吸收入 \(C_p\)，得式\eqref{g3:eq:moser-positive-power}。

	对负幂取 \(q<0\)，测试 \(\eta^2v^{q-1}\)。由于 \(v\ge\varepsilon\)，该复合及其导数有界，测试合法；整理弱方程并对 \(q-1\) 取绝对值，得到
	\[
		\int|\mathrm D(\eta v^{q/2})|^2
		\le C\left(1+\frac{q^2}{(q-1)^2}\right)
		\int v^q|\mathrm D\eta|^2.
	\]
	令 \(w=v^{-1}\)，从 \(w^p\) 出发用同一 Sobolev 迭代，得到 \(\|w\|_{\infty;Q_1}\le C_p\|w\|_{p;Q_2}\)。初始积分因 \(w\le\varepsilon^{-1}\) 而有限；迭代常数不使用这个上界。对所得不等式取倒数即为式\eqref{g3:eq:moser-negative-power}。
\end{proof}

正幂与负幂估计还不能直接连接：一般正函数可在大部分区域很小，却在一小块上很大。下面的平均振荡引理排除这种失控，它是此处所需的 John--Nirenberg 估计。\footnote{Fritz John（1910--1994），德国出生的美籍数学家，研究偏微分方程与不适定问题。Louis Nirenberg（1925--2020），加拿大出生的美籍数学家，研究线性与非线性偏微分方程及其几何应用。}

\begin{lemma}\label{g3:lem:dyadic-exponential-oscillation}
	设 \(Q_0\) 为立方体，\(g\in L^1(Q_0;\mathbb R)\)，且对 \(Q_0\) 及它的每个二进子立方体 \(Q\)，都有
	\[
		\frac1{|Q|}\int_Q|g-g_Q|\le B,
		\qquad g_Q=\frac1{|Q|}\int_Qg,
	\]
	其中 \(B>0\)。则存在只依赖 \(n\) 的 \(c_n,C_n>0\)，使
	\[
		\frac1{|Q_0|}\int_{Q_0}
		\exp\left(\frac{c_n}{B}|g-g_{Q_0}|\right)\le C_n.
	\]
\end{lemma}
\begin{proof}
	在 \(Q_0\) 的二进网格中，选取使 \(|g-g_{Q_0}|\) 的平均大于 \(2B\) 的极大子立方体 \(Q_j\)。它们互不相交，且
	\[
		\sum_j|Q_j|\le\frac1{2B}\int_{Q_0}|g-g_{Q_0}|\le\frac12|Q_0|.
	\]
	每个 \(Q_j\) 的父立方体平均不超过 \(2B\)，故
	\(|g_{Q_j}-g_{Q_0}|\le2^{n+1}B\)。在其余点上，二进平均收敛到函数值，因而 \(|g-g_{Q_0}|\le2B\) 几乎处处。

	所用的二进微分性质可由同一最大立方体论证取得：对可积 \(h\)，二进平均的最大值超过 \(t\) 的集合测度不超过 \(\|h\|_1/t\)。用连续函数在 \(L^1\) 中逼近 \(g\)；连续函数的细网格平均趋于其点值，而误差的最大平均及误差本身超过 \(t\) 的集合测度至多为 \(2\|h\|_1/t\)。令逼近误差趋零，再令有理 \(t\downarrow0\)，便得上述几乎处处收敛。

	在每个 \(Q_j\) 内改减去它自己的平均，重复同一选择。第 \(m\) 代坏立方体的总体积至多为 \(2^{-m}|Q_0|\)。沿包含某点的立方体链，连续两个平均之差至多为 \(2^{n+1}B\)；离开第 \(m\) 代坏集合的点因而满足
	\[
		|g-g_{Q_0}|\le(m+1)2^{n+1}B
	\]
	几乎处处。把 \(Q_0\) 按离开坏集合的代数分层，取 \(c_n>0\) 足够小使 \(e^{c_n2^{n+1}}<2\)，则指数积分除以 \(|Q_0|\) 不超过
	\(\sum_{m\ge0}2^{-m}e^{c_n(m+1)2^{n+1}}\)，是收敛的几何级数。
\end{proof}

\termidx[Moser Harnack budengshi]{Moser--Harnack 不等式}
\foreignidx{Moser--Harnack inequality}
\begin{theorem}[Moser--Harnack 不等式]\label{g3:thm:moser-harnack}
	设 \(n\ge2\)、\(r>0\)、\(x_0\in\mathbb R^n\)，记 \(Q_a(x_0)=x_0+(-a,a)^n\)，\(A\) 为 \(Q_{4r}(x_0)\) 上实可测矩阵场，满足
	\(\xi\cdot A\xi\ge\lambda|\xi|^2\)、\(|A\xi|\le\Lambda|\xi|\)，\(\lambda>0\)。若非负 \(u\in H^1(Q_{4r}(x_0))\) 满足
	\[
		\int_{Q_{4r}(x_0)} A\mathrm Du\cdot\mathrm D\phi\,\mathrm dx=0
		\quad\text{对所有 }\phi\in C_c^\infty(Q_{4r}(x_0)),
	\]
	则
	\begin{equation}\label{g3:eq:moser-harnack}
		\mathop{\mathrm{ess\,sup}}_{Q_r(x_0)}u
		\le C(n,\lambda,\Lambda)
		\mathop{\mathrm{ess\,inf}}_{Q_r(x_0)}u.
	\end{equation}
\end{theorem}
\begin{proof}
	缩放后可取 \(r=1\)、\(x_0=0\)。仍令 \(v=u+\varepsilon\)。测试函数 \(\phi=\eta^2/v\) 合法，因为 \(1/v\) 是有界导数的复合；展开弱方程，得
	\[
		\int\eta^2 A\mathrm D(\log v)\cdot\mathrm D(\log v)
		=2\int\eta A\mathrm D(\log v)\cdot\mathrm D\eta.
	\]
	吸收交叉项后有
	\begin{equation}\label{g3:eq:moser-log-energy}
		\int\eta^2|\mathrm D\log v|^2
		\le C(\lambda,\Lambda)\int|\mathrm D\eta|^2.
	\end{equation}
	\(\log v\) 在紧子集上属于 \(H^1\)：函数本身由 \(|\log v|\le C_\varepsilon+v\) 控制，导数为 \(\mathrm Dv/v\)。

	对 \(Q_2\) 的任意二进子立方体 \(Q\)，其同心二倍立方体 \(2Q\) 都含于 \(Q_4\)。选 \(\eta=1\) 于 \(Q\)、支于 \(2Q\)，且 \(|\mathrm D\eta|\le C/\ell(Q)\)。式\eqref{g3:eq:moser-log-energy}给出
	\(\int_Q|\mathrm D\log v|^2\le C\ell(Q)^{n-2}\)。结合凸域 Poincaré 不等式\cref{g3:prop:poincare-convex-mean}及 Hölder 不等式，得到
	\[
		\frac1{|Q|}\int_Q|\log v-(\log v)_Q|
		\le C\ell(Q)\left(\frac1{|Q|}\int_Q|\mathrm D\log v|^2\right)^{1/2}
		\le B,
	\]
	其中 \(B\) 只依赖 \(n,\lambda,\Lambda\)。上一引理遂给出一个固定 \(p_0>0\)，使
	\[
		\left(\frac1{|Q_2|}\int_{Q_2}v^{p_0}\right)
		\left(\frac1{|Q_2|}\int_{Q_2}v^{-p_0}\right)\le C.
	\]
	这里分别把 \(e^{\pm p_0(\log v-(\log v)_{Q_2})}\) 控制在指数绝对振荡之下；两个平均中的常数因子相互抵消。

	将这个乘积界代入正幂估计\eqref{g3:eq:moser-positive-power}与负幂估计\eqref{g3:eq:moser-negative-power}，得到
	\(\mathop{\mathrm{ess\,sup}}_{Q_1}v\le C\mathop{\mathrm{ess\,inf}}_{Q_1}v\)。全部常数不依赖 \(\varepsilon\)，而 \(\operatorname*{ess\,inf}(u+\varepsilon)=\operatorname*{ess\,inf}u+\varepsilon\)，上确界亦然；令 \(\varepsilon\downarrow0\) 即得结论。
\end{proof}

\begin{corollary}\label{g3:cor:measurable-elliptic-holder}
	设 \(\Omega\subset\mathbb R^n\) 开，\(n\ge2\)，\(A:\Omega\to\mathbb R^{n\times n}\) 可测，且存在 \(0<\lambda\le\Lambda\)，使 \(\xi\cdot A(x)\xi\ge\lambda|\xi|^2\)、\(|A(x)\xi|\le\Lambda|\xi|\) 对几乎处处 \(x\) 及所有 \(\xi\) 成立。满足 \(\int_\Omega A\mathrm Du\cdot\mathrm D\phi=0\)（\(\phi\in C_c^\infty(\Omega)\)）的实弱解 \(u\in H^1_{\mathrm{loc}}(\Omega)\) 有局部 Hölder 连续代表。若 \(\Omega\) 连通且此代表非负，则它或者处处为正，或者恒为零。
\end{corollary}
\begin{proof}
	先说明任意符号的弱解也局部有界。在正幂能量证明中，把测试函数改为 \(\eta^2 u\min\{|u|,M\}^{q-2}\)，或者先用 \((u^2+\delta^2)^{1/2}\) 作光滑化，再令 \(\delta\downarrow0\)。按 \(u>0\)、\(u<0\) 两部分计算，其导数同样具有正的主项，得到 \(|u|\) 的式\eqref{g3:eq:moser-power-energy}，故 \(u\) 局部有界。

	在一个紧含的 \(Q_{4r}\) 上记 \(m=\operatorname*{ess\,inf}u\)、\(M=\operatorname*{ess\,sup}u\)。\(u-m\) 与 \(M-u\) 都是非负弱解。设内立方体上的本性上下确界为 \(a,b\)，两次 Harnack 不等式给出
	\[
		b-m\le C(a-m),\qquad M-a\le C(M-b).
	\]
	相加整理，得到
	\[
		\operatorname*{ess\,osc}_{Q_r}u
		\le\vartheta\operatorname*{ess\,osc}_{Q_{4r}}u,
		\qquad \vartheta=\frac{C-1}{C+1}<1.
	\]
	迭代缩小立方体，振荡按半径的 \(\beta=-\log\vartheta/\log4>0\) 次幂衰减（必要时取更小的 \(\beta<1\)）。嵌套本性值域的共同极限定义每个中心点的代表；重叠立方体的振荡界给出其 Hölder 连续性，在几乎处处的微分点它与 \(u\) 相等。最后若非负连续代表有一个零点，Harnack 使邻近内立方体上恒零；零集既开又闭，连通性给出恒零。
\end{proof}

\(n=1\) 时无需迭代。弱方程 \((a u')'=0\) 给出 \(a u'=c\) 几乎处处，故
\(u(t)=u(t_0)+c\int_{t_0}^t a(s)^{-1}\,\mathrm ds\)。由 \(\lambda\le a\le\Lambda\)，在外区间上非负便使内区间满足同型的最大值与最小值比较不等式，其常数只依赖椭圆性比。高维证明中各方向的能量与 Sobolev 提升，正承担了这一维积分公式的作用。

\subsection{椭圆系统中最大值原理的失效}
\label{g3:subsec:1207:05}
丛值和矩阵系统通常没有与换标架相容的“正部”，而一个分量的方程又会受到其他分量的驱动。即使主部是最简单的正 Laplace 型，能量正定也不保证按分量保持次序。

\begin{example}\label{g3:exm:elliptic-system-no-order}
	在 \((0,\pi)\) 上考虑两分量实系统
	\[
		PU=-U''+\begin{pmatrix}2&1\\1&2\end{pmatrix}U
		=\begin{pmatrix}0\\\sin x\end{pmatrix},
		\qquad U(0)=U(\pi)=0.
	\]
	主符号为 \(\xi^2I\)，零阶矩阵的特征值为一和三，所以二次型严格正定。令 \(U(x)=v\sin x\)，方程等价于
	\[
		\begin{pmatrix}3&1\\1&3\end{pmatrix}v
		=\begin{pmatrix}0\\1\end{pmatrix},
		\qquad v=\begin{pmatrix}-1/8\\3/8\end{pmatrix}.
	\]
	右端逐分量非负，边界值为零，解的第一分量却在内部严格为负。由能量法可知该解唯一：齐次解与自身配对后，各非负项之和为零。因此失败来自真实的次序问题，而非解的非唯一性。
\end{example}

对完全解耦的标量方程可以逐分量应用最大值原理；某些具有特定非对角符号的合作系统也有另外的比较定理。一般椭圆系统仍使用本章已经建立的 \(L^2\)、\(L^p\) 与适用主部的 Hölder 估计。标量 Moser 证明中的负幂和对数，不能直接换成一个向量的同名运算；需要对丛值模长另行推出标量微分不等式时，联络和零阶曲率项都必须进入计算。
% 纲要条目排为编号小节标题，并进入目录与书签；说明段仍保留于纲要注释中。


\begin{chaptersummary}
能量方法在完备空间中给出弱解。Lax--Milgram 定理使用有界性与强制性；约束极小问题还须检查弱闭性、弱下半连续性和极小化序列的有界性。零边界 Poincaré 不等式消去常数核，闭流形上则须令右端消灭常数，并固定解的均值。这些条件决定可解性，尚不替代导数估计。

\ProseParagraphBreak

主符号可逆使高频得到与算子阶数相同的导数增益。Fourier 方法先证明常系数 \(L^2\) 估计，冻结与交换子把它局部化；差商直接作用于能量弱解，兼容 Sobolev 阶数的逆算子则从任意分布解开始提升。奇异积分补足 \(L^p\) 乘子控制，\(1<p<\infty\) 时得到相应的内部估计与正则性。Hölder 核估计、带尺度插值及小球吸收建立标量主部的 Schauder 结论，系数正则性决定可继续提升的次数。

\ProseParagraphBreak

Cauchy--Riemann 型算子的实主符号可逆，因而服从一阶系统的局部 \(W^{1,p}\) 理论；二维 \(p>2\) 时的连续嵌入使非线性复合可控。全局估计仍保留零阶项、核以及边界和端部条件。实标量方程另有次序工具：屏障把极值变成严格的法向信息，Moser 迭代与对数平均振荡把正解的上、下界联系起来。耦合系统的显式反例说明，椭圆性与能量正定本身并不保证逐分量比较。
\end{chaptersummary}

\BookExercises

\begin{exercise}\label{g3:ex:skew-energy-equation}
	在 \(H=\mathbb R^2\) 上令
	\[
		A=\begin{pmatrix}1&-b\\ b&1\end{pmatrix},\qquad
		a(u,v)=\langle Au,v\rangle,\qquad f=(1,0),\qquad b\ne0.
	\]
	计算方程 \(a(u,v)=\langle f,v\rangle\) 的解，以及 \(J(v)=\tfrac12a(v,v)-\langle f,v\rangle\) 的极小点。再设 \(\Omega\subset\mathbb R^n\) 为有界开集，\(B\in W^{1,\infty}(\Omega;\mathbb R^n)\)、\(\operatorname{div}B=0\)，\(c\in L^\infty(\Omega)\)、\(c\ge0\)。证明对每个 \(F\in H^{-1}(\Omega)\)，弱方程
	\[
		-\sum_i\partial_i^2u+B\cdot\mathrm Du+cu=F,
		\qquad u\in H_0^1(\Omega;\mathbb R)
	\]
	唯一可解。解释漂移项如何同时影响方程而不改变对角能量。
\end{exercise}
\begin{solution}
	\(A\) 的行列式为 \(1+b^2\)，故解为
	\(u=(1,-b)/(1+b^2)\)。但 \(a(v,v)=|v|^2\)，所以 \(J\) 的唯一极小点为 \((1,0)\)。反对称矩阵被二次能量消去了，而方程的逆矩阵仍含 \(b\)。

	在 \(H_0^1\) 上取
	\[
		a(u,v)=\int_\Omega\bigl(\mathrm Du\cdot\mathrm Dv+(B\cdot\mathrm Du)v+cuv\bigr)\,\mathrm dx.
	\]
	系数有界与 Cauchy--Schwarz 不等式使该型连续。对紧支光滑 \(u\)，分部积分给出
	\(\int(B\cdot\mathrm Du)u=-\tfrac12\int(\operatorname{div}B)u^2=0\)。用 \(H_0^1\) 的定义作光滑逼近，\(B\) 有界使这一恒等式传到所有 \(u\in H_0^1\)。因此
	\(a(u,u)\ge\|\mathrm Du\|_2^2\ge(1+C_P^2)^{-1}\|u\|_{H^1}^2\)。Lax--Milgram 给出唯一解。对一般的不同测试 \(u,v\)，漂移积分并不为零，所以它仍真实地存在于方程中。
\end{solution}

\begin{exercise}\label{g3:ex:galerkin-energy-approximation}
	设 \(H\) 为实或复 Hilbert 空间，\(a\) 满足\cref{g3:thm:lax-milgram}的连续性常数 \(C\) 与强制常数 \(\alpha\)，\(u\) 是右端 \(F\in H^\dagger\) 的解。若 \(V_N\subset H\) 为有限维子空间，证明存在唯一 \(u_N\in V_N\) 使 \(a(u_N,v)=F(v)\) 对所有 \(v\in V_N\) 成立，并证明
	\[
		\|u-u_N\|\le\frac C\alpha\inf_{v\in V_N}\|u-v\|.
	\]
	若 \(a\) 另为 Hermitian，证明在能量范数 \(\|w\|_a=a(w,w)^{1/2}\) 下，上式中的常数可取一。说明嵌套子空间之并稠密时的收敛结论。
\end{exercise}
\begin{solution}
	有限维子空间闭，限制的型保留原来的界，所以 Lax--Milgram 给出 \(u_N\)。相减得到 \(a(u-u_N,w)=0\) 对 \(w\in V_N\) 成立。对任意 \(v\in V_N\)，记 \(e=u-u_N\)，则
	\[
		\alpha\|e\|^2\le\operatorname{Re}a(e,e)
		=\operatorname{Re}a(e,u-v)\le C\|e\|\|u-v\|.
	\]
	若 \(e=0\) 已得结论，否则约去 \(\|e\|\) 再取下确界即可。

	Hermitian 情形中，\(a\) 是等价的 Hilbert 内积，误差与 \(V_N\) 正交。因此对每个 \(v\in V_N\)，
	\(\|u-v\|_a^2=\|u-u_N\|_a^2+\|u_N-v\|_a^2\)，给出最佳逼近性质。若子空间嵌套且并稠密，每个固定精度的逼近向量从某个 \(N\) 起都在 \(V_N\) 中，故右端下确界趋零，\(u_N\to u\) 于 \(H\)。
\end{solution}

\begin{exercise}\label{g3:ex:point-source-energy}
	设 \(0<a<1\)。证明点值 \(F(v)=v(a)\) 是 \(H_0^1(0,1;\mathbb R)\) 上的有界线性泛函，求 \(-u''=\delta_a\) 的零边界弱解，并计算 \(\int_0^1|u'|^2\)。证明它不属于 \(H^2(0,1)\)。再用本章 Fourier 归一化证明，在 \(\mathbb R^n\) 上
	\[
		\delta_0\in H^{-s}(\mathbb R^n)\quad\Longleftrightarrow\quad s>n/2.
	\]
\end{exercise}
\begin{solution}
	一维 \(H^1\) 函数有绝对连续代表。对零端点函数，\(|v(a)|=|\int_0^av'|\le\sqrt a\|v'\|_2\)，故点值有界。分段线性函数
	\[
		u(x)=\begin{cases}(1-a)x,&0\le x\le a,\\ a(1-x),&a\le x\le1\end{cases}
	\]
	连续且两端为零，属于 \(H_0^1\)。其导数在 \(a\) 处的跳跃为 \(-1\)，分段分部积分给出 \(\int u'v'=v(a)\)，所以它是唯一弱解。直接计算
	\(\int|u'|^2=a(1-a)^2+(1-a)a^2=a(1-a)=u(a)\)。二阶分布导数为 \(-\delta_a\)，不能由 \(L^2\) 函数表示，故不属于 \(H^2\)。

	在全空间，\(\widehat{\delta_0}=(2\pi)^{-n/2}\)。因而负阶范数有限等价于
	\(\int_{\mathbb R^n}(1+|\xi|^2)^{-s}\,\mathrm d\xi<\infty\)。零附近无奇性，无穷远的极坐标积分与 \(\int_1^\infty r^{n-1-2s}\,\mathrm dr\) 同敛散，恰要求 \(2s>n\)。等号为对数发散。这说明点源是否属于能量对偶取决于维数。
\end{solution}

\begin{exercise}\label{g3:ex:div-curl-elliptic}
	对 \(u=(u_1,u_2):\mathbb R^2\to\mathbb R^2\)，令
	\[
		Pu=\bigl(\partial_1u_1+\partial_2u_2,\,
		\partial_1u_2-\partial_2u_1\bigr).
	\]
	计算 \(\sigma_1(P)\) 与 \(P^*P\)，证明 \(P\) 椭圆，并证明
	\[
		\|Pu\|_2^2=\|\mathrm Du\|_2^2\qquad(u\in C_c^\infty(\mathbb R^2;\mathbb R^2)).
	\]
	指出 \(\partial_1\) 单独作用于标量函数时在二维为什么不具有相应的全梯度估计。
\end{exercise}
\begin{solution}
	按 \(\mathrm D_j=-\mathrm i\partial_j\) 约定，
	\[
		\sigma_1(P)(\xi)=\mathrm i\begin{pmatrix}\xi_1&\xi_2\\-\xi_2&\xi_1\end{pmatrix},
		\qquad\sigma_1(P)^*\sigma_1(P)=|\xi|^2I_2.
	\]
	因此非零 \(\xi\) 时符号可逆，且常系数乘法给出
	\(P^*P=-(\partial_1^2+\partial_2^2)I_2\)。分部积分于是给出题设能量等式；也可直接用 Plancherel 逐频率计算。

	\(\partial_1\) 的符号是 \(\mathrm i\xi_1\)，在 \((0,\xi_2)\)、\(\xi_2\ne0\) 时为零。更具体地，取非零紧支光滑 \(\phi,\psi\)，令 \(u_N(x_1,x_2)=\phi(x_1)\psi(x_2)e^{\mathrm iNx_2}\)。其 \(\partial_1\) 范数保持不变，而 \(\partial_2\) 范数至少为 \(N\|\phi\psi\|_2-\|\phi\psi'\|_2\)，趋于无穷。实值估计若成立也会分别控制实、虚部，从而控制该复值函数，故实情形同样失败。
\end{solution}

\begin{exercise}\label{g3:ex:elliptic-order-sharp}
	设 \(P\) 为 \(\mathbb R^n\) 上 \(m\ge1\) 阶常系数椭圆系统，取非零 \(\phi\in C_c^\infty\)、非零 \(v\in\mathbb C^r\)、单位向量 \(\xi_0\)。令
	\[
		u_N(x)=N^{-m}\phi(x)e^{\mathrm iNx\cdot\xi_0}v.
	\]
	证明 \(\|Pu_N\|_2+\|u_N\|_2\) 有界，而 \(\|u_N\|_{H^{m+1}}\to\infty\)。由此说明二阶椭圆方程对 \(L^2\) 右端的普遍估计不能额外多得一个导数。
\end{exercise}
\begin{solution}
	每次求导落在指数上贡献一个 \(N\)，落在 \(\phi\) 上则不增加 \(N\)。所有不超过 \(m\) 阶的项因此在乘 \(N^{-m}\) 后一致有界，得到 \(Pu_N\) 的 \(L^2\) 界；\(\|u_N\|_2=N^{-m}\|\phi\|_2|v|\)。

	另一方面，频率平移公式给出 \(\widehat u_N(\xi)=N^{-m}\widehat\phi(\xi-N\xi_0)v\)。选有限的 \(R\)，使 \(c=\int_{|\eta|\le R}|\widehat\phi(\eta)|^2\,\mathrm d\eta>0\)。若 \(N\ge2R\)，该球中有 \(|\eta+N\xi_0|\ge N/2\)，故
	\[
		\|u_N\|_{H^{m+1}}^2
		\ge N^{-2m}(N/2)^{2m+2}c|v|^2.
	\]
	右侧为固定正数乘 \(N^2\)，所以范数趋于无穷。额外正则性可以来自右端或特殊解的附加信息，但不能由这一组统一的 \(L^2\) 数据界给出。
\end{solution}

\begin{exercise}\label{g3:ex:poisson-linear-constraint}
	设 \((M,g)\) 为连通、紧且无边界的光滑黎曼流形，\(f,\eta\in L^2(M;\mathbb R)\) 均为零均值，且 \(\eta\ne0\)。令 \(w,z\in H^1(M)\) 分别为 \(\Delta_gw=f\)、\(\Delta_gz=\eta\) 的零均值弱解。给定 \(a\in\mathbb R\)，求能量
	\[
		J(u)=\frac12\int_M|\mathrm du|^2\mu_g-\int_Mfu\,\mu_g
	\]
	在约束 \(\int_Mu\,\mu_g=0\)、\(\int_M\eta u\,\mu_g=a\) 下的唯一极小点，并写出它满足的弱方程。
\end{exercise}
\begin{solution}
	以 \(z\) 测试其方程，得到
	\(b=\int_M\eta z\,\mu_g=\int_M|\mathrm dz|^2\mu_g>0\)。若等号为零，\(z\) 是零均值常数，故为零；方程便迫使 \(\eta=0\)，与假设矛盾。于是取
	\[
		t=\frac{a-\int_M\eta w\,\mu_g}{b},\qquad u=w+tz.
	\]
	它满足两项约束，并满足 \(\Delta_gu=f+t\eta\)。若 \(v\) 为任意另一可行函数，则 \(h=v-u\) 为零均值且 \(\int\eta h=0\)，因此
	\[
		J(v)-J(u)=\frac12\int_M|\mathrm dh|^2\mu_g
		+\int_M(\Delta_gu-f)h\,\mu_g
		=\frac12\int_M|\mathrm dh|^2\mu_g\ge0.
	\]
	中间的配对按弱方程理解。等号使 \(h\) 为零均值常数，故 \(h=0\)，给出唯一性。约束产生的额外右端是 \(t\eta\)，并保留了闭流形上的零均值兼容条件。
\end{solution}

\begin{exercise}\label{g3:ex:nonelliptic-missing-direction}
	在 \(\mathbb R^2\) 上考虑 \(P=-\partial_1^2\)。取非零 \(\eta,\zeta\in C_c^\infty(\mathbb R)\)，令 \(u_j(x,y)=\eta(x)\zeta(y)e^{\mathrm ijy}\)。证明 \(\|u_j\|_2+\|Pu_j\|_2\) 一致有界，而 \(\|u_j\|_{H^2}\to\infty\)。再构造一个不光滑的分布 \(T\)，满足 \(PT=0\)。说明这两个现象如何对应于主符号缺失的方向。
\end{exercise}
\begin{solution}
	指数因子的模长恒为一，故
	\[
		\|u_j\|_2=\|\eta\|_2\|\zeta\|_2,\qquad
		\|Pu_j\|_2=\|\eta''\|_2\|\zeta\|_2.
	\]
	另一方面，
	\[
		\partial_y^2u_j
		=\eta(x)e^{\mathrm ijy}\bigl(\zeta''+2\mathrm ij\zeta'-j^2\zeta\bigr),
	\]
	所以
	\[
		\|\partial_y^2u_j\|_2
		\ge\|\eta\|_2
		\bigl(j^2\|\zeta\|_2-2j\|\zeta'\|_2-\|\zeta''\|_2\bigr)
		\longrightarrow\infty.
	\]
	因此无法用右端和零阶范数统一控制全部二阶导数。

	定义 \(T(\phi)=\int_{\mathbb R}\phi(x,0)\,\mathrm dx\)。由于测试函数紧支，
	\[
		(PT)(\phi)=-\int_{\mathbb R}\partial_x^2\phi(x,0)\,\mathrm dx=0.
	\]
	这个分布支于直线 \(\{y=0\}\) 而非零，不能由光滑函数表示：若光滑代表在直线外为零，连续性便迫使它处处为零。主符号为 \(\xi_1^2\)，在非零方向 \((0,\xi_2)\) 消失；函数列的高频振荡和分布的奇性都恰好位于未被算子控制的 \(y\) 方向。
\end{solution}

\begin{exercise}\label{g3:ex:jump-coefficient-regularity}
	在区间 \((-1,1)\) 上令
	\[
		a(x)=
		\begin{cases}1,&x<0,\\2,&x>0,\end{cases}
		\qquad
		u(x)=
		\begin{cases}x,&x\le0,\\x/2,&x>0.\end{cases}
	\]
	证明 \(u\in H^1(-1,1)\) 且 \(-\partial_x(a\,\partial_xu)=0\) 弱成立，但 \(u\notin H^2(-1,1)\)。计算 \(u''\) 的分布表达式，并解释这里哪一个系数条件失效。
\end{exercise}
\begin{solution}
	\(u\) 连续且分段线性，是 Lipschitz 函数，因此属于 \(H^1\)，其弱导数为
	\[
		u'=\mathbf1_{(-1,0)}+\frac12\mathbf1_{(0,1)}.
	\]
	连续性保证第一次分部积分没有原点的跳跃项。几乎处处有 \(a\,u'=1\)，所以对任意 \(\phi\in C_c^\infty(-1,1)\)，
	\[
		\int_{-1}^1a\,u'\phi'\,\mathrm dx=\int_{-1}^1\phi'\,\mathrm dx=0.
	\]
	再次求导得到 \(u''=-\tfrac12\delta_0\)。点质量不能由 \(L^2\) 函数表示，故 \(u\notin H^2\)。系数 \(a\) 始终介于一和二之间，椭圆下界保持成立，但它不属于 \(W^{1,\infty}_{\mathrm{loc}}\)。其差商在跨越原点时大小为 \(1/|h|\)，因此差商证明所需的统一系数界失效。保持连续的是通量 \(a\,u'\)，不是导数 \(u'\)。
\end{solution}

\begin{exercise}\label{g3:ex:point-source-sobolev-threshold}
	在 \(\mathbb R^n\)、\(n\ge1\) 上，令 \(\Delta_0=-\sum_j\partial_j^2\)。证明方程
	\[
		(1+\Delta_0)u=\delta_0
	\]
	有唯一的缓增分布解，且
	\[
		u\in H^s(\mathbb R^n)\quad\Longleftrightarrow\quad s<2-\frac n2.
	\]
	说明 \(u\) 在原点以外光滑，并证明它不可能在整个 \(\mathbb R^n\) 上光滑。将这个阈值与 \(\delta_0\) 的负 Sobolev 阶数比较。
\end{exercise}
\begin{solution}
	按本章的单位等距归一化，\(\widehat{\delta_0}=(2\pi)^{-n/2}\)。因此解必须满足
	\[
		\widehat u(\xi)=(2\pi)^{-n/2}(1+|\xi|^2)^{-1}.
	\]
	右侧为缓增光滑函数，定义一个缓增分布；乘以 \(1+|\xi|^2\) 恢复所需右端，故存在。该乘数在 \(\mathcal S'\) 上可用光滑倒数反转，故也唯一。

	其 \(H^s\) 范数平方有限恰好等价于
	\[
		\int_{\mathbb R^n}(1+|\xi|^2)^{s-2}\,\mathrm d\xi<\infty.
	\]
	原点附近没有问题；无穷远处的径向积分与
	\(\int_1^\infty r^{2s-4+n-1}\,\mathrm dr\) 可比，故充要条件为 \(s<2-n/2\)，等号时对数发散。同样计算给出 \(\delta_0\in H^t\) 当且仅当 \(t<-n/2\)，恰好相差两阶。

	在原点以外，右端为零，\cref{g3:cor:elliptic-smoothness}保证 \(u\) 光滑。若它在原点也光滑，则 \((1+\Delta_0)u\) 会是光滑函数，与非零点质量矛盾。这是局部正则性精确跟随右端奇性的例子。
\end{solution}

\begin{exercise}\label{g3:ex:interior-distance-constant}
	在单位圆盘 \(B_1\subset\mathbb R^2\) 中令
	\(u_j(r,\theta)=r^j\cos(j\theta)\)，其中 \(j\ge1\)。证明它们调和，并计算
	\[
		\|u_j\|_{L^2(B_1)}^2,\qquad
		\|\mathrm Du_j\|_{L^2(B_R)}^2
		\quad(0<R\le1).
	\]
	说明对每个固定 \(R<1\)，梯度与外层零阶范数的比值随 \(j\) 一致有界；但若要求同一个估计延伸到 \(R=1\)，则这样的常数不存在。
\end{exercise}
\begin{solution}
	\(u_j\) 是多项式 \(\operatorname{Re}(x+\mathrm iy)^j\)。直接用极坐标 Laplace 表达式计算，
	\[
		\partial_r^2u_j+r^{-1}\partial_ru_j+r^{-2}\partial_\theta^2u_j
		=\bigl(j(j-1)+j-j^2\bigr)r^{j-2}\cos(j\theta)=0
	\]
	于穿孔圆盘成立；由于它本来就是光滑多项式，原点处也成立。积分得到
	\[
		\|u_j\|_{L^2(B_1)}^2
		=\pi\int_0^1r^{2j+1}\,\mathrm dr=\frac{\pi}{2j+2}.
	\]
	而
	\[
		|\mathrm Du_j|^2
		=|\partial_ru_j|^2+r^{-2}|\partial_\theta u_j|^2
		=j^2r^{2j-2},
	\]
	故 \(\|\mathrm Du_j\|_{L^2(B_R)}^2=\pi jR^{2j}\)。两范数比值的平方为 \(2j(j+1)R^{2j}\)。固定 \(R<1\) 时，指数衰减控制这个二次多项式，因而上确界有限；取 \(R=1\) 则随 \(j\) 发散。内部估计中距外边界的正距离具有实际作用，即使算子系数恒定、方程右端恒为零，也不能删去这个条件。
\end{solution}

\begin{exercise}\label{g3:ex:anisotropic-fundamental-solution}
	设 \(n\ge3\)，\(A\) 为实对称正定矩阵，\(L_A=-\sum_{i,j}a^{ij}\partial_i\partial_j\)。用本章的 \(\Gamma\) 写出 \(L_A\) 的基本解，并验证其归一化因子。若 \(f\in C_c^\infty(\mathbb R^n)\)、\(u=\Gamma_A*f\)，证明对 \(1<p<\infty\)，\(\|\mathrm D^2u\|_p\le C(A,n,p)\|f\|_p\)。说明其中对 \(A\) 的依赖可由最小、最大特征值控制。
\end{exercise}
\begin{solution}
	令 \(B=A^{1/2}\)、\(y=B^{-1}x\)。若 \(u(x)=v(y)\)，链式法则给出 \(L_Au=-\sum_i\partial_{y_i}^2v\)。体积换元为 \(\mathrm dx=(\det B)\mathrm dy\)，所以
	\[
		\Gamma_A(x)=(\det A)^{-1/2}\Gamma(A^{-1/2}x).
	\]
	对测试函数 \(\varphi\)，换元后
	\[
		\int\Gamma_A(x)L_A\varphi(x)\,\mathrm dx
		=\int\Gamma(y)\left(-\sum_i\partial_{y_i}^2\right)(\varphi(By))\,\mathrm dy
		=\varphi(0),
	\]
	故归一化恰使右端为单位点质量。

	令 \(F(y)=f(By)\)，则 \(v(y)=u(By)=\Gamma*F(y)\)。二阶导数的 Fourier 乘子为 \(-\xi_i\xi_j/|\xi|^2\)，因而 \(\|\mathrm D_y^2v\|_p\le C(n,p)\|F\|_p\)。这一恒等式也可由基本解的二阶核公式直接识别；卷积的远处衰减使其没有额外的多项式项。又有
	\(\mathrm D_x^2u=B^{-1}(\mathrm D_y^2v)B^{-1}\)，体积换元的因子在两边范数中抵消，所以可取常数不超过 \(C(n,p)\|B^{-1}\|^2=C(n,p)/\lambda_{\min}(A)\)。对于采用一般坐标范数或同时控制原坐标与新坐标的区域版本，\(\|B\|=\lambda_{\max}(A)^{1/2}\) 同样控制区域的形变。
\end{solution}

\begin{exercise}\label{g3:ex:lp-endpoint-poisson}
	在 \(B_1\subset\mathbb R^2\) 取光滑截断 \(\eta\)，在 \(B_{1/4}\) 为一、在 \(B_{1/2}\) 外为零。令
	\[
		u(x,y)=\eta(x,y)(x^2-y^2)\log\sqrt{x^2+y^2},\qquad u(0)=0.
	\]
	证明 \(u\in W^{2,p}(B_1)\) 对每个有限 \(p\ge1\) 成立，\(-\partial_x^2u-\partial_y^2u\in L^\infty(B_1)\)，但 \(u\notin W^{2,\infty}(B_1)\)。由此解释 \(p=\infty\) 不能直接加入本章二阶 \(L^p\) 估计。
\end{exercise}
\begin{solution}
	在 \(r<1/4\) 中，\(u=r^2\log r\cos(2\theta)\)。极坐标 Laplace 公式给出
	\[
		(\partial_x^2+\partial_y^2)u=4\cos(2\theta)\qquad(r>0).
	\]
	远离零点的截断项均光滑有界。零附近 \(u=O(r^2|\log r|)\)、\(\mathrm Du=O(r|\log r|)\)、\(\mathrm D^2u=O(1+|\log r|)\)。在穿孔球上作一阶或二阶分部积分，内边界项分别受 \(Cr^3|\log r|\)、\(Cr^2|\log r|\) 控制而趋零，故这些经典导数在原点没有附加点质量，确为弱导数。因为 \(\int_0^{1/4}r(1+|\log r|)^p\,\mathrm dr<\infty\)，全部有限指数的结论成立。

	将 \((x^2-y^2)\log r\) 对 \(x\) 求两次导数，可写成 \(2\log r+b(\theta)\)，其中 \(b\) 有界。因此 \(\partial_x^2u\) 在任意足够小圆环上都有越来越大的绝对值，其本质上界为无穷。右端有界仍不能控制全部二阶导数的上确界；这与 Hölder 右端的 Schauder 估计不矛盾，这里的 \(\cos(2\theta)\) 在原点没有连续延拓。
\end{solution}

\begin{exercise}\label{g3:ex:cr-degenerating-symbol}
	对 \(a>0\)，在 \(\mathbb R^2\) 上考虑复标量算子 \(D_a=\partial_s+\mathrm i a\partial_t\)。求其主符号的逆，并证明存在 \(C_a\) 使
	\[
		\|\mathrm Du\|_{L^p}\le C_a\|D_au\|_{L^p},\qquad u\in C_c^\infty(\mathbb R^2;\mathbb C),\quad1<p<\infty.
	\]
	再证明，当 \(a\downarrow0\) 时任何这样的常数都至少按 \(1/a\) 发散。
\end{exercise}
\begin{solution}
	主符号为 \(\mathrm i\sigma-a\tau\)，其逆为 \((-\mathrm i\sigma-a\tau)/(\sigma^2+a^2\tau^2)\)，对非零实 \((\sigma,\tau)\) 可逆。各一阶导数乘以这个逆，得到光滑零阶齐次乘子；\cref{g3:cor:constant-elliptic-lp}给出所需估计。

	取非零实函数 \(\chi\in C_c^\infty(\mathbb R)\)，令 \(u_N(s,t)=\chi(s)\chi(t)e^{\mathrm iNt}\)。记 \(c=\|\chi(s)\chi(t)\|_p>0\)。三角不等式给出
	\[
		N^{-1}\|\partial_tu_N\|_p\longrightarrow c,
		\qquad N^{-1}\|D_au_N\|_p\longrightarrow ac.
	\]
	将题设估计用于 \(u_N\)，除以 \(N\) 后令 \(N\to\infty\)，得到 \(C_a\ge1/a\)。\(a=0\) 时主符号在纯 \(t\) 余方向上消失，失去椭圆性；常数的发散正是这个退化的定量表现。
\end{solution}

\begin{exercise}\label{g3:ex:critical-surface-composition}
	在二维单位圆盘中取光滑径向截断 \(\eta\)，在 \(r<e^{-3}\) 为一、在 \(r>e^{-2}\) 为零，令 \(u(r)=\eta(r)(\log(1/r))^{1/3}\)。证明 \(u\in W^{1,2}(B_1)\)，但对全空间光滑函数 \(F(t)=e^{t^4}\)，有 \(F(u)\notin L^1(B_1)\)。这说明曲面的临界 \(W^{1,2}\) 空间为何不能无条件承接任意光滑复合。
\end{exercise}
\begin{solution}
	令 \(t=\log(1/r)\)。零阶积分由 \(\int_3^\infty t^{2/3}e^{-2t}\,\mathrm dt<\infty\) 控制；一阶积分为常数乘
	\[
		\int_0^{e^{-3}}\frac{1}{r(\log(1/r))^{4/3}}\,\mathrm dr
		=\int_3^\infty t^{-4/3}\,\mathrm dt<\infty.
	\]
	穿孔圆盘上的分部积分内边界项至多为 \(Cr(\log(1/r))^{1/3}\)，趋零，所以计算出的导数确是全圆盘弱导数。截断区域没有奇性，故 \(u\in W^{1,2}\)。

	另一方面，零附近 \(F(u)=\exp(t^{4/3})\)。其积分包含
	\[
		2\pi\int_3^\infty\exp(t^{4/3}-2t)\,\mathrm dt=\infty.
	\]
	\(F\) 虽光滑，但 \(u\) 的值域无界，\(F'\) 在这一值域上也没有界。\(p>2\) 时的连续嵌入在有界光滑域上给出有界值域，正是本章复合命题可固定导数界的原因。
\end{solution}

\begin{exercise}\label{g3:ex:cr-scalar-gauge}
	设 \(\Omega\subset\mathbb C\) 开，\(\phi\in C^\infty(\Omega;\mathbb C)\)，记 \(D=\partial_s+\mathrm i\partial_t\)、\(D_\phi=D+(D\phi)\)。证明对任意局部 \(W^{1,p}\) 函数 \(u\)，
	\[
		D_\phi u=e^{-\phi}D(e^\phi u).
	\]
	据此描述 \(D_\phi u=0\) 的全部分布解。若某段边界上要求 \(u\) 的迹为实数，说明变换 \(v=e^\phi u\) 后的边界条件；何时它仍是取实值条件？
\end{exercise}
\begin{solution}
	光滑乘法的弱链式法则给出 \(D(e^\phi u)=e^\phi Du+e^\phi(D\phi)u\)，故恒等式成立；局部光滑乘数也作用于分布，同一公式在分布意义下成立。若 \(D_\phi u=0\)，则 \(Dv=0\)。\cref{g3:thm:cr-local-w1p}及其光滑性推论使 \(v\) 光滑，随后方程 \(Dv=0\) 就是经典 Cauchy--Riemann 方程，即 \(v\) 全纯。因而全部分布解恰为 \(u=e^{-\phi}h\)，其中 \(h\) 全纯。

	在可定义迹的光滑边界片，乘数与迹相容，所以 \(v\) 的迹位于实直线 \(e^{\phi(s,t)}\mathbb R\subset\mathbb C\)。这个直线等于 \(\mathbb R\) 当且仅当 \(\operatorname{Im}\phi\in\pi\mathbb Z\) 于该边界片；特别地，\(\phi\) 实值时保留实边界条件。局部算子的变换同时改变边界数据，故算子在内部共轭并不自动说明两套原边值问题相同。
\end{solution}

\begin{exercise}\label{g3:ex:schauder-holder-sharpness}
	设 \(0<\alpha<1\)。在 \((-1,1)\) 上求满足 \(u''=|x|^\alpha\)、\(u(0)=u'(0)=0\) 的解，并证明它属于 \(C^{2,\alpha}\)，但不属于任何 \(C^{2,\beta}\)，\(\alpha<\beta\le1\)。说明二阶椭圆方程保持右端 Hölder 指数的估计在这个模型中不能提高。
\end{exercise}
\begin{solution}
	两次积分得到
	\[
		u(x)=\frac{|x|^{2+\alpha}}{(1+\alpha)(2+\alpha)}.
	\]
	零点处前两阶导数由两侧极限连续延伸，且 \(u''=|x|^\alpha\)。凹性给出 \(|s^\alpha-t^\alpha|\le|s-t|^\alpha\) 对非负 \(s,t\) 成立，再用 \(\bigl||x|-|y|\bigr|\le|x-y|\)，可得 \([u'']_\alpha\le1\)。另一方面，取 \(y=0\)，有
	\[
		\frac{|u''(x)-u''(0)|}{|x|^\beta}=|x|^{\alpha-\beta}\longrightarrow\infty.
	\]
	因此更大的半范数无界。右端的精确振荡指数已经出现在解的二阶导数中。
\end{solution}

\begin{exercise}\label{g3:ex:continuous-rhs-not-c2}
	在平面小圆盘 \(B_{1/2}\) 上，对 \(0<r=\sqrt{x^2+y^2}<1/2\) 令
	\[
		u(x,y)=(x^2-y^2)\log\log(e/r),\qquad u(0)=0.
	\]
	证明 \(u\in C^1(B_{1/2})\)，并计算 \(f=(\partial_x^2+\partial_y^2)u\) 在穿孔圆盘上的表达式。证明令 \(f(0)=0\) 后，\(f\) 连续、分布方程仍成立，但 \(u\notin C^2\)。由此辨明把 Schauder 条件中的 Hölder 连续改成一般连续所遇到的困难。
\end{exercise}
\begin{solution}
	记 \(L(r)=\log(e/r)\)、\(\psi(r)=\log L(r)\)。有 \(u=r^2\psi(r)\cos2\theta\)，并且
	\[
		\psi'=-\frac1{rL},\qquad
		\psi''=\frac1{r^2L}-\frac1{r^2L^2}.
	\]
	函数大小为 \(O(r^2\log L)\)，梯度大小为 \(O(r\log L+r/L)\)，均以足够速度趋零，因此 \(u\) 在零点可微、导数为零，且梯度连续。极坐标 Laplace 公式给出
	\[
		f=(r^2\psi''+5r\psi')\cos2\theta
		=-\left(\frac4L+\frac1{L^2}\right)\cos2\theta.
	\]
	这在零点一致趋零，故 \(f\) 连续。对紧支测试函数在穿孔区域分部积分，内圆边界项的绝对值由
	\(C r(\sup_{|x|=r}|\mathrm Du|+\sup_{|x|=r}|u|)\) 控制，趋于零，因此没有额外点质量，分布方程在整个圆盘成立。

	沿正 \(x\) 轴两次求导，得到
	\[
		u_{xx}(x,0)=2\log L(x)-\frac3{L(x)}-\frac1{L(x)^2},
	\]
	它趋于正无穷，故 \(u\) 不能为 \(C^2\)。同时沿该轴，\(|f(x,0)|/x^\alpha\to\infty\) 对每个 \(\alpha>0\) 成立。连续右端的振荡未必足以使二阶势核在零点产生连续导数。
\end{solution}

\begin{exercise}\label{g3:ex:maximum-spectral-threshold}
	在 \((0,\pi)\) 上考虑 \(P_\mu u=-u''-\mu u\)，其中 \(\mu\in\mathbb R\)。先对 \(\phi\in C_c^\infty(0,\pi)\) 证明
	\[
		\int_0^\pi(|\phi'|^2-|\phi|^2)\,\mathrm dx
		=\int_0^\pi\sin^2x\left|\left(\frac\phi{\sin x}\right)'\right|^2\mathrm dx.
	\]
	由此证明：若 \(\mu<1\)、\(u\in H_0^1(0,\pi)\) 弱满足 \(P_\mu u=f\ge0\)，其中 \(f\in L^2\)，则 \(u\ge0\)。再分别说明 \(\mu=1\) 时的唯一性失效及 \(\mu>1\) 时的非负保持失效。
\end{exercise}
\begin{solution}
	令 \(w=\phi/\sin x\)。展开 \(\phi'=w'\sin x+w\cos x\)，积分交叉项，并用
	\((\sin x\cos x)'=\cos^2x-\sin^2x\)，得到所述恒等式；\(w\) 支于区间内部，所以边界项为零。因此 \(\|\phi\|_2\le\|\phi'\|_2\)，由 \(H_0^1\) 的定义延伸到整个空间。

	记 \(v=u_-\)。测试弱方程以 \(-v\)，得到
	\[
		\int_0^\pi(|v'|^2-\mu v^2)\,\mathrm dx
		=-\int_0^\pi fv\,\mathrm dx\le0.
	\]
	左端至少为 \((1-\mu)\|v\|_2^2\)，故 \(v=0\)。当 \(\mu=1\) 时，\(\sin x\) 是非零的齐次零边界解；当 \(\mu>1\) 时，取 \(f=\sin x\)，则
	\(u=\sin x/(1-\mu)\) 是一个严格为负的解。因而允许负的零阶项后，比较性质取决于它与域上能量下界的相对大小。
\end{solution}

\begin{exercise}\label{g3:ex:harnack-ellipticity-dependence}
	设 \(0<\varepsilon<1\)，在 \((-4,4)\) 上令
	\[
		a_\varepsilon(x)=
		\begin{cases}\varepsilon,&|x|<1,\\1,&1\le|x|<4,\end{cases}
		\qquad
		u_\varepsilon(x)=\int_{-4}^x a_\varepsilon(t)^{-1}\,\mathrm dt.
	\]
	证明 \(u_\varepsilon\) 是正的 \(H^1\) 弱解，满足 \((a_\varepsilon u_\varepsilon')'=0\)。计算它在 \((-1,1)\) 上的本性上确界与本性下确界之比，并解释这个例子对 Harnack 常数依赖的意义。
\end{exercise}
\begin{solution}
	对每个固定 \(\varepsilon>0\)，\(u_\varepsilon\) 为绝对连续函数，导数是有界的 \(a_\varepsilon^{-1}\)，故属于 \(H^1(-4,4)\)，且在开区间中为正。通量 \(a_\varepsilon u_\varepsilon'=1\) 几乎处处，测试任意内部紧支函数即得弱方程；系数在两个界面上的跳跃没有造成通量跳跃。

	在 \((-1,1)\) 上，\(u_\varepsilon(x)=3+(x+1)/\varepsilon\)。因此
	\[
		\frac{\operatorname*{ess\,sup}_{(-1,1)}u_\varepsilon}
		{\operatorname*{ess\,inf}_{(-1,1)}u_\varepsilon}
		=\frac{3+2/\varepsilon}{3}\longrightarrow\infty.
	\]
	区域及外尺度固定，系数上界也固定为一，但椭圆性下界趋零。故比较常数不能对这样的退化系数族保持统一。
\end{solution}

\begin{exercise}\label{g3:ex:cooperative-system-comparison}
	设常数 \(a>|b|\)，在 \((0,\pi)\) 上取
	\[
		C=\begin{pmatrix}a&b\\b&a\end{pmatrix},\qquad
		P U=-U''+CU.
	\]
	当 \(b\le0\) 时，证明任何零边界弱解 \(U\in H_0^1(0,\pi;\mathbb R^2)\) 在右端 \(F\in L^2\) 逐分量非负时也逐分量非负。当 \(b>0\) 时，以 \(F=(0,\sin x)\) 给出相反例子。说明主部及能量正定之外，什么附加结构控制了这里的次序性质。
\end{exercise}
\begin{solution}
	令 \(W=( (U_1)_-,(U_2)_-)\)、\(V=( (U_1)_+,(U_2)_+)\)，则 \(U=V-W\)，每个分量有 \(V_iW_i=0\)。测试向量弱方程以 \(-W\)，得到
	\[
		\int_0^\pi |W'|^2+W\cdot CW-V\cdot CW\,\mathrm dx
		=-\int_0^\pi F\cdot W\,\mathrm dx\le0.
	\]
	对角乘积 \(V_iW_i\) 为零，而 \(b\le0\) 使交叉项 \(V\cdot CW\le0\)。又因 \(C\) 的最小特征值为 \(a-|b|>0\)，左端至少为
	\(\int |W'|^2+(a-|b|)|W|^2\)，故 \(W=0\)。

	若 \(b>0\)，令 \(U=v\sin x\)，则
	\[
		v=(I+C)^{-1}\begin{pmatrix}0\\1\end{pmatrix}
		=\frac1{(1+a)^2-b^2}\begin{pmatrix}-b\\1+a\end{pmatrix}.
	\]
	分母为正，第一分量却为负。这里两种系统都有相同的正定标量主部，零阶能量也都正定；决定分量正锥是否保持的是非对角耦合的符号。
\end{solution}
